% ============================================================
%  H. Høffding, Relation som Kategori:
%  en erkendelsesteoretisk Undersøgelse
%  Det Kgl. Danske Videnskabernes Selskab,
%  Filosofiske Meddelelser I, 3.
%  København: Bianco Lunos Bogtrykkeri, 1921.
%
%  DIPLOMATIC TRANSCRIPTION (Danish)
%  Status: COMPLETE. Printed pp. 3–109 (body 3–108, Indhold 109), transcribed
%  2026-08-10 from the scan described below. 107 of 107 pages present, no gaps,
%  no duplicates; braces and guillemets balanced; compiles with 0 errors and 0
%  missing characters.
%
% ------------------------------------------------------------
%  THE WITNESS
% ------------------------------------------------------------
%  The Royal Danish Academy's own digitisation, downloaded 2026-08-10 from
%    http://publ.royalacademy.dk/books/109/739
%    -> uploads/2019-07-23/AFL 2/F_1_00_00_1920-1924_6255/F_1_03_00_1921_5831.pdf
%  114 pp., 408x612 pt, produced by ABBYY FineReader Server (2019). Kept
%  untracked in this directory as scan.pdf; the root .gitignore excludes
%  texts/**/scan*.pdf.
%
% ------------------------------------------------------------
%  PAGE MAP — VERIFIED MECHANICALLY 2026-08-10
% ------------------------------------------------------------
%      printed 3–109  =  PDF + 2   (PDF 5–111), uniform.
%
%  Verified by reading the running-head numeral off the text layer of every
%  PDF page 5–111: 105 of 107 return exactly the expected numeral. The two
%  exceptions are explained and neither is an offset change —
%      PDF   5 = printed 3    opening of §1, carries no running head
%      PDF 112 = printed 110  blank verso, no text at all
%  No duplicated leaves anywhere in the volume.
%
%  Unpaginated matter, outside the map:
%      PDF   1      series half-title (Filosofiske Meddelelser I, 3)
%      PDF   2      Videnskabernes Selskab series advertisement
%      PDF   3      title leaf — a SECOND setting of the same 1921 title,
%                   differing from PDF 1 only in capitalisation of the
%                   commission line
%      PDF   4      blank verso
%      PDF 112      blank verso after the Indhold
%      PDF 113–114  back-matter advertisements for the Selskab's Skrifter
%
%  The body is printed pp. 3–108; printed p. 109 carries the Indhold. The
%  title leaves are unnumbered but count as pp. 1–2, so there is no printed
%  p. 1 or 2 to transcribe.
%
% ------------------------------------------------------------
%  CONVENTIONS
% ------------------------------------------------------------
%  * Diplomatic. 1921 orthography exactly as printed — Kategori, Erkendelse,
%    absolut, Horizonten, ogsaa, Maalestok. Not modernised; in particular the
%    pre-1948 aa is kept and nouns stay capitalised.
%  * Quotation marks are GUILLEMETS »…« as printed, set \guillemotleft /
%    \guillemotright. This book does NOT use the low-high Danish pair „…“ that
%    the nineteenth-century books in this repo use.
%  * Em dash ---.
%  * Emphasis, three categories, all three verified against the images:
%      ITALIC (kursiv), chiefly the titles of works cited in the footnotes and
%        the occasional foreign phrase. Italic that is an author's stress ->
%        \emph{} (8 sites); italic that is a title or a foreign phrase set in
%        italic by convention -> \textit{} (74 sites).
%      SMALL CAPS, the footnote house style for proper names throughout —
%        Paul Barth, Brochard, Deussen, Croce, McTaggart, Poincaré, Ettore
%        Galli, Mach, Meyerson, Hobbes, Clerk Maxwell, Einstein, Whitehead,
%        Langevin, Kroman, Haldane, Westermarck, La Harpe, Hoefer and others ->
%        \textsc{}. Body-text names are ordinarily plain roman; the exceptions
%        are Rickert and Droysen (pp. 77–89), the book's only small caps outside
%        a note. NOTE that no mechanical witness sees small caps — ABBYY tags
%        them as ordinary roman — so these came off the images page by page.
%      LETTERSPACED (spærret). Confined to structure: the lettered subsection
%        heads of §4 (and §3's, which are set on their own line) and the whole
%        of the Indhold. It is NOT used for stress inside body text anywhere in
%        this book. That is a checked finding, not an assumption: spacing.py was
%        swept over all 107 pages and its 35 candidate runs were adjudicated
%        against 600 dpi crops. Every one on a page without a head proved to be
%        justification noise — word gaps opened, letters inside the words at
%        normal advance. p. 81 is the control: its spærret head sits two lines
%        from three flagged body words and the two look nothing alike.
%  * A word broken across a page turn is set WHOLE, with the division hyphen
%    dropped, a trailing % swallowing the line end, and the page marker at the
%    exact point of the break — see p. 6 ("Over" | "sigt"). Where the hyphen
%    belongs to the word it is kept.
%  * The printer's signature lines at the foot of pp. 33, 49, 65, 81, 97
%    ("Vidensk. Selsk. Filosof. Medd. I, 3." with a signature numeral) and the
%    signature marks 1*, 3*, 4*, 5*, 6*, 7* on pp. 3, 35, 51, 67, 83, 99 are
%    omitted as printer's apparatus, and logged at the site. The form varies:
%    the sheets signed on pp. 33 and 49 read "I. 3." with a full stop, those on
%    pp. 65, 81 and 97 read "I, 3." with a comma.
%  * Running heads ("Nr. 3. Harald Høffding:" verso, "Relation som Kategori."
%    recto) and folios are not transcribed.
%  * Footnotes are numbered 1, 2, … and RESTART ON EVERY PRINTED PAGE. Use
%    plain \footnote{}, and put \setcounter{footnote}{0} immediately after the
%    page marker of every printed page that carries a note.
%  * Page markers: % --- p. N --- on its own line where printed page N begins.
%    check.py depends on them; every page has one.
%  * Printer's defects are transcribed AS PRINTED and logged in a % comment at
%    the site, never silently corrected. The guillemet balance is therefore
%    expected to sit at 0 unless a defect is logged; see RESUME-NOTES.
%
% ------------------------------------------------------------
%  ERRATA AND DEFECTS LOGGED
% ------------------------------------------------------------
%  All transcribed AS PRINTED and logged again in a % comment at the site. None
%  is silently corrected. Verified at 600 dpi unless noted.
%
%   p.  12  "Cio … cio" for Ciò, in both places; "filosofia" with a plain i;
%           "Mc. Taggart" set as two words, for McTaggart.
%   p.  13  a stray low dot between "eller" and "Kræfter" — a loose sort or a
%           broken comma, not punctuation the sense wants. Not transcribed.
%   p.  39  "Gasus" for Gauss; a parenthesis opened at "(Discours" is never
%           closed.
%   p.  43  "henægtes" for benægtes; "Russel"/"Russel's" with one l, against
%           "Bertrand Russell" earlier in the same note. Both kept.
%   p.  44  the in-text reference is a superscript 1 but the note at the foot —
%           the page's only note — is numbered 2. Set here as 1, matching the
%           in-text mark; the printed 2 recorded at the site.
%   p.  46  "Forholdet" broken across a line with NO hyphen, "For" / "holdet".
%           Set here as one word.
%   p.  47  a note ends "men ogsaa p. 96 f)." — a closing parenthesis with no
%           opening one.
%   p.  53  a full stop where the sense wants a comma ("det andet.").
%   p.  59  "dévanciers" as printed.
%   p.  75  "Belydning" for Betydning (note 1).
%   p.  81  "under hvike" for hvilke.
%   p.  89  "alle Maader, paa hvilken".
%   p.  92  a stray mid-height mark between "det" and "er"; a broken or turned
%           sort, verified at 900 dpi, not punctuation.
%   p.  94  a like mark after "Relationsbegrebet."; and "Den Modsætning, her
%           foreligger", the relative pronoun absent as printed.
%   p. 100  a small raised tick sloping like an acute stands after "Eller" at
%           cap height; a stray or broken sort, the sentence running straight on.
%   p. 101  the note prints "socièté" — GRAVE on the first e, acute on the
%           second, checked against the unmistakable acute of "Poincaré" three
%           lines above.
%
%  The guillemet balance is nonetheless 0: none of the above disturbs it.
%
% ------------------------------------------------------------
%  STRUCTURE
% ------------------------------------------------------------
%  EVERY head below was read off the printed page at 300–600 dpi, not off the
%  Indhold. That distinction did real work here: the OCR of the Indhold renders
%  the Greek sub-section marks α β γ δ as Latin a b y d, and an earlier draft of
%  this file, built from that OCR before the scan was in hand, carried a
%  lettering for §3d that was in genuine doubt. It is now confirmed.
%
%      p.   3   1. Relationsbegrebets Historie.
%      p.  16   2. Relationsbegrebet i Psykologien.
%      p.  26   3. Relationsbegrebet i Erkendelsesteorien.
%      p.  26      a. Relationsbegrebet og de andre fundamentale Kategorier.
%      p.  33      b. Relationsbegrebet og de formale Kategorier.
%      p.  33         α.   [bare]
%      p.  40         β.   [bare]
%      p.  44         γ.   [bare]
%      p.  45      c. Relationsbegrebet og de reale Kategorier.
%      p.  45         α.   [bare]
%      p.  55         β.   [bare]
%      p.  58      d. Relationen Subjekt---Objekt.
%      p.  58         α. Indledning.
%      p.  62         β. Omkring Kopernikanismen.
%      p.  67         γ. Filosofiske Synspunkter.
%      p.  72         δ. Den nye Relativitetsteori.
%      p.  77   4. Relationsbegrebet i Etiken.
%      p.  77      a. Historieforskningen og Relationsbegrebet.
%      p.  81      b. Etiken og de psykologiske Relationer.
%      p.  87      c. Etiken og de historiske Relationer.
%      p.  90      d. Etiken og de individuelle Relationer.
%      p.  92      e. Kant's Etik og Relationsbegrebet.
%      p.  94      f. Religionsfilosofien og Relationsbegrebet.
%      p.  99   5. Relationsbegrebet i Kosmologien.
%      p.  99      a. b. c. d. at pp. 99, 103, 105, 107 — bare LATIN letters
%      p. 109   Indhold
%
%  THREE FEATURES OF THE PRINTING WORTH RECORDING, because a reader collating
%  this edition against the scan will meet all three:
%
%  1. The book is inconsistent about its Greek marks, and the inconsistency is
%     the printer's, not the transcriber's. In §3b and §3c the marks are BARE:
%     "α. Gennem Udformning af de formale Kategorier…" — the letter, a stop, and
%     the sentence running on at ordinary paragraph indent, with no title
%     anywhere on the page. In §3d every mark carries a printed title run in
%     with an em dash: "β. Omkring Kopernikanismen.¹ --- En af de Grunde…".
%     This was checked twice, by two agents working independently, once
%     specifically to test whether the first four titles had been missed. They
%     had not. The titles the Indhold gives the bare marks are therefore in this
%     edition's table of contents only, via \addcontentsline, and never in the
%     body — which is what the book does.
%  2. The lettered subsection heads differ likewise. §3's (a. b. c. d.) sit on
%     their own line, letterspaced. §4's (a. … f.) are run in, letterspaced,
%     followed by an em dash and then the text. In both cases the head is set
%     here as a heading and the run-in dash dropped, logged at the site.
%  3. §5 subdivides with bare LATIN letters a. b. c. d. (pp. 99, 103, 105, 107),
%     roman, no titles, no dash, at ordinary indent — and the Indhold does not
%     list them at all. They are left inline, with no sectioning command.
%
%  HEAD AGAINST INDHOLD. All 24 Indhold entries agree with the printed heads in
%  page number and, with two exceptions, in wording. Both readings are kept:
%      §4d  body "Etiken og de individuelle Relationer" / Indhold "Etik og de
%           individuelle Relationer". Verified at 400 dpi on both pages.
%      throughout: the Indhold's entries carry no terminal full stop, the
%           printed heads all do.
% ============================================================

\documentclass[12pt]{article}
\usepackage{fontspec}
\usepackage{unicode-math}
\setmainfont{Libertinus Serif}
\setmathfont{Libertinus Math}
\usepackage[danish]{babel}
\usepackage[protrusion=true]{microtype}
\usepackage[margin=1.4in]{geometry}
\usepackage{setspace}
\usepackage{fancyhdr}
\usepackage[colorlinks=true,linkcolor=black,urlcolor=black,
  pdftitle={Relation som Kategori},
  pdfauthor={Harald Høffding}]{hyperref}
\setlength{\headheight}{15pt}
\pagestyle{fancy}
\fancyhf{}
\fancyhead[LE,RO]{\thepage}
\fancyhead[RE]{\textit{Relation som Kategori}}
\fancyhead[LO]{\textit{\rightmark}}
\setlength{\parindent}{1.5em}
\setlength{\parskip}{0pt}
\setstretch{1.3}

% Section numbering to match Høffding's original:
%   1.  Main sections   (arabic)
%   a.  Subsections     (alpha)
%   α.  Sub-subsections (Greek letters written into the heading text)
\renewcommand{\thesection}{\arabic{section}.}
\renewcommand{\thesubsection}{\alph{subsection}.}

\title{\textbf{Relation som Kategori}\\[0.5em]
  \large en erkendelsesteoretisk Undersøgelse}
\author{Harald Høffding}
\date{Det Kgl.\ Danske Videnskabernes Selskab,
  \textit{Filosofiske Meddelelser} I, 3.\\
  København: Bianco Lunos Bogtrykkeri, 1921}

\begin{document}
\maketitle
\tableofcontents
\bigskip


% ============================================================
%  BATCH 1 — §1 Relationsbegrebets Historie
% ============================================================
% --- p. 3 ---
% PAGE MAP CONFIRMED BY EYE: printed p. 3 = PDF p. 5.
% This page carries NO running head and NO folio -- §1 opens here and the only
% thing at the foot is the signature mark "1*". That is why the mechanical
% running-head sweep skipped it; the offset is unchanged.
% The head is printed "1. Relationsbegrebets Historie." -- with the final stop,
% and agreeing with the Indhold apart from that stop.
% The first paragraph opens with a two-line DROP INITIAL "N" (the text layer
% loses it, returning "ærværende Afhandling er..."). Transcribed as ordinary
% text; no \lettrine.
\section{Relationsbegrebets Historie.}\label{sec:1}

\noindent Nærværende Afhandling er, ligesom en tidligere Afhandling om
\guillemotleft Totalitet som Kategori\guillemotright{} (Videnskabernes
Selskabs Skrifter, 1917), et Forsøg paa videre Udførelse af et enkelt Punkt i
min Bog \guillemotleft Den menneskelige Tanke\guillemotright{} (1910), i
hvilken det var min Bestræbelse at give en Fremstilling af de Hovedformer,
under hvilke menneskelig Tænkning arbejder, og i Sammenhæng dermed af de
Hovedproblemer, paa hvilke den arbejder. Sammenhængen mellem Tankeformer og
Tankeproblemer viser sig tydelig deri, at Tankeformerne ligger til Grund ved
de Spørgsmaal der stilles, Spørgsmaal, der, naar de stilles i en bestemt og
skarp Form, bliver til Problemer.

Efterat Interessen for Erkendelsesteori paany var bleven vakt i Slutningen af
forrige Aarhundrede, var det især den menneskelige Erkendelses første
Grundsætninger, der tildrog sig Opmærksomheden. Man drøftede Muligheden af at
begrunde eller udlede dem, og man drøftede Muligheden af deres fuldstændige
Gennemførelse paa alle Omraader. Men saadanne Grundsætninger indeholder visse
Grundbegreber, som mere eller mindre uvilkaarlig anvendes dels af den
saakaldte sunde Menneskeforstand (sens commun), den førvidenskabelige Tænken,
dels af den videnskabelige Tænken, saaledes som denne lægger sig for Dagen i
Videnskabernes Historie. Baade ved de Spørgsmaal, Tanken stiller, og ved
% --- p. 4 ---
de Svar, den anerkender, lægger visse Former sig for Dagen, som det er
Kategorilærens, den deskriptive Erkendelsesteoris Opgave at fremdrage og
ordne. De to Hovedkilder for Kategorilæren er det psykologiske Studium af den
\guillemotleft sunde Menneskeforstand\guillemotright{} og det historiske
Studium af den egentlige Videnskabs Udvikling gennem forskellige Stadier, og
disse Kilder staar i stadig Vexelvirkning med hinanden, idet Videnskaben
successivt udvikler sig af den sædvanlige Bevidstheds uvilkaarlige Eftertanke,
især naar der er opstaaet Tvivl, og idet paa den anden Side Videnskabens
Tankegang og Resultater langsomt faar Indflydelse paa den sædvanlige
Bevidstheds Form og Indhold, saa at den sunde Menneskeforstand ikke er den
samme til alle Tider.

Medens den tidligere Afhandling behandlede Totalitetsbegrebet og de med det
sammenhængende Tanketendenser, skal nærværende Afhandling søge at belyse
Relationsbegrebet og de Principer for Forskningen, der hænger sammen med det.
De to Begreber peger gensidig hen til hinanden. Totalitetsbegrebet gør sig
gældende ad to forskellige Veje. Al Erkendelse bestaar i Sammenfatten, en
uvilkaarlig Stræben efter at danne Helheder, hvor saadanne ikke er umiddelbart
givne. Men selve det umiddelbart Givne, de foreliggende Emner kan fremtræde
(og fremtræder ved nærmere Undersøgelse maaske altid) som Helheder, og de kan
kun forstaas ved Analyse af Sammenhængen mellem deres Elementer og med andre
Helheder og Elementer. Den menneskelige Erkendelse danner Helheder, hvor
saadanne ikke er umiddelbart givne, og opløser dem, hvor de umiddelbart
foreligger. I begge Tilfælde viser Helhedsbegrebet sig at hænge sammen med
Relationsbegrebet. Der kan kun dannes Helheder ved, at Emner bringes i
bestemte Forhold til hverandre, og givne Helheder forstaas kun ved Paavisning
af de indre og ydre
% --- p. 5 ---
Relationer, der er afgørende for deres Bestaaen. Relation viser sig derfor at
være en væsentlig Tankeform paa alle Omraader. At tænke er at finde bestemte
Forhold, sætte et Emne i Forhold til et andet, og de mest almindelige
Relationer viser sig at være Kontinuitet og Diskontinuitet, Lighed og Forskel.
Disse almindelige eller fundamentale Kategorier antager speciellere Former i
de formale Kategorier (Identitet, Analogi, Negation, Rationalitet), i de reale
Kategorier (Kausalitet, Totalitet, Udvikling) og i de ideale Kategorier
(Værdibegreberne). Den nærmere Paavisning heraf er given i
\guillemotleft Den menneskelige Tanke\guillemotright.

Relationsbegrebet vil nu i det Følgende modtage Belysning gennem Undersøgelse
af den Maade, hvorpaa det gør sig gældende ved Opstillingen og Drøftelsen af
de filosofiske Hovedproblemer. I mit Universitetsprogram af 1902 og i
\guillemotleft Den menneskelige Tanke\guillemotright{} (p. 246---254) har jeg
hævdet, at der er fire saadanne Problemer: det psykologiske, det
erkendelsesteoretiske, det etiske og det kosmologiske, og har paavist
Analogierne mellem dem. ---

\bigskip
% The printing leaves a full blank line here (line pitch 15.0 pt, this gap
% 29.7 pt) -- a thought-break, not a new section. Rendered as \bigskip.
Et Overblik over Kategorilærens Historie (se \guillemotleft Den menneskelige
Tanke\guillemotright{} p. 140---153) viser en afgjort Tendens til at stille
Relationsbegrebet i første Linie. Hos Aristoteles, den første, der har
opstillet en Kategorilære, er Relation (πρὸς τι) ganske vist kun én blandt ti
Kategorier, men en nærmere Betragtning viser, at de fleste (og i Grunden alle)
andre aristoteliske Kategorier udtrykker særlige Arter af Relation. Selv
nævner Aristoteles som Exempler paa Relation: Ligestorhed og Uligestorhed,
Aarsag og Virkning, Erkendelse og det Erkendte. Men i senere græsk Tænkning
kom Relation til at spille en aldeles afgørende Rolle, og her trænger den i
\guillemotleft Den menneskelige Tanke\guillemotright{} givne Over%
% --- p. 6 ---
\setcounter{footnote}{0}
sigt til en vigtig Supplering. Karakteristisk for denne senere græske
Tænkning, inden Mystiken gjorde sin Indflydelse gældende, var den
kritisk-skeptiske Retning, den slog ind paa. Platon staar, trods sin
spekulative Stræben som Forgænger for denne Retning, dels ved de
Indvendinger, han i \guillemotleft Parmenides\guillemotright{} lader komme til
Orde mod Idelæren, dels i den Vægt, han i senere Dialoger lægger paa, at
Erkendelse udtrykker sig i Dommen, altsaa i Bestemmelse af Begrebers
indbyrdes Forhold (smlgn. mine \guillemotleft Bemærkninger om den platoniske
Dialog Parmenides\guillemotright{} i Videnskabernes Selskabs filosofiske
Meddelelser, 1920). Men medens Platon trods alt var overbevist om, at
Sandheden var funden og var udtrykt i en Verden af evige Ideer, kom den senere
akademiske Skole til det Resultat, at den ikke blot ikke var funden, men ikke
kunde findes; hvad man kunde naa ved fornuftig Brug af Erfaringer, var kun
Sandsynlighed. Den betydeligste Mand af denne Skole er Karneades
(2. Aarh.\ f.\ Kr.), som med Rette er kaldt den mægtigste Aand i Grækenland i disse
Aarhundreder. Af afgørende Betydning for Relationsbegrebet blev det, at han
opkastede Spørgsmaalet om, hvori vort Kriterium for Sandhed og Virkelighed
bestaar. Overfor Stoikerne, der slog sig til Ro med, at der gaves
Forestillinger, i hvilke Sandheden aldeles umiddelbart kunde gribes
% Greek read off the image at 400 dpi. The printing uses the curly kappa (ϰ)
% of its Greek fount throughout; transcribed here as ordinary κ.
(φαντασίαι καταληπτικαί)\footnote{Om dette Udtryks Fortolkning se
\textsc{Paul Barth}: \textit{Die Stoa}, p. 66 f.}, hævder han, at da flere
indbyrdes stridende Forestillinger kan optræde med dette Præg, kan man ikke
stole paa nogensomhelst enkelt Forestilling. En Forestillings Gyldighed kan
kun bero paa dens Overensstemmelse med andre Forestillinger. Den maa ikke blot
være ἀπερίστατος, ikke modsiges af andre Forestillinger, men den skal finde
Bekræftelse ved sin Sammenhæng med andre Forestillinger, ligesom Lægerne
% --- p. 7 ---
\setcounter{footnote}{0}
ved deres Diagnoser ikke holder sig til noget enkelt Kendemærke, men støtter
sig til Sammenhængen mellem alle Kendemærker (συνδρομὴ τῶν σημείων).
Fremdeles maa der i alle vigtige Tilfælde kræves en \guillemotleft analyserende
Gennemgang\guillemotright{} af alle Forhold, der har med Forestillinger at
gøre (saa at Forestillingen bliver en φαντασία
διεξωδευμένη).\footnote{\textit{Sextus Empiricus} ed. Bekker, p. 225---232.}

Det vilde være urigtigt at kalde Karneades Skeptiker. Han tror ikke paa, at
nogensomhelst Forestilling har Krav paa Gyldighed bortset fra dens Forhold til
andre Forestillinger, men han henviser herved egentlig kun til den Vej, den
saakaldte sunde Menneskeforstand uvilkaarlig gaar, naar en Tvivl rejser sig,
og han henviser til Erfaringsvidenskabens Metode. I sin Fremstilling af
Karneades' Lære sammenligner Sextus hans Metode med de empiriske Lægers, og
der er ogsaa udtalt den Formodning, at Karneades staar i Gæld til denne
Lægeskole.\footnote{\textsc{Victor Brochard}: \textit{Les Sceptiques Grecs}.
1889, p. 183.} I det af Karneades hævdede Sandhedskriterium lægger
Relationsbegrebet sig tydelig for Dagen, iøvrigt paa en Maade, der allerede
antydes af Platon (Staten X, p. 602 C---603 A), saaledes som vi paa et senere
Sted vil faa Lejlighed til at omtale.

Hvis nu Sandhedskriteriet er det modsigelsesfrie Forhold mellem vel prøvede og
gennemtænkte Forestillinger, saa stilles derved en Opgave, der aldrig vil
kunne løses absolut. Der er ingen Grænser for, i hvor mange Forhold en
Forestilling skal undersøges, for at den kan staa sin Prøve. Det viser sig
her, hvad der atter og atter vil vise sig i det Følgende, at Relationsbegrebet
paa én Gang betinger en Grænse og en Opgave. Formedelst dette Begrebs
fremherskende Betydning i vor Erkendelse udtrykker enhver vunden Forstaaelse
visse bestemte Forhold og gælder
% --- p. 8 ---
kun for dem, er forsaavidt \guillemotleft relativ\guillemotright. Men Relation
som Kategori fører paa den anden Side stedse ud over enhver vunden
Forstaaelse, idet den henpeger paa Vigtigheden eller Nødvendigheden af at søge
nye Relationer, der kan gøre Bestemmelsen af det forstaaede Emne
fuldstændigere og dermed selve Forstaaelsen dyberegaaende. Hvad der øver
Modstand og stanser eller begrænser Tanken, stiller den netop derved nye
Opgaver, yder nye Impulser. Charles Renouvier har derfor med Rette talt om
Relation som Metode (la méthode des relations). I Relationsbegrebets Historie
lige til den nyeste Tid brydes de to Tendenser, Grænse og Opgave, Modstand og
Arbejde, Afslutning og Udvidelse. Og begge Tendenser udspringer af Relation
som Kategori.
% "la méthode des relations" is set ROMAN here, not italic (italics.py reports
% no italic anywhere on p. 8, and the image agrees). Left unmarked.

Allerede i Oldtiden viser denne Brydning sig. Akademikerne efter Karneades'
Tid var mest optagne af Afslutningen og nærmede sig derved Stoikerne og deres
Tro paa en én Gang for alle given Sandhed. Skeptikerne derimod mente, at naar
man aldrig kan blive færdig med at vinde den fulde Sandhed, har man ikke Ret
til at paastaa, at den er. Desuden kan der jo spørges, ved hvilket Kriterium
vi kan afgøre, at det af Karneades hævdede Kriterium er det rette. --- Medens
Dogmatikerne forkaster Relationsbegrebet og klamrer sig til visse
Forestillinger som definitive, hævder Skeptikerne ud fra Relationsbegrebet, at
der ingen Sandhed er, ti en relativ Sandhed er ingen Sandhed.

Den mest fremragende af disse Skeptikere, Ainesidemos (sandsynligvis i første
Aarh.\ f.\ Kr.), har givet en systematisk Fremstilling af de forskellige
Forhold, der efter hans Mening umuliggør sand Erkendelse. Han henviser til de
levende Væseners, særlig Menneskenes Forskelligheder, hvad Opfattelsesevne
angaar, --- til Forskellen mellem vore forskellige Sanser indbyrdes, hvorved
\guillemotleft en og samme Ting\guillemotright{}
% --- p. 9 ---
\setcounter{footnote}{0}
kommer til at staa som forskellig for forskellige Sanser, --- til den Maade,
paa hvilken vore Tilstande og vore Standpunkter i Rum og Tid, særlig
Afstandene fra Genstanden, betinger vor Opfattelse, --- og til Indflydelsen af
Vaner og Love. I det hele opstiller han 10 Arter af Forhold, der gør vor
Erkendelse usikker. Sextus Empirikus (ed. Bekker, p. 10) viser, at disse 10
Arter (τρόποι) kan føres tilbage til tre, idet enten det opfattende og
dømmende Subjekt (ὁ κρίνων), eller det opfattede og bedømte Objekt
(τὸ κρινόμενον), eller ogsaa begge paa én Gang kan bestemme Opfattelsen. Og
disse tre Arter af Forhold fører han igen tilbage til det almene
Relationsbegreb (ὁ πρὸς τι τρόπος). Her opstilles for første Gang Relation som
almindelig Kategori.

Med Urette har man i den Rolle, Drøftelsen af Virkelighedskriteriet spiller
henimod Oldtidens Slutning, set et Tegn paa Svækkelse af
Forskertrangen.\footnote{\textsc{Deussen}: \textit{Die Philosophie der
Griechen}. 1911, p. 410.} Det er netop nye og betydningsfulde Tanker, der
kommer frem i den akademiske og den skeptiske Skole, Synspunkter, som Platon
og Aristoteles i deres intellektuelle Begejstring ikke havde ladet komme til
deres Ret. Den Indsigt, at Afgørelsen af, hvad der kan anerkendes som
Virkelighed, bestemmes dels ved Forholdet mellem Emner indbyrdes, dels ved
Forholdet mellem Emnerne og det opfattende og tænkende Subjekt, er fra den
akademiske og skeptiske Skole gaaet over til den nyere Tids Filosofi og
Naturvidenskab\footnote{\textit{Den menneskelige Tanke}, p. 104 f.}, saaledes
som det i det Følgende vil vise sig. I Oldtiden selv dannede Drøftelsen af
dette Spørgsmaal Indledning til en ren empirisk Videnskab, særlig i
Lægeskolerne, og en videnskabelig Bevægelse i empirisk Retning var maaske, som
allerede
% --- p. 10 ---
\setcounter{footnote}{0}
bemærket, medvirkende ved selve Spørgsmaalets principielle Udspring. ---

I nyere Tid er det især Kriticisme og Positivisme, der har lagt Vægt paa
Relationsbegrebet. Kant, hvis Kategorilære er den betydningsfuldeste siden
Aristoteles, kom til sin Filosofi gennem den Indsigt, at Relationens Kategori
er den vigtigste af alle, idet vor Erkendelse overalt gaar ud paa at forbinde
Emnerne gennem Bestemmelse af deres indbyrdes Forhold. Det er to Arter af
Forhold, der især bliver af Betydning, Afhængighedsforhold og
Lighedsforhold, og det er i Arbejdet paa at tænke saadanne Forhold, at
Erkendelsen aabenbarer sig som
Syntese.\footnote{\textit{Om Kontinuiteten i Kants filosofiske
Udviklingsgang} (Vid.\ Selsk.\ Skr.\ 1893) p. 53.} Ogsaa hos Kant gør sig
Brydningen mellem de to Synspunkter for Relationsbegrebet gældende. Han har en
vis Tilbøjelighed til at betragte det som en Ufuldkommenhed ved vor
Erkendelse, at vi stiller alting i Relationer, og han underforstaar derfor en
\guillemotleft Ting i sig selv\guillemotright{} som liggende udenfor alle
Relationer. Nogle Kantianere, især William Hamilton, har betonet dette endnu
stærkere og sat det som en gennemgaaende Begrænsning for vor Erkendelse, at
der stedse maa arbejdes ved at sætte et Emne i Forhold til et andet. Paa
lignende Maade erklærede Herbart, at vor Forholdssætten kun kan give en
tilfældig Opfattelse af Tingene, og at sand Realitet slet ikke staar i
Relationer. Men hos Kant var dog alligevel det væsentlige Synspunkt, at det
Karakteristiske for al Videnskab er det i bestemte Love udtrykte Forhold
mellem Emnerne. Og fra dette Synspunkt bliver Relationens Nødvendighed for
Erkendelsen netop en Kilde til Fylde og Rigdom. --- Jo flere Relationer der
kan findes
% --- p. 11 ---
og bestemmes, des større Enhed og des større Mangfoldighed vil der paa én Gang
være i vor Erkendelse.

For Kants spekulative Efterfølgere blev Relationen mere og mere en Skranke,
man stræbte ud over ved at søge en højeste, altomfattende Enhed, indenfor
hvilken alle Relationer, særlig alle Modsætninger, skulde staa som ensidige
Opfattelser. Fichte staar her som en karakteristisk Overgangstype. Han har et
klart Blik for, at naar vi søger at finde og at bestemme Forhold, er det en
aandelig Aktivitet, som lægger sig for Dagen. Uden aandelig Virksomhed ingen
Forholdssætten: \guillemotleft Die reine Tätigkeit ist Bedingung des
Beziehens,\guillemotright{} (Wissenschaftslehre\textsuperscript{2}, p. 244).
% The superscript 2 on "Wissenschaftslehre" is an EDITION mark, not a footnote
% reference: p. 11 carries no note. Set with \textsuperscript, not \footnote.
Relation er ikke noget Modtaget eller Paatvunget, men forudsætter en Tankeakt.
Og naar der ved en saadan Akt sættes en Skranke, betyder det blot, at ny
Aktivitet er nødvendig. Aandeligt Arbejde, det teoretiske som det praktiske,
søger at naa ud over Skranker, og den fuldstændige Skrankefrihed staar som det
ophøjede Ideal, der aldrig kan virkeliggøres i nogen bestemt Tilstand. I
Fichtes senere Filosofi blev dog hans Problem et andet. Istedetfor at spørge,
hvorledes vi gennem Sætten af Forhold og gennem Kamp med de saaledes satte
Skranker kan naa stedse rigere Sandhed, spørger han nu paa nyplatonisk Vis,
hvorledes der overhovedet kan gives en Stræben efter Sandhed, naar der til
Grund for Alt ligger en absolut Væren.

Hegels Filosofi gik i sin historiske Orientering ud fra et Stade, paa hvilket
de menneskelige Tanker hvilede trygt ved hverandres Side, uden at der gjorde
sig nogen Trang gældende til at bestemme deres indbyrdes Forhold. En saadan
Trangs Opstaaen er efter hans Opfattelse et Opløsningsfænomen. Det er en
ulykkelig Bevidsthed, i hvilken Tankernes indbyrdes Harmoni er brudt. Da
behøves
% --- p. 12 ---
\setcounter{footnote}{0}
der et strengt aandeligt Arbejde, for at de enkelte Tanker, trods eller netop
ved deres indbyrdes Modsætninger, kan komme til at staa som Led i en stor
Helhed. Det sker ved Hjælp af den dialektiske Metode, der i Grunden betyder en
Systematisering af Relationsbegreber. Den enkelte Tanke fører ud over sig selv
til andre Tanker derved, at dens Begrænsning indses. De Modsigelser, der
fremkommer ved Tankernes Isolation, falder herved bort. --- En kritisk
Undersøgelse godtgør nu dog, at det i Grunden er Hegel selv, der isolerer
Begreberne for saa sejrrig bagefter at paavise deres højere Enhed. Saaledes
naar han lige i Begyndelsen af sin
Logik\footnote{\textit{Wissenschaft der Logik}. 1812. I, p. 22---26. --- Nyere
Hegelianere stiller sig her kritisk overfor Mesteren. Se \textsc{B. Croce}:
\textit{Cio che è vivo e cio che è morto della filosofia di Hegel} (1907),
p. 21---23. \textsc{Mc. Taggart}: \textit{A Commentary on Hegels Logic}
(1910), p. 17.}
% PRINTER'S READINGS in this note, transcribed as printed and NOT repaired:
%   "Cio ... cio" -- Italian requires the grave (Ciò); the printing has none, in
%   both places. Read at 600 dpi.
%   "filosofia" -- plain i, not the "filosofía" the text layer returns.
%   "Mc. Taggart" -- McTaggart set as two words with a stop after "Mc.",
%   the surname in small capitals.
stiller \guillemotleft Væren\guillemotright{} som noget for sig og
\guillemotleft Intet\guillemotright{} som noget for sig, for saa at vise, at
de ikke kan undvære hinanden. Enhver Væren har jo sin Grænse, ud over hvilken
der Intet er, naar ikke en ny Form for Væren kan findes. Og Begrebet Intet
betyder ikke andet end netop en saadan Grænse for Væren. Det er Erfaringen om
begrænset Væren, som Hegel opløser i to Begreber, der staar i Modsigelse til
hinanden. --- Men skønt Hegel selv saaledes ofte skaber de Modsigelser, hans
Dialektik skal overvinde, og skønt det ogsaa gaar altfor let for ham at komme
fra Relation til Relation, er hans Filosofi dog et stort Udtryk for den
Sandhed, at man kun kan tænke Tilværelsen ved at arbejde sig frem fra Relation
til Relation, og et Udtryk for, at i selve dette Arbejde er det Højeste, der
kan tænkes, tilstede som fremaddrivende.
% spacing.py flags four "runs" on this page (Isolation/herved, hvilken/Intet,
% Modsigelse/hinanden, skønt/Hegel). All four are loose justification, not
% spærret: checked on the image. No letterspaced emphasis anywhere in pp. 3--15.

Med fuld Bevidsthed har den franske Nykantianer Charles Renouvier stillet
Relationsbegrebet i Spidsen for
% --- p. 13 ---
\setcounter{footnote}{0}
sin Kategorilære. At alt er relativt --- siger han --- dette Skepticismens
stærke Ord, var den kritiske Filosofis sidste Ord i Oldtiden, og det maa være
den moderne Metodes første Ord. Og han finder i denne Kategori en hel Metode
indeholdt (la méthode des relations): det gælder at bestemme vore Emner
gennem saa mange Relationer som muligt. --- Og Kriticismen mødes her med
Positivismen. For Comte ligger Relationsbegrebet indeholdt i selve Begrebet
positiv Videnskab, der overalt søger at sætte bestemte, i Love udtrykte
Relationer i Stedet for Teologiens og Metafysikens absolute, men ubestemte
Væsener eller Kræfter.
% PRINTER'S DEFECT: a stray low dot stands between "eller" and "Kræfter"
% ("eller .Kræfter."), confirmed at 600 dpi. It is a loose sort or a broken
% comma, not punctuation the sense wants; not transcribed, logged here.
Fra denne Side betragter Comte udtrykkelig Kant som sin Forgænger, ligesom han
ogsaa mødes med ham deri, at Erkendelsesarbejdet opfattes som bestaaende i at
forbinde Emner (tout se réduit toujours à lier) --- dels gennem Klassifikation
(par similitude), dels gennem Afledning (par
filiation).\footnote{\textit{Discours sur l'esprit positif} (1844), p. 20.}
Som Dauriac\footnote{\textit{Contingence et catégorie} (Revue de Métaphysique
et de Morale. 1916), p. 499.} har paavist, er Comte's Klassifikation af
Videnskaberne anlagt efter samme Princip, som senere ligger til Grund for
Renouvier's Ordning af Kategorierne, nemlig saaledes, at der skrides frem fra
det Simple og Fundamentale til det Komplexe og Afledede.
% A second stray low mark sits immediately after "Revue" in note 2 -- same
% character of defect as the one before "Kræfter"; not transcribed.

Goblot\footnote{\textit{Traité de Logique} (1918), p. 138.} indvender mod
ethvert Forsøg paa at opstille en Kategorilære, at vor Viden endnu er for
provisorisk, baade hvad de sidste Slutninger og hvad de fundamentale Principer
angaar, til at et System af Grundbegreber skulde kunne opstilles. Denne
Indvending rammer ikke Renouvier, som er klar over, at der ikke kan føres
noget rationelt Bevis for Kategoriernes Gyldighed eller Kategorilistens Fuld%
% --- p. 14 ---
\setcounter{footnote}{0}
stændighed. De findes ved Analyse af Erfaringen, opstilles paa Prøve og maa
stedse paany bekræftes ved Erfaring.\footnote{Se Séailles: \textit{La
philosophie de Renouvier}, p. 91 f.} Renouvier bebrejder derfor Kant, at han
har forsøgt en deduktiv Begrundelse af Kategorilæren. I ejendommelig
Modsætning hertil har Renouvier's Discipel
Hamelin\footnote{\textit{Essai sur les éléments principaux de la
représentation}. (1907).}, alt imedens han fastholder sin Forgængers
Kategoritavle, søgt at give en dialektisk Udledning ved at vise, at enhver
Kategori fører over til sin Modsætning. Men han forveksler, som saa ofte
Hegel, Begrebers Korrelation med et indbyrdes Modsigelses- eller
Modsætningsforhold mellem dem; et saadant kommer kun frem, naar de enkelte
Begreber tages ud af deres Korrelation gennem en isolerende Abstraktion, som
saa ophæves igen og derved frembringer et Skin af Dialektik. Hamelin's
dialektiske Overgange er da ogsaa meget kunstige. --- Ud over en analytisk
Metode kommer man, som Renouvier meget rigtig har set, ikke i Kategorilæren.

Kategorierne søges og findes ved Analyse af det faktiske Arbejde, vor Tanke
øver, naar der stilles bestemte Spørgsmaal til den. Enhver Kategoritavle kan
derfor kun have en foreløbig Karakter, og de speciellere Kategorier kan ikke
udledes af de mere almene og omfattende, skønt de ved nærmere Undersøgelse kan
vise sig at være speciellere Former af dem. Som overalt, hvor vi har at gøre
med Begreber, der er uddragne af Erfaring (i dette Tilfælde af det faktiske
Tankearbejde), er der et omvendt Forhold mellem Indhold og Omfang. Jo flere
Elementer et Begreb indeholder, des snevrere en Kreds af Emner gælder det for.
Ved Kategorilæren faar dette en særlig vigtig Betydning ved Forholdet mellem
Værdibegreberne og de mere
% --- p. 15 ---
universelle (fundamentale, formale og reale) Kategorier. Men, som jeg i femte
Kapitel af Afhandlingen \guillemotleft Totalitet som
Kategori\guillemotright{} har søgt at vise, da al Værdi staar i Relation til en
vis Helhed og dennes Eksistensbetingelser, er der alligevel en indre
Sammenhæng mellem Værdibegreberne og de andre Kategorier.

Det er nu Opgaven for nærværende Afhandling at belyse Relationsbegrebet
nærmere ved at se det i dets Betydning for de forskellige filosofiske
Problemer. Det vil herved være hensigtsmæssigt ikke blot at holde sig til
selve Begrebet Relation, men ogsaa at undersøge det Princip, der svarer til
Begrebet, Relationsprincipet, som jeg (\guillemotleft Den menneskelige
Tanke\guillemotright, p. 255) har formuleret saaledes: Ethvert Emne skal
betragtes i saa mange Relationer som muligt. Hvor vi ingen Relationer kan
finde, faar vi ingen Erkendelse. Det forholdsløse er derfor med Rette af Kant
erklæret for at ligge udenfor al Videnskab og kun at være et rent negativt
Begreb. Vor Tænkning virker som en Passer: begge dennes Ben maa finde et Sted,
og vi faar da Kundskab gennem Forholdet mellem disse to Placeringer.
% The Relationsprincip itself ("Ethvert Emne skal betragtes ...") is set ROMAN,
% not italic and not spærret -- italics.py finds nothing on p. 15 and the image
% agrees. Left unmarked.

Naar Afhandlingen siges at handle om \guillemotleft Relation\guillemotright{}
og ikke om \guillemotleft Relativitet\guillemotright, er det, fordi Udtrykket
\guillemotleft Relativitet\guillemotright{} populært gerne misforstaas, idet
man mener (ligesom paa deres Vis de gamle Skeptikere), at naar noget er
relativt, vil det i Grunden sige, at det er ugyldigt eller ligegyldigt, ikke
har nogen objektiv Betydning. I Modsætning til denne Misforstaaelse maa
hævdes, at al Sandhed finder sit Udtryk i Paavisning af aldeles bestemte
Relationer og kun gælder indenfor dem, men saa ogsaa gælder absolut dér
(saafremt da Relationerne er bestemte rigtig). At al Nødvendighed, som
allerede Kant lærte, er hypotetisk, altsaa gældende i Relation til sine
Forudsætninger, vil ikke sige,
% Paragraph runs on across the page break into p. 16 -- batch 2 continues it.
% ============================================================
%  BATCH 2 — §1 concludes; §2 Relationsbegrebet i Psykologien (head at p. 16)
% ============================================================
% --- p. 16 ---
at saadan Nødvendighed ikke skulde være virkelig Nødvendighed. Ved al
Nødvendighed maa der spørges om Forudsætningerne, thi Nødvendighed bestaar i et
Forhold mellem Grund og Følge. Et saadant Forhold bliver kun til ved
Tankevirksomhed. Og ogsaa fra dette Synspunkt er Udtrykket
\guillemotleft{}Relativitet\guillemotright{} uheldigt ved at vække
Forestillingen om en hvilende Egenskab eller Tilstand. Men Relation bestaar i
Refereren, i en bestemt Indstilling af Tankens to Passerben. Der ligger i
Begrebet Relation en Anvisning til Tankearbejde.

% The head sits roughly two fifths down printed p. 16; §1's closing paragraph
% (at saadan Nødvendighed ... Tankearbejde.) runs on above it. Head read off
% the page image, not off the Indhold; the two agree exactly.
\section{Relationsbegrebet i Psykologien}
\label{sec:2}

Ethvert Forsøg paa at give en endog blot foreløbig Forklaring af, hvad der menes
ved Bevidsthedsliv eller Sjæleliv, vil komme til at lægge særlig Vægt paa den
Kendsgerning, at alt, hvad der sanses, forestilles, føles og villes, frembyder
visse bestemte Forhold, uden hvilke det ikke kan tænkes. Der er ganske vist
gjort Tilløb til at opfatte Sjælelivet som et Kaos af relationsløse Elementer;
men saadanne Tilløb har, trods den Energi, der undertiden ligger til Grund for
dem, ikke kunnet føre til en forstaaelig Psykologi. Sagen er, at fra sin første
Opstaaen er ethvert sjæleligt Element vævet ind i en Sammenhæng, ved hvilken det
bestemmes, og gennem hvilken det selv kan øve en bestemmende Indflydelse i
Sjælelivet. Det er herved, at Sjælelivet udgør en Helhed, saa at ikke mindst her
den ovenfor omtalte Forbindelse mellem Begreberne Totalitet og Relation gør sig
gældende. I saadanne Fænomener som Erindring, Sammenligning, Forundring og
Foretrækken lægger Bevidsthedslivets Karakter sig klart for Dagen, og i dem er
der netop mere end et Kaos, idet de forskellige
% --- p. 17 ---
\setcounter{footnote}{0}%
Elementer, som psykologisk Analyse kan skelne i dem, helt igennem bestemmes i
deres Ejendommelighed ved deres indbyrdes Forhold. Og alle Erfaringer peger hen
paa, at mere primitive og simple sjælelige Fænomener maa opfattes i Analogi med
dem; i hvert Tilfælde har Sproget ikke dannet særlige Betegnelser, der kunde
udtrykke Forskellen. Vi vil i det Følgende komme til at se, at dette hænger
sammen med den Grad af Bevidsthed, der kan være tilstede overfor selve
Bevidsthedslivets Forhold.

Det er Relationsbegrebets gennemgaaende Betydning i Bevidsthedslivet, der gør
det berettiget, som især Leibniz og Kant har indskærpet, at karakterisere
Bevidsthed som Syntese. En Relation kan kun gøre sig gældende, naar dens Led
sammenfattes; Relationen er jo netop Udtryk for en Sammenhæng, for Leddenes
Henhøren til en Helhed.

Den allerede omtalte dobbelte Karakter af Relationen, som begrænsende og som
videreførende, fremtræder tydelig i Psykologien. Ethvert sjæleligt Elements
Ejendommelighed beror tildels paa dets Forhold til andre sjælelige Elementer og
er forsaavidt begrænset ved dette Forhold. Men netop ved Forholdsbestemtheden
peger ethvert Element, naar det optræder pludselig og uden tydelig Sammenhæng
med andre Elementer, ud over sig selv og afføder en Trang til
Supplering.\footnote{I min Afhandling om Platon's
  \guillemotleft{}Parmenides\guillemotright{} (Vid.\ Selsk.\ filos.\
  Meddelelser.\ 1920.\ Kap.\ 3) har jeg gjort opmærksom paa, at medens Platon
  fandt det vanskeligt at finde en logisk Overgang fra et Begreb (en Ide) til et
  andet, nærede han som Psykolog ingen Tvivl om, at selv en pludselig
  opdukkende og i højeste Grad værdifuld Oplevelse havde sit bestemte Sted i
  Oplevelsernes Række, betinget ved, at en hel Række af Bestræbelser og
  Erfaringer var gaaet forud i bestemt Orden.}

Hvad Enkeltheder angaar, maa jeg henvise til min Fremstilling af Psykologien,
der netop er bygget paa den
% --- p. 18 ---
Grundtanke, at Forholdslovens Gyldighed for de sjælelige Fænomener godtgør
Bevidsthedens Karakter som Syntese. Kun nogle Hovedtræk skal der her mindes om.

Hvad Sansefornemmelserne angaar, stod det længe hen, hvorvidt Relationsbegrebet,
og med det Syntesebegrebet gjaldt for dem. Trods Hobbes' energiske Paastand, at
det at have en eneste uforanderlig Fornemmelse vilde være Et med slet ingen
Fornemmelse at have, var man tilbøjelig til at opfatte Fornemmelserne som et
Kaos, som først \guillemotleft{}Forstanden\guillemotright{} skulde ordne. I
Løbet af det 19.\ Aarhundrede kom man til at se Sagen anderledes. Chevreul's
Opdagelse af Kontrastvirkningen viste, at Fornemmelsernes Kvalitet beroede paa
Forholdet til forudgaaende eller samtidige Fornemmelsers Kvalitet, og Fechner
paaviste et analogt Forhold for den Intensitet, med hvilken de enkelte
Fornemmelser gør sig gældende. --- Man har været tilbøjelig til at opfatte det
som en Ufuldkommenhed, at vi ikke har, hvad man har kaldet en
\guillemotleft{}absolut\guillemotright{} Fornemmelse af noget, og selv Helmholtz
betragtede successiv Kontrastvirkning som en Illusion. Men man kan ikke paavise
nogen \guillemotleft{}Normalfornemmelse\guillemotright{}, hvis Kvalitet skulde
være den \guillemotleft{}rigtige\guillemotright{}. Og i hvert Tilfælde vilde en
saadan Normalfornemmelse selv være betinget ved visse Forhold, som man saa af en
eller anden Grund gav Fortrin for andre Forhold. --- Vi staar allerede her, ved
disse elementære Fænomener, overfor noget, der vil vise sig at gælde ved de
højeste intellektuelle Funktioner. ---

Gaar vi fra Fornemmelserne til Forestillingerne (der i deres simpleste Form er
reproducerede Fornemmelser), viser det sig, at Betingelsen for, at der i et
bestemt Øjeblik skal fremstaa en vis bestemt Forestilling, er det Forhold, i
hvilket denne Forestilling staar til andre Forestillinger.
% --- p. 19 ---
Man har vel forsøgt at tillægge enhver Forestilling en Slags
Selvopholdelsestrang, saa at den, saasnart andre Forestillinger (eller
Fornemmelser) ikke lægger Beslag paa den psykiske Energi, \guillemotleft{}af sig
selv\guillemotright{} skulde træde op i Bevidstheden. Men selv saa er der jo
bestemte Forhold --- Hæmning eller Ikkehæmning --- hvoraf Forestillingens
Optræden afhænger. Og selve den Selvstændighed, hvormed den enkelte Forestilling
gør sig gældende, betinges, som Erfaring viser, ved dens Forhold til det hele
øvrige Bevidsthedsliv og kan kun forstaas gennem dette Forhold. Der vil altid
være noget Unaturligt, og tilsidst Umuligt i at holde en Forestilling isoleret
fra alt Andet i Bevidstheden. Hver enkelt Forestilling har en Tendens til at
fremkalde andre Forestillinger efter bestemte Love. Atter her aabenbarer
Sjælelivets syntetiske Karakter sig. De fleste Psykologer er da ogsaa enige om,
at de saakaldte Associationslove er forskellige Former for Helhedsassociation,
d.\ v.\ s.\ for Trangen til Genfremkaldelse af hele den Bevidsthedstilstand, i
hvilken den enkelte Forestilling engang har været Led. ---

Paa Følelseslivets Omraade har man opdaget Forholdsloven, længe før man blev
klar over dens Betydning for andre Sider af Sjælelivet. Tydeligst fremtræder
Relationens Betydning ved Forundring, hvad enten denne er en særegen Følelse
eller er Led i mere komplicerede Følelser. Men ved al Lyst og Ulyst, Glæde og
Sorg er Forholdet mellem den forhaandenværende Tilstand og forudgaaende Tilstande
af afgørende Indflydelse. Den Lykke, der føles, beror ikke blot paa den bestemte
Størrelse af den Vinding eller det Fremskridt, der erfares, men ogsaa paa
Forholdet til det Følelsesniveau, som allerede forud var tilstede. Jo flere
forskellige Relationer, der er afgørende for en Følelse, des ejendommeligere og
rigere vil denne Følelse
% --- p. 20 ---
\setcounter{footnote}{0}%
være. Ved enhver saadan rigt udformet Følelse vil det blive Spørgsmaalet, om den
kan hævde sig overfor de nye Relationer, Sjælelivet kan stedes i. Saaledes f.\
Eks.\ naar Humoristen mødes af tragiske Skæbner\footnote{\textit{Den store
  Humor}, p.~102---104.}, eller naar en tilkæmpet Trøst skal staa sin Prøve
overfor nye Tilskikkelser.\footnote{\textit{Oplevelse og Tydning},
  p.~150---152.} ---

Ligesaalidt som der under aldeles ensformige indre og ydre Tilstande kunde
opstaa Fornemmelse, Forestilling eller Følelse, ligesaalidt kunde der være nogen
Villen. At ville er at foretrække og forudsætter altsaa, at der gør sig indre
eller ydre Forskelle gældende. Allerede Refleksbevægelse og Instinkt forudsætter
Forandringer i den indre organiske Tilstand eller i de ydre Forhold. Reaktion
forudsætter Aktion. Det gælder baade for den uvilkaarlige og for den vilkaarlige
Villen. Der gives Valg af alle Grader, men en Forskelsrelation forudsættes
stedse. Al Overvejelse, den uvilkaarlige saavel som den vilkaarlige, gaar kun ud
paa at faa saadanne Forskelsrelationer saa tydelig frem som muligt. De
forskellige Handlingsmuligheder maales efter deres Forhold til Sjælelivets
øjeblikkelige eller hidtidige Tilstand. Der kan her staa en Kamp, hvis Udfald
beror paa Relationerne til Individets Karakter og Fortid. Man har undertiden
opfattet Motiverne som absolut selvstændige, af Villien selv uafhængige
Elementer. Men ethvert Motiv er selv en Foretrækken, en Villen. Hvad der skal
kunne blive Motiv for et Individ, beror paa Individets Karakter og Forhistorie.
Motiver er aldrig relationsløse. Det er af afgørende Betydning, paa hvilket
Punkt af et Individs Udvikling en Handlingsmulighed optræder; derpaa
% --- p. 21 ---
vil det bero, om et vist Motiv overhovedet kan blive til, og om det kan blive
det sejrende.

% The printing leaves one full blank line here (measured at 400 dpi: 2x the
% normal baseline-to-baseline distance), a division stronger than a paragraph
% break. The following paragraph is indented as usual. Set as \bigskip.
\bigskip
Men Et er, at der gør sig en Forholdsmæssighed gældende overalt i Bevidstheden,
et Andet er, hvorvidt Bevidstheden selv bliver opmærksom paa de Relationer, der
i de enkelte Tilfælde er afgørende. Mange psykologiske Problemer opstaar ved, at
man ikke opdager Relationerne, idet disse maaske kun opdages ved objektiv,
historisk Undersøgelse, en Undersøgelse, Individet selv ikke, i hvert Tilfælde
ikke i selve Oplevelsens Tid, er istand til at foretage.

Kontrastforhold gør sig, baade hvad Kvaliteter og Intensiteter angaar, fra først
af gældende, uden at vi mærker det. Vi er optagne af det Liv og den Styrke,
hvormed det ved Kontrasten fremhævede nye Element fremtræder. Saaledes glemmer
vi, naar en Farve staar \guillemotleft{}mættet\guillemotright{}, let den
Relation, hvorved Mætheden betinges. Vi vender os mod en ny Følelse, der træder
frem med Liv og Fylde efter en Følelse af modsat Art, og glemmer den Besejrede
over Sejrherren. Et Forsæt eller en afgørende Beslutning kan melde sig som
\guillemotleft{}et Slag paa Hovedet\guillemotright{}, idet vi glemmer den
uvilkaarlige eller vilkaarlige Overvejelse, der har banet Vejen for dem gennem
et Sammenspil af Erindringer, Følelser og Drifter. Vi mærker ikke altid det
Associationsforløb, en Forestilling skylder sin Tilblivelse; den kan da staa som
en pludselig, umotiveret Indskydelse.

Vi har i det hele en Tilbøjelighed til at dvæle ved visse fremspringende Punkter
i vort Sjæleliv og overse de Relationer, der betinger Overgangen fra et saadant
Punkt til et andet. William James, der fortræffelig har belyst denne
Ejendommelighed ved Sjælelivet, skelner mellem Hvilepladser og Flyvepladser
indenfor Bevidsthedens Strøm og
% --- p. 22 ---
\setcounter{footnote}{0}%
mener, at vi har en Tilbøjelighed til at holde os til Hvilepladserne. Tydeligere
udtrykker han sig, naar han taler om substantiviske og transitive Dele af vort
indre Liv og mener, at vi er tilbøjelige til at tillægge hine en absolut
uafhængig Bestaaen, idet Overgangene ikke er lette at følge og ofte slet ikke
direkte kan iagttages, men maa sluttes at have været
der.\footnote{\textit{Principles of Psychology}. I, p.~243 ff.} Saadanne
substantiviske Elementer haves netop i kontrastbestemte Kvaliteter eller
Intensiteter, i umiddelbar Genkendelse (Genkendelseskvaliteten) og i en
pludselig opdukkende Forestilling, en udpræget Følelse, en overraskende
Beslutning.

James har desuden gjort opmærksom paa, at Relationer undertiden mærkes som
Kvaliteter, saaledes ved de sjælelige Tilstande, der i Sproget udtrykkes ved
\guillemotleft{}og\guillemotright{}, \guillemotleft{}men\guillemotright{},
\guillemotleft{}trods\guillemotright{} o.\ lgn. Det er en Art Fornemmelse (eller
Følelse) af Forbindelse eller af Modsætning, der kan være ligesaa umiddelbar som
de sædvanlig saakaldte Fornemmelser. Den saakaldte Genkendelseskvalitet er af
samme Art.

Herhen kan ogsaa regnes den mærkelige Fornemmelse af Retningen i vort Sjæleliv,
vi kan have. For at benytte Sammenligningen med Passeren kan vi sige, at medens
det ene Passerben staar fast (bestemt ved vor Fortid og dens Resultater), kan vi
mærke, at det andet Passerben er spændt ud i en vis Retning, uden at vi er klare
over, paa hvilket bestemt Sted det vil komme til at staa. Det er snart mere en
ubestemt Forventning, snart mere en ubestemt Higen, der her optræder.

Paa Overgangen mellem Psykologi og Erkendelsesteori ligger det Spørgsmaal, om og
hvorledes en Relation, man er bleven sig mere eller mindre klart bevidst, kan
overføres til andre Led end de oprindelige. En vigtig Maade,
% --- p. 23 ---
\setcounter{footnote}{0}%
hvorpaa dette kan ske, er følgende. Vi ordner uvilkaarlig vore Emner efter
Lighed og Forskel. Derved dannes der Rækker, og det kan da vise sig, at i en
saadan Række kan et Led skydes ud og Forholdet dog blive ved at gælde mellem
dets tidligere Forgænger i Rækken og dets tidligere Efterfølger. I Rækken ABC
kan maaske B skydes ud og Forholdet mellem A og C være det samme som mellem A og
B.\footnote{Smlgn.\ \textit{Den menneskelige Tanke}, p.~162 f.} Det er et
saadant Tilfælde, den formelle Logik konstruerer som et Forhold mellem
Præmisser og Konklusion. Præmisserne bliver f.\ Eks.: A = B og B = C, og
Konklusionen bliver A = C. Overhovedet gør Relation --- som netop derfor hører
til de fundamentale Kategorier --- sig gældende forud for al Begrebsdannelse,
Dømmen og Slutten. Allerede indenfor Sanse-, Erindrings- og Fantasianskuelse
ordnes og omordnes Elementer paa Grund af deres indbyrdes Forhold, uden at der
endnu er Tale om Tænken i snevrere Betydning af Ordet. Jeg henviser her til mine
Undersøgelser om Forholdet mellem Anskuen og Dømmen i første Hovedafsnit af
\guillemotleft{}Den menneskelige Tanke\guillemotright{} (især p.~53---59). Af
særlig Interesse er her den Maade, paa hvilken Anskuelsen artikuleres aldeles
uvilkaarlig og ubevidst, uden at en Dom dannes. Der gør sig en Stræben gældende
i Bevidsthedslivet efter at dvæle i Anskuelse saa længe som muligt uden at danne
Begreber og Domme eller drage Slutninger.

I Rækken ABC blev A staaende, men C kunde indtræde i samme Relation til A, som B
havde haft. Spørgsmaalet opstaar nu, om man ikke kan flytte begge Passerben, saa
at de sættes paa nye Steder, men Forholdet imellem dem vedbliver at være det
samme. I simpleste Tilfælde
% --- p. 24 ---
foreligger saadan Flytning, naar der foreligger en Relation mellem to nye Led (X
og Y), og denne Relation saa genkendes som den samme, man allerede kender mellem
to tidligere Led (A og B). Saadan Genkendelse kan --- ligesom Genkendelse af
Kvaliteter --- være enten umiddelbar eller middelbar; i sidste Tilfælde
forudsætter den udtrykkelig Udmaaling, eller ogsaa er der flere Relationer, der
tjener som Mellemled mellem de to Relationer, der identificeres. Den Intuition,
i hvilken det efter Sagnet gik op for Newton, at Forholdet mellem det faldende
Æble og Jorden kunde være det samme som Forholdet mellem Planeterne og Solen,
har først været umiddelbar; men den blev middelbar gennem de Maalinger og
Beregninger, der førte til den definitive Opstilling af hans Gravitationslære. I
det hele opstaar mange Spørgsmaal og Problemer netop ved, at Muligheden af
Relationsidentiteter ad en eller anden Vej optræder for Bevidstheden. Baade
Bestemmelsesspørgsmaal, som skyldes manglende Relationer, og
Afgørelsesspørgsmaal, som skyldes Strid mellem forskellige mulige Relationer,
kan føre Tænkningen videre. Allerede indenfor den sædvanlige Bevidsthed (sens
commun) kan saadanne Spørgsmaal med Forsøg paa Besvarelse optræde, og hvad vi
kalder Videnskab beror kun paa streng Begrundelse og Formulering af
Spørgsmaalene og skarp Undersøgelse af, om Svaret virkelig er en Løsning.
% The phrase sens commun is set in ROMAN here, not italic: checked on the page
% image and confirmed by italics.py, which finds no italic run on printed
% p. 24. Left unmarked accordingly.

Det er en stigende Grad af udtrykkelig Bevidsthed om Relationer, der fører over
fra uvilkaarlig Tænkning til logisk og matematisk Tænkning, der tumler med
Relationer som simple Kvaliteter, der kan blive Led i nye Rækker, og disse igen
i Rækkers Rækker. Et Eksempel paa Overgangen viser Udviklingen af Talbegrebet
fra Talkvalitet gennem
% --- p. 25 ---
\setcounter{footnote}{0}%
Løbenummer til Antal, og derfra igen til de forskellige
Funktionsbegreber.\footnote{\textit{Den menneskelige Tanke}, p.~189---193.}

% As on p. 21, the printing leaves one full blank line here (measured at
% 400 dpi). Set as \bigskip; the following paragraph is indented as usual.
\bigskip
En bestemt Grænse for Tænkningen vilde der foreligge, dersom der skulde gives
Rækker, hvis Led uden videre skulde kunne ombyttes, fordi Ligheds- eller
Forskelsforholdet var absolut det samme mellem alle Leddene. Saadanne Rækker
vilde være enten kaotiske Forskelsrækker eller absolute Identitetsrækker. Om de
forekommer i Tilværelsen, kan ikke paa Forhaand afgøres. Men fra Psykologiens
Standpunkt maa det siges, at de forekommer ikke i det faktiske Sjæleliv, hvor
der i det Højeste kan findes Tilnærmelser til dem. Det er en Grundlov i
Psykologien, at Addendernes Orden ikke er ligegyldig. Hvor der synes at
foreligge Kaos eller absolut Identitet mellem psykiske Elementer, vil det være
Psykologens Opgave at finde, i første Tilfælde, Ligheder eller bestemte
Afhængighedsforhold mellem de tilsyneladende absolut forskellige Elementer, i
andet Tilfælde, Forskelligheder og dermed følgende Afhængighedsforhold indenfor
det tilsyneladende absolut Ens. Om Kaos og om absolut Identitet kan der kun
tales ud fra bestemte Synspunkter, og i Relation til dem.\footnote{Anf.\ Skr.\
  p.~165; 170.} Altsaa Relationsprincipet gør sig ogsaa her gældende.
Tankearbejdet forudsætter et stadigt Vekselforhold eller en stadig Brydning
mellem Kontinuitet og Diskontinuitet, Lighed og Forskel; kun ved saadan Brydning
faar Tanken sine bestemte Opgaver og kan begynde at arbejde. Hume's Hævden af
kaotiske Elementer som det egentlig Givne gjorde, at efter hans Opfattelse kunde
intet Tankearbejde komme til at begynde. Spinoza's Hævden af de evige Attributer
% ============================================================
%  BATCH 3 — §3 Erkendelsesteorien (head at p. 26), §3a, §3b (head at p. 33)
% ============================================================
% --- p. 26 ---
\setcounter{footnote}{0}
som absolute Identiteter forudsatte, at alt Tankearbejde var
gjort. De to Tænkere endte, hver paa sin Vis, i Relationer,
der egentlig slet ikke var Relationer.

\section{Relationsbegrebet i Erkendelsesteorien.}\label{sec:3}
% p. 26: printed head reads "3. Relationsbegrebet i Erkendelsesteorien." —
% bold, centred, on its own line. The numeral comes from the counter.

\subsection{Relationsbegrebet og de andre fundamentale Kategorier.}\label{sec:3a}
% p. 26: printed head reads "a. Relationsbegrebet og de andre fundamentale
% Kategorier." — letterspaced (spærret) throughout, centred, set over two
% lines, NOT run in with the paragraph that follows. Matches the Indhold
% (p. 109) word for word. spacing.py's flag on "andre fundamentale" is the
% head's own letterspacing, not emphasis.

Den Analyse, der skal føre til Udfinden af Tankens
Former og Forudsætninger, maa ikke blot rettes mod det
i strengere Forstand videnskabelige Arbejde, men ogsaa
mod den uvilkaarlige Tænken, vi alle øver i det daglige
Liv, den, der er karakteristisk for, hvad man kalder \guillemotleft den
sunde Menneskeforstand\guillemotright. Af denne uvilkaarlige Tænken
har den videnskabelige Tænken, der gaar ud fra bestemte
Begreber og søger nøjagtig Begrundelse, Skridt for Skridt
udviklet sig, og der kan ved Overgangen til Videnskab,
saa stor og paafaldende Modsætningen end kan være, ikke
være nogen fuldstændig Forandring i Maaden at tænke
paa. Kant saa dette, da han udtalte, at selv \guillemotleft der gemeine
Verstand\guillemotright{} har visse \guillemotleft aprioriske Erkendelser\guillemotright{}
(Kritik der reinen Vernunft\textsuperscript{2} p.~3), og viser det endnu bestemtere i sin
Lære om Syntesen som Bevidsthedens Grundform, udsprungen
af \guillemotleft en Funktion, uden hvilken vi ikke vilde have
Erkendelse, skønt vi kun sjælden bliver os den bevidst\guillemotright{}
(Kritik der reinen Vernunft\textsuperscript{1} p.~78). Vi har ovenfor set,
% p. 26: the raised 2 and 1 after "Vernunft" are EDITION numerals, not
% footnote marks — checked on the image at 300 dpi. The one footnote on this
% page hangs on "Meyerson" below. "Kritik der reinen Vernunft" is set roman,
% not italic.
hvilken Betydning dette Begreb har for Psykologien. Ligesom
de andre fundamentale Kategorier er dette Begreb et
saadant, i hvilket Psykologi og Erkendelsesteori mødes. I
nyeste Tid har Meyerson\footnote{\textit{De l'explication dans
les sciences.} 1921. I p.~IX---XIV.} stærkt og paa meget lærerig
Maade fremhævet, at det er ganske den samme Tænkning,
% --- p. 27 ---
der øves i Naturvidenskab, Filosofi og sund Menneskeforstand
(sens commun). Dog gaar han ikke nærmere ind
paa de Trin og Veje, ad hvilke den praktiske Tænkning
successivt nærmer sig til Videnskaben. Han fremhæver
saaledes ikke Tvivlens Betydning som urovækkende ved,
at den praktiske Orientering i Verden ikke altid slaar til
eller ogsaa fører til modstridende Muligheder, overfor hvilke
man staar raadvild. Det er i Tvivlens Skole, at Tanken
kommer til at stille strengere Fordringer til sig selv.

Hvis man vilde definere Ordet \guillemotleft Fornuft\guillemotright{} saaledes, at
det kan passe paa den førvidenskabelige saavel som paa
den videnskabelige Tænken, kunde man sige, at den er et
Indbegreb af Kategorier, særlig af fundamentale Kategorier,
som maa anvendes paa alle Trin. Til disse fundamentale
Kategorier er at regne de tre Begrebspar Syntese-Relation,
Kontinuitet-Diskontinuitet og Lighed-Forskel.

Relation er (ligesom Syntese) en Grundform, der stadig
kommer igen ved alle specielle Tankeformer. Det kan da
ikke undre, at den særlig gør sig gældende indenfor hvert
af de nævnte Begrebspar. Saaledes er der Relation mellem
Syntese og Relation. Disse to Begreber kan nemlig træde
i Modsætning til hinanden, saaledes at Sammenfatten og
Enhed kommer til at raade paa Bekostning af Forholdsbestemtheden,
eller, omvendt, Leddene i Relationen gør sig
saa selvstændig gældende, at de ikke synes at være Led i
samme sammenfattede Hele. Paa den anden Side bliver
Trangen til Syntese stærkere, jo stærkere Leddenes indbyrdes
Spænding er. Naar Kant siger, at vi kun sjælden
bliver os den Funktion bevidst, der ligger til Grund for
Syntesen, saa er det, fordi det først er, naar Relationen
antager Karakter af Modsætning, Splid og Spænding, at
det sammenfattende Arbejde ret kan komme til Bevidst%
% --- p. 28 ---
hed. Heri ligger netop Tvivlens Betydning. Og den stadige
Korrelation mellem Syntese og Relation viser sig deri, at
Tvivl og Problemer kun eksisterer, saalænge forskellige
Emner, trods nok saa skarp Modsætning eller Modsigelse
mellem dem, sammenholdes i udtrykkelig Relation til hverandre.
Et Problem falder bort, naar man lader et af de
stridende Led falde bort, eller naar man svækker Modsætningen
mellem dem.

Der kan opstaa en Frygt for Relationernes Mangfoldighed
og Skarphed, som stiller den trygge Hvilen i uvilkaarlige
Synteser som særlig tillokkende. Derved fremkaldes
en saadan Reaktion, som Hamann, Herder og Jacobi øvede
overfor Kant's kritiske Filosofi, og som James og Bergson
betegner overfor Nutidens Videnskab. Og naar Tagore opfordrer
Vestevropæerne til at optage indisk Tankegang,
ligger deri netop en Opfordring til at se bort fra alle skarpe
og spændende Relationer, der er blevne udformede i den
evropæiske Civilisation, og til at vende tilbage til en relationsløs
Tilstand. Men saa stor og dyb end den Enhedssøgen
er, der udmærker Vedantalæren, hvis moderne Profet
Tagore er, saa er der ingen Vej tilbage. Enten er de
Relationer, der gør sig gældende, kunstige og ubegrundede,
og saa vil Kritiken fælde dem, eller ogsaa er de grundede
i faktiske Forhold, og da kan de kun magtes ved et aandeligt
Arbejde. Tvivl magtes kun enten ved, at dens Ubegrundethed
paavises, eller ved Problemløsning, eller ved
Paavisning af, at her foreligger et uløseligt Grænsespørgsmaal.
Naar Descartes netop fra Tvivlens Mulighed slutter
til Tankens Eksistens, saa betød det en Erkendelse af den
nøje Sammenhæng mellem Kategorierne Syntese og Relation.

Forholdet mellem Kontinuitet og Diskontinuitet er analogt
det mellem Syntese og Relation, eller er en speciellere
% --- p. 29 ---
Form for dette. Syntesen gaar des lettere, jo færre Afbrydelser
der sker. Om muligt optages uvilkaarlig saadanne
Afbrydelser i Anskuelsen (idet denne artikuleres bestemtere)
eller i Dommen (idet dennes Led faar Tilføjelser, der begrænser
eller udvider). Hvor Kontinuiteten raader, faar Anskuelse
og Eftertanke Karakteren af en jævnt henglidende
Strøm uden Hvirvler og Vandfald. Man spørger ikke, da
ingen Stansning eller Forundring motiverer noget Spørgsmaal.
Eller, hvis man spørger, lyder Spørgsmaalet: Hvorfor
ikke? Hvorfor ikke blive ved som hidtil, naar der er
Energi nok til, at Ekspansionen kan fortsættes? Men i
Livet og i Videnskaben gør mere eller mindre skarpe Arbejdsdelinger
sig gældende, og derigennem Modstand og
Modsætning, man ikke uden videre kan komme udenom.
Man maa da forsøge at finde Mellemled mellem de Led,
der er komne til at staa isolerede i Forhold til hverandre.
Eller man maa forsøge at paavise en Kontinuitet, der ligger
dybere end den Diskontinuitet, der har trængt sig frem,
en Kontinuitet, indenfor hvilke de foreliggende Spring eller
% p. 29: "indenfor hvilke" as printed (one would expect "hvilken", the
% antecedent being the singular "en Kontinuitet"). Left uncorrected.
Brud kan komme til at staa som afledede i Forhold til
den Lov eller den Type, den nye Kontinuitet frembyder.
Fra rent psykologisk Side er dette allerede omtalt ovenfor;
det er overhovedet Psykologiens Opgave at slaa Bro mellem
de ofte saa skarpe Modsætninger og Katastrofer, Sjælelivet
kan frembyde for den umiddelbare Iagttagelse. Særlig
er her James' Paavisning af \guillemotleft transitive\guillemotright{} Tilstande, i
Modsætning til de \guillemotleft substantiviske\guillemotright, af Interesse. Psykologisk
Diskontinuitet kan overhovedet kun betyde et Problem,
aldrig betegne en Løsning. Analogt betyder i Biologien
Mutationslæren, der hævder pludselig, af ydre Forhold
uafhængig, Opstaaen af nye organiske Arter, nye Problemer.
Selve Mutationslærens Grundlægger, Hugo de Vries,
% --- p. 30 ---
\setcounter{footnote}{0}
har straks udtalt, at Mutationerne betyder ligesaa mange
Problemer. Den moderne Kvantalære i Fysiken, ifølge
hvilken Energi ikke udløses paa kontinuerlig Maade, men
i Spring\footnote{\textsc{Poincaré}: \textit{L'hypothèse des Quanta.}
(Dernières Pensées. Chap.~6).}, søges nærmere begrundet gennem Paavisning af
% p. 30: the author's name in this note is set in SMALL CAPITALS in the
% printing; transcribed \textsc{}. The italic is the title only.
matematiske Love, ifølge hvilke den springvise Udløsning
staar som nødvendig. Der søges i hvert Tilfælde teoretiske
Synspunkter til Forstaaelse af Diskontinuiteten. Allerede
naar Springene kan ordnes i en vis Række, begynder Tanken
sit Arbejde paa at finde ny Kontinuitet, selv om denne
foreløbig er mere formal end real.

De Problemer, der her ligger, blev allerede rejst i den
ældste græske Filosofi. Parmenides' \guillemotleft Ene\guillemotright{} betød paa én
Gang Kontinuitet og Identitet; disse var for ham ikke forskellige.
Derimod hævdedes Diskontinuiteten af Herakleitos
og Demokritos. Gennem Platonismen og dens Modstandere
strækker Striden sig helt op i nyere Tid. Kant ansaa denne
Strid for en rationel Nødvendighed. Hans Lære om Antinomierne
gaar ud paa her at paavise en absolut Grænse
for Fornuften. Den kritiske Filosofis store Grundlægger saa
ikke, at Diskontinuitet aldrig kan blive Tankens sidste Ord
uden i Form af et sidste Problem, og hans \guillemotleft Antiteser\guillemotright, der
netop hævder dette Synspunkt, har derfor afgjort Ret overfor
hans \guillemotleft Teser\guillemotright, der hævder absolute Diskontinuiteter.

Kontinuitet og Diskontinuitet er Korrelata, der gensidig
supplerer hinanden. De betegner forskellige Synspunkter
og Operationer, og Videnskabens Historie viser, hvorledes
snart den ene, snart den anden af disse Kategorier kommer
til at staa i første Linie, men saaledes, at Kampen mellem
dem stedse vil begynde paany. Ingen Forsker har belyst
dette Forhold bedre end Henri Poincaré, især i følgende
Udtalelse: \guillemotleft Denne Kamp vil vare, saa længe man vil dyrke
% --- p. 31 ---
\setcounter{footnote}{0}
Videnskab, saa længe Menneskeheden vil tænke, thi den
udspringer af to uforsonlige Tendenser i den menneskelige
Aand, to Tendenser, som den ikke kan miste uden at ophøre
at eksistere, nemlig Trangen til at begribe, og vi begriber
kun det Begrænsede, og Trangen til at anskue, og
vi kan kun anskue Udstrækningen, som er ubegrænset.
Selv om denne Kamp ikke skulde ende med den ene af
de Kæmpendes sluttelige Sejr, saa vilde den dog derfor
ikke være ufrugtbar, thi i hver ny Kamp er Kamppladsen
en anden; hver Gang gøres der derfor et Skridt fremad og
gøres der en Erobring, ikke for den ene af de Kæmpende,
men for Menneskeheden.\guillemotright\footnote{\textit{Les conceptions
nouvelles de la matière} (i Sammelværket \guillemotleft Le
Matérialisme actuel\guillemotright, 1913), p.~67. --- Jfr. ogsaa den interessante
Diskussion i det franske filosofiske Selskab 13.~Marts 1910 (Bulletin de la
société française de philosophie. Vol.~X), især Perrin's, Brunschwicg's og
Meyerson's Udtalelser (p.~116---120).} --- Til denne smukke og træffende
Udtalelse maa jeg dog føje den Bemærkning, at Forholdet
mellem Anskuen og Tænken ogsaa kan have en
anden Karakter end den af Poincaré angivne. Det er ofte
netop Anskuelsen, der fastholder begrænsede Billeder, medens
Tanken, i Kraft af Loverkendelse, ser enhver Anskuelses
Begrænsning og Muligheden af nye Led i Billedernes
Række. Desuden er der flere Overgangsformer mellem
Anskuen og Tænken, som Tildels er antydede i det
Foregaaende, og der vil i hvert enkelt Tilfælde være bestemte,
i Erfaringen givne Emner, som afgør, om Anskuen
eller Tænken skal føre Ordet.

Hvor Kontinuiteten kan fastholdes, vil Betingelsen være,
at der ikke gør sig store Forskelligheder gældende i Tid,
Sted, Intensitet, Størrelse eller Kvalitet mellem de Emner, Bevidstheden
er optagen af. Hvor der raader Lighed, gaas der
let og, i hvert Tilfælde tilsyneladende, kontinuerlig over fra
% --- p. 32 ---
\setcounter{footnote}{0}
det ene Emne til det andet. Lighed og Forskel er et Begrebspar,
der frembyder Analogi med Parret Kontinuitet
og Diskontinuitet. Men medens det om Kontinuiteten kan
siges, at den ikke mærkes, før den er forbi, har Lighed
en noget mere bevidst Karakter. Den staar ikke i det nøje
Forhold til Anskuen, som Kontinuiteten gør. --- Jo større Ligheden
er, des bedre vil Kontinuiteten kunne bevares. Det
er først ved udtrykkelig, maaske vilkaarlig Sammenligning,
at man opdager Forskellen mellem de forskellige Grader
eller Arter af Lighed. Jeg har (i min Psykologi og i \guillemotleft Den
menneskelige Tanke\guillemotright{} p.~64 f.) skelnet mellem fem Arter af
Lighed: Identitet, Dækningslighed, sammensat Lighed, Kvalitetslighed
og Forholdslighed (Analogi). I denne Række træder
Forskelselementet stedse mere og mere frem, og det
er væsentlig gennem tiltagende Forskel eller aftagende Lighed,
at vi kommer over fra Kontinuitet til Diskontinuitet.

Relationerne Lighed og Forskel er af saa væsentlig Betydning,
at det er berettiget at definere Tænkning som
Sammenligning. At de to Relationer behersker al Tænkning
og alle Kategorier, indsaa allerede Porfyrios, da han
fandt det nødvendigt at skrive en Indledning til det aristoteliske
Skrift om Kategorierne\footnote{\textit{Isagoge in categorias}
(Scholia in Aristotelem ed.\ Brandis, p.~1).}, i hvilken han forklarede
Begreberne Slægt, Forskel, Art, Individ og Egenskab. Det
var denne \guillemotleft Indledning\guillemotright, der kom til at ligge til Grund for
de mange Drøftelser i Middelalderen om Almenbegrebernes
objektive Gyldighed (Striden mellem Nominalister og Realister).
De anførte fem Begreber danner en fremadskridende
Række, hvad Forskelligheder angaar, altsaa fra det Abstrakte
eller Almene til det Konkrete eller Specielle. For
Middelalderens Tænkning, der overvejende var klassifikatorisk
og syllogistisk, var Spørgsmaalet om, hvilket af disse
% --- p. 33 ---
Begreber der skulde anvendes i det enkelte Tilfælde, og
med hvilken Ret det anvendtes, af største Betydning. Den
nyere Videnskab foretager her en stor Reduktion. For den
er det ikke Forholdet mellem højere og lavere Begreber,
der er afgørende, men alt kommer an paa, om et Begreb
(et Emne) kan træde i Stedet for et andet i alle Henseender
eller ikke. Det er Identitet, der stiller de andre Arter
af Lighed i Skygge. Og netop ved Identitetsprincipets grundlæggende
Karakter som Nerven i al Slutten, og ved den
deraf betingede skarpe Bestemmelse af Identitetsbegrebet,
adskiller moderne Videnskab sig fra den sædvanlige Bevidsthed,
der ikke tager det saa nøje med Forskellen mellem
Identitet og blot Lighed. Parmenides har iøvrigt allerede
set, at Tankens Nerve ligger her, skønt det først var
Leibniz, der lagde Identitetsbegrebet til Grund for Logiken,
fordi det gjorde Substitution mulig. Det er i denne Sammenhæng
ligegyldigt, om man definerer Identitet som den
størst mulige Lighed eller som den mindst mulige Forskellighed;
begge Definitioner vidner om det korrelate Forhold
mellem Kategorierne Lighed og Forskel. --- Men ved
det afgørende Vendepunkt, der betegnes ved Identitetsbegrebets
Herredømme, kommer vi over fra de fundamentale
til de formale Kategorier.

\subsection{Relationsbegrebet og de formale Kategorier.}\label{sec:3b}
% p. 33: printed head reads "b. Relationsbegrebet og de formale Kategorier."
% — letterspaced (spærret), centred, on its own single line, not run in.
% Matches the Indhold (p. 109). spacing.py's flags on "Relationsbegrebet"
% and "formale Kategorier." are the head's own letterspacing, not emphasis.

% p. 33: the sub-subsection mark is GREEK ALPHA — α — read off the image at
% 300 dpi and re-checked at 600 dpi; unambiguous (round bowl, no ascender,
% and set in a face distinct from the surrounding roman). The text layer
% renders it "ct". It carries NO head text: the printed page runs straight on,
% "α. Gennem Udformning af de formale Kategorier (Identitet, ...", set at the
% ordinary paragraph indent. The Indhold (p. 109) titles this subdivision
% "Rene Kvalitetsrækker"; that wording appears nowhere on p. 33, so it is
% carried into the table of contents only, not into the text.
\addcontentsline{toc}{subsubsection}{α.~Rene Kvalitetsrækker}%
\label{sec:3b-alpha}%
α. Gennem Udformning af de formale Kategorier (Identitet,
Analogi, Negation, Rationalitet) bliver det muligt at
faa nøjagtig bestemte Relationer.

Identitetsbegrebet dannes som en naturlig Fortsættelse
af den Tendens, som allerede den praktiske Tænken har
til at danne Rækker af Emner, der ordnes efter tiltagende
Ligheder eller aftagende Forskelligheder. Dets Dannelse
% p. 33: the foot of the page carries the signature line
% "Vidensk. Selsk. Filosof. Medd. I. 3." with the signature numeral 3 at the
% outer edge; not transcribed. Likewise the "3*" at the foot of p. 35.
% --- p. 34 ---
muliggør Synspunkter, ud fra hvilke der ikke blot kan
vindes Oversigt over Relationernes Verden, men ved hvis
Hjælp den ene Relation maaske kan udledes af den anden.
Først paa dette Punkt træder den erkendelsesteoretiske og
den psykologisk-historiske Betragtningsmaade i afgjort Modsætning
til hinanden. Psykologisk og historisk gør der sig
i det højeste Tilnærmelser gældende til Identitet, og det
rene Identitetsbegreb afvises maaske endogsaa, for at de
Forskelligheder og Forandringer, Iagttagelsen frembyder,
kan komme til deres fulde Ret. Erkendelsesteoretisk betragtes
Forholdet fra den modsatte Side: Forskelligheder
og Forandringer staar som Brud paa den Identitet, som
er al Videnskabs første Forudsætning, og som betinger,
hvad der bliver et videnskabeligt Problem.

Identitetsbegrebets grundlæggende Betydning for videnskabelig
Tænken beror paa, at dets Anvendelse er Forudsætning
for de to vigtigste Tankefunktioner: Klassifikation
og Bevisførelse.

Inddeling af de forskellige Emner, som Iagttagelse viser
os, kan kun blive af afgørende Betydning, naar der er et
Synspunkt, ud fra hvilket Emnerne, trods al Forskellighed,
kan betragtes som identiske. I Slægts- og Artsbegreber
er der givet saadanne Synspunkter. De enkelte Emner, for
hvilke saadanne Begreber gælder, kan --- set fra vedkommende
Begrebs Synspunkt --- betragtes som et og samme
Emne, hvilket lægger sig for Dagen derved, at hvilketsomhelst
af dem kan bruges som Eksempel paa Begrebet.
Forskellighederne mellem de individuelle Emner falder dermed
ikke bort, men kan netop ved at blive forenede under
samme Synspunkt træde frem som skarpe Modsætninger.
Der er maaske end ikke en eneste Egenskab af
dem, Begrebet angiver som \guillemotleft fælles\guillemotright, der forekommer paa
% --- p. 35 ---
aldeles ens Maade i de enkelte Emner, der eksemplificerer
Begrebet. Men der raader dog en Analogi mellem de forskellige
Emner, idet Forholdet mellem deres Egenskaber
gennemgaaende er det samme.

Aristoteles slog med stort Eftertryk fast, baade at Videnskab
kun angaar det Almene, Indholdet af Slægts- og
Artsbegreber, og at det, der virkelig eksisterer, altid er individuelt.
Om det Individuelle kunde man efter hans Mening
ikke danne noget Begreb. Han har derved, uden at
være sig det klart bevidst, stillet et stort Problem, der
under forskellige Former strækker sig gennem hele den
videnskabelige Tænknings senere Historie. Idet vi her foreløbig
kun betragter Spørgsmaalet fra Klassifikationens Synspunkt,
maa der henvises til, at der dog kan dannes Begreber,
hvis Indhold er de Egenskaber eller Forhold, som
trods alle Forskelligheder i et individuelt Emnes enkelte
Tilstande kan udtrykke Emnets Identitet med sig selv.
Saadanne typiske Individualbegreber vil der mere paa
Aandsvidenskabens end paa Naturvidenskabens Omraade
være Mulighed for og Grund til at danne. Mellem de forskellige
Tilstande af samme Emne raader der, ligesom
mellem de forskellige Eksempler paa samme Almenbegreb,
kun Analogi, ikke Identitet.

Hvor hverken Analogi eller Identitet længere kan paavises,
staar vi overfor et Kaos. Men selve Begrebet Kaos
faar kun virkelig Betydning gennem Modsætningsforholdet
til Identitet. Der er Forskel mellem det Kaos, der staar
for den sædvanlige Bevidsthed, naar den overvældes af
Forskellighederne, og det Kaos, som statueres af en Tanke,
der i alvorligt Arbejde har søgt efter Analogier og Identiteter,
og nu tilsidst nedlægger Arbejdet. Og selv i dette
sidste Tilfælde kan det aldrig med Sikkerhed hævdes, at
% --- p. 36 ---
\setcounter{footnote}{0}
vi staar overfor en kaotisk Forskelsrække. Det er strengt
taget ikke berettiget at sige andet, end at det hidtil ikke
er lykkedes at finde Enhedssynspunkter for de Emner, der
synes at danne saadanne Rækker. Ligesom Identitet kan
defineres ved den størst mulige Lighed eller den mindst
mulige Forskellighed, saaledes kan Kaos defineres som den
størst mulige Forskellighed eller den mindst mulige Lighed.
Ligheden er saa lille, at den hidtil ikke bemærkedes,
ligesom vi ved de Emner, vi betragter som identiske, kan
sige, at Forskellighederne er saa smaa, at de ikke bemærkes.

Begrebet Negation dukker op her som Udtryk for, at
en Afbrydelse af Tankearbejdet --- her af det inddelende
Tankearbejde --- viser sig nødvendig, da ingen flere Lighedspunkter
kan findes. Klassifikationernes Historie viser,
at man i saadanne Tilfælde ofte danner negative Begreber
(eller Begreber, der kun indeholder negative Kendemærker),
som naar Lamarck inddelte Dyrene i Hvirveldyr og Hvirvelløse;
senere fandt Cuvier positive Kendemærker for de
Hvirvelløse og Mulighed for en positiv Inddeling af dem.

Den kaotiske Forskelsrække saavel som den absolute
Identitetsrække betegner to ideale Grænselinier for Erkendelsen,
Grænselinier, som aldrig kan paavises i virkelig
Erkendelse. Den ene af disse Grænser vilde betegne Erkendelsens
Umulighed\footnote{Stuart Mill forsøgte at føre al Slutten
tilbage til en associativ Overgang fra en Enkelthed til en derfra forskellig
Enkelthed. Men en saadan Overgang er allerede rent psykologisk kun mulig,
naar den første Enkelthed genkendes; kun da kan det forstaas, hvorledes den
kan fremkalde den anden Enkelthed, med hvilken den forhen fulgtes. Det er
dog snarest en Analogislutning, der gør, at der kan sluttes \guillemotleft fra
Enkelthed til Enkelthed\guillemotright: som a i en tidligere Iagttagelse forholdt
sig til b, vil a' nu forholde sig til b'. Men saa maa der jo dog være Lighed
mellem a og a'. Smlgn. \textit{Den nyere Filosofis Historie} (Første Udg.\ II,
p.~377; Anden Udg.\ II, p.~414). (I tredie Udgave: III, p.~119). Smlgn.\ i
nyeste Tid: \textsc{Ettore Galli}: \textit{Nel regno del cognoscere.} (1919)
p.~159---167.}, den anden dens Afslutning. Mellem
% p. 36: this note runs over the page break — its last four lines are printed
% at the foot of p. 37, which carries no note of its own. Set here in one
% \footnote{}. The letters a, b, a', b' are roman, not italic, in the printing;
% the primes are printed as right single quotes. The author's name in the
% second citation is in SMALL CAPITALS, "Stuart Mill" in the first is not.
% --- p. 37 ---
begge ligger en Skala af Tankerækker (af hvilke nogle af de
vigtigste er fremstillede i \guillemotleft Den menneskelige Tanke\guillemotright{}
p.~162---174), der betegner forskellige Trin, ad hvilke Erkendelsen
arbejder sig hen imod sit Ideal. ---

Saalænge Erkendelsen kun tilstræber en Klassifikation,
lader den Forskellighederne bestaa og søger blot at naa en
ordnet Oversigt over dem. Ved Bevisførelse derimod faar
Identitetsbegrebet aktiv eller produktiv Betydning, afgiver
et Middel til at paavise nye Sandheder, maaske finde nye
Emner uden stadig at støtte sig til Iagttagelse. Her opdukker
Rationalitetens Kategori, der fremtræder i Forholdet
mellem Grund og Følge. Dens Forudsætning er en Identitet
mellem to Emner, der berettiger til overalt at sætte det
ene i Stedet for det andet. De to Præmisser, af hvilke Slutningen
fremgaar, maa indeholde et Begreb, der er identisk
i dem begge. Allerede Klassifikationen gør Slutning mulig,
efter at der er dannet identiske Slægts- og Artsbegreber
(eller typiske Individualbegreber). Men Videnskaben søger
at gaa et Skridt videre. Den danner ikke blot ordnede
Rækker af kendte Led, der kun kan fortsættes ved Hjælp
af nye Iagttagelser, men ogsaa Rækker, hvis Fortsættelse
kan ske efter samme Lov, efter hvilken den er begyndt.
Her fremtræder Forskellen mellem Logik og Matematik.
Matematikeren opstiller ikke blot Identitetsprincipet som
gældende for de Enere, med hvilke han arbejder, altsaa
1=1, men han opstiller ogsaa den Forudsætning, at han
af disse Enere kan danne en Række, som opstaar ved, at
Ener paa stedse samme Maade føjes til Ener, altsaa
1 + 1 + 1 . . . En Forudsætning som denne sidste kender
den rene Logik ikke. For den har Gentagelse af et Emne
% --- p. 38 ---
\setcounter{footnote}{0}
ingen tankemæssig Betydning, saa megen psykologisk eller
pædagogisk Betydning saadan Gentagelse kan have. Den
holder sig til Identitetsprincipet (og de dermed mere eller
mindre direkte sammenhængende logiske Principer). Den
matematiske Grundforudsætning er derimod (ved Siden af
Identitetsprincipet) Gentagelsens positive og produktive Betydning.
Det er denne Forudsætning, Poincaré har kaldet
loi de recurrence, og i hvilken han finder Betydningen af,
% p. 38: "loi de recurrence" is set ROMAN in the printing, not italic —
% checked on the image; italics.py reports no italic run here either. Left
% unmarked rather than given \textit{}.
hvad Kant har kaldt en apriorisk-syntetisk Dom.\footnote{Smlgn.\ herom
\textit{Totalitet som Kategori} (Vid.\ Selsk.\ Skr.\ 1917)}
% p. 38: the note ends without a full stop after the closing parenthesis,
% as printed.

Medens Klassifikationen lod de kvalitative Forskelligheder
ligge, som de var, og blot ordnede dem, søger den
moderne Naturvidenskab at bearbejde eller udrense Kvalitetsforskellighederne
saaledes, at der af dem kan dannes
Rækker af lignende Art som de nys omtalte, altsaa rent
formale Rækker, i hvilke nye Led stedse kan konstrueres
efter samme Lov, som allerede har givet de foregaaende
Led deres Plads. Der er fire Arter af Kvalitetsforskelligheder,
der særlig har gjort en saadan Udrensning mulig,
nemlig Tid, Tal, Grad og Sted. For umiddelbar Iagttagelse
fremtræder de ikke som \guillemotleft rene\guillemotright, men som Kvaliteter. Forskellen
mellem Tider, mellem Steder, mellem Antal og
mellem Grader, er ligesaavel kvalitativ som Forskellen
mellem Gult og Blaat, Haardhed og Blødhed. Videnskabens
Historie viser, hvilket Tankearbejde der har været nødvendigt,
inden de nævnte fire Kvalitetsrækker har kunnet omdannes
til Rækker, i hvilke hvert enkelt Led er fuldt bestemt
ved sin Plads i Rækken, en Plads, der selv er bestemt
ved Forholdet til de andre Led. Hvad der udfylder
Øjeblikkene eller sker paa de forskellige Steder, hvad der
tælles eller grademaales, er nu ligegyldigt. Hvert enkelt
% ============================================================
%  BATCH 4 — §3b concludes; §3c (head at p. 45)
%  A head is listed in the Indhold at p. 55 — find it.
% ============================================================
% --- p. 39 ---
\setcounter{footnote}{0}
Led i Rækken betragtes, som om det havde samme Indhold som de
andre, saa at kun Pladsen gør Forskel. I korte Træk har jeg skildret
denne Udrensningsproces i \guillemotleft Den menneskelige
Tanke\guillemotright{} (p.~186---201), og jeg skal ikke gaa nærmere ind
paa den her. Jeg skal kun tilføje, at de fire Udrensningsprocesser ikke
foregaar uafhængig af hverandre, men gensidig støtter hverandre, og at
dette bliver muligt ved, at Kategorierne overhovedet opstaar sporadisk;
en vis Udvikling af den ene Kategori kan derved faa Betydning for
Udviklingen af den anden. Saaledes har Tidsbegrebet udviklet sig med
Støtte af Talbegrebet og Rumsbegrebet. Og det synes, som om
Talbegrebet overhovedet har været til stor Hjælp ved Udrensningen af de
tre andre Kategorier. Den Proces, ved hvilken Ener føjes til Ener, er
simplere og lettere end den, hvorved Øjeblik føjes til Øjeblik, Grad til
Grad og Sted til Sted. Allerede Kant saa, at der i Tællen havdes det
simpleste Eksempel paa \guillemotleft Synthesis des
Gleichartigen\guillemotright . Comte har endnu stærkere fremhævet dette,
og ogsaa Gauss har udtalt, at Tallet i langt højere Grad end Rummet er
et Produkt af vor egen Aand.%
% The note's proper names are set in SMALL CAPS in the printing (Kant,
% Comte, Gasus, Mach); rendered \textsc{}. This is the footnote house
% style throughout the batch — body-text names are plain roman.
% PRINTER'S DEFECT (1): »Gasus« as printed, for Gauss.
% PRINTER'S DEFECT (2): the parenthesis opened at »(Discours« is never
% closed. Transcribed as printed.
% The superscript 1 after »Irrtum« is an edition number, not a note mark.
\footnote{I de af Erdmann udgivne \guillemotleft
Refleksioner\guillemotright{} siger \textsc{Kant}, at det egentlig kun er
Tal og Størrelser, der kan konstrueres ud af ren Fornuft, fordi der ved
dem blot behøves Gentagelse. (Smlgn. \textit{Totalitet som Kategori},
p.~36). --- \textsc{Comte} fandt i de rene Talspekulationer
\guillemotleft la véritable origine de tout le système
scientifique\guillemotright{} (\textit{Discours sur l'esprit positif.}
1844, p.~100. --- \textsc{Gasus} erklærede i et Brev Tallet for et
Produkt af vor Aand, medens Rummet ogsaa har en Realitet udenfor denne.
(Se Citaterne hos E. \textsc{Mach}: \textit{Erkenntnis und
Irrtum}\textsuperscript{1}. p.~384).}

Medens typiske Individualbegreber, som angaar konkrete Emner, kun kan
udvikles under stadigt Henblik til og Laan fra Iagttagelse, er der
gennem de nævnte Udrensningsprocesser tilvejebragt Mulighed for en
Dannelse af Tankerækker, der kan fortsættes efter samme Princip, som
% --- p. 40 ---
\setcounter{footnote}{0}
ligger til Grund for dens Dannelse, og ud fra hvilke der kan stilles Krav
til de Emner, Iagttagelsen frembyder.

% Sub-subsection mark read off the page image at 600 dpi: GREEK BETA (β),
% upright, in the roman face, followed by a full stop. It stands bare at
% the head of an ordinary indented paragraph — the printed page gives NO
% title. The title below is the Indhold's (p. 109, set spærret), recorded
% in the ToC only, never printed into the text.
\addcontentsline{toc}{subsubsection}{β.~Identitet og Relation}%
\label{sec:3b-beta}%
β.~Ved Dannelsen af de omtalte rene Kvalitetsrækker ligger overalt
Identitetsbegrebet til Grund, idet de er karakteriserede ved, at
Begreberne Tid, Tal, Grad og Sted er frigjorte for alt Indhold, der
udfylder Tiden, udgør Tallet, frembyder Grad, indtager Sted; med
Indholdet bortfalder alle reale Forskelligheder indenfor Rækken. Naar de
endnu maa kaldes Kvalitetsrækker, kommer det af, at Forskellene mellem
Tid, Tal, Grad og Sted indbyrdes bliver ved at bestaa. Der er, som der
senere vil blive Lejlighed til at omtale, en Tendens til at lade Tallet
blive raadende tilsidst og gøre al Erkendelse til et System af Ligninger.
Dog brydes her, som Meyerson har
vist\footnote{\textit{De l'Explication dans les Sciences.} I p.~264;
II p.~167; 204---220.}, Algebra og Geometri om Forrangen. Selve
Spørgsmaalet hører ind under de reale Kategorier.

Men saa høj en Grad af Eksakthed, som den Erkendelse, der arbejder med
de nævnte Rækker, end kan naa, saa melder Relationsbegrebet sig selv i
disse rene Regioner. Dette viser sig især paa tre Punkter.

Identitet er selv en Relation, et Emnes eller en Tankes Forhold til sig
selv under forskellige Omstændigheder. Hvilken Definition af Identitet
man end benytter, saa undgaas Forholdsbestemtheden ikke. Enten man taler
om den størst mulige Lighed eller den mindst mulige Forskellighed, saa
forsvinder Skæret af Relation ikke. Og særlig stærkt bliver dette Skær,
naar man benytter Leibniz's eksperimentale Definition: identisk er de
Tanker eller Emner, der uden Ophævelse af Gyldigheden (salva veritate)
kan sættes i Stedet for hinanden overalt.%
% This note begins on printed p. 40 and RUNS OVER onto printed p. 41,
% where it occupies the whole footnote block (p. 41 carries no note of
% its own). Given here entire, at its mark.
\footnote{Som omtalt i \guillemotleft Totalitet som
Kategori\guillemotright{} (p.~41 f.) har der hos danske Matematikere
gjort sig forskellige Opfattelser gældende overfor det Spørgsmaal, om
deres Videnskab kræver Grundsætninger, der udtrykker absolut Identitet,
eller om Tilnærmelse til en saadan er tilstrækkelig. Denne Modsætning
minder om de to forskellige Definitioner paa Identitet --- som størst
mulig Lighed eller som mindst mulig Forskel. I det sidste Arbejde, der
foreligger fra \textsc{Zeuthens} Haand, udtaler han
(\textit{Hvorledes Matematiken i Tiden fra Platon til Euklid blev
rationel Videnskab.} Vid. Selsk. Skr. 1917, p.~80), at Anvendeligheden af
Euklid's geometriske System beror paa, \emph{at} de empiriske Linier har
de i Postulaterne angivne Egenskaber, hvad Euklid forudsætter, eller om
man vil, den afhænger af, \emph{hvorvidt} de har den. Ved dette
\guillemotleft hvorvidt\guillemotright{} hentydes sikkert til den
Opfattelse, som \textsc{Hjelmslev} har gjort gældende, og ifølge hvilken
der ikke behøves ideale Definitioner som de euklidiske, men at man kan
gaa ud fra Linier og Figurer, som kan tegnes.} Ved at lade det komme an
paa Mu%
% --- p. 41 ---
% No footnote of its own; the whole footnote block on this page is the
% run-over of note 2 of p. 40, set above at its mark.
ligheden af Substitution i alle Tilfælde benytter man tydelig nok netop
Relationen mellem vedkommende Tanker (Emner) og andre; Gyldigheden
bevares kun, hvis saadan Relation ikke ændres, naar Substitutionen sker.
Selv ved Identiteten 1 = 1 kan Spørgsmaalet blive, om vi virkelig raader
over Enere, der overalt kan sættes i Stedet for hverandre. Ved
Diskussionen om Fechners Forsøg paa at give Weber's Lov en strengt
matematisk Formulering kom dette Spørgsmaal frem.

Men for det andet forudsætter Opstillingen og Anvendelsen af
Identitetsbegrebet altid et bestemt Synspunkt, ud fra hvilket det ene
gælder. Ved Klassifikation er det tydeligt nok, at Identiteten kun
gælder visse Kendemærker, der stadig kommer igen, saaledes at der i det
mindste er Analogi mellem deres Former i de enkelte Tilfælde. Men ogsaa
den Identitet, der er Nerven i enhver Bevisførelse, gælder kun fra et
vist Synspunkt. Identiteten er stedse naaet gennem Udrensning, d.~v.~s.
ved Bortseen fra Synspunkter, der ikke kommer den Sag ved, som drøftes.
Naar de forskellige Naturkræfter bestemmes som Energier, d.~v.~s. maales
ved Evnen til at udføre et Arbejde, er derved til%
% --- p. 42 ---
\setcounter{footnote}{0}
vejebragt et Enhedssynspunkt, ligesom naar ellers højst forskellige Ting
betragtes ene og alene under Synspunktet Tyngde. Ligetil
Spektralanalysens Fremkomst betragtede man i Henhold til Newtons Fysik
Kloderne i Verdensrummet væsentlig under Tyngdens Synspunkt. Det er ved
at anlægge et rent kvantitativt Synspunkt --- et Synspunkt, der igen kun
bliver muligt ved, at man holder sig til visse bestemte Egenskaber ---
at der kan skaffes Identitet tilveje mellem ellers højst uensartede
Emner. Ligesom i de beskrivende Videnskaber Kvalitet kan reduceres til
Identitet formedelst Slægts- og Artsbegreber, saaledes kan i de eksakte
Videnskaber Kvalitet reduceres til Identitet formedelst
Kvantitet.\footnote{Smlgn. \textit{Den menneskelige Tanke},
p.~177---180.} Platon har givet dybsindige Antydninger i begge
Retninger; hvad det sidste Punkt angaar, optog Galilei paa genial Maade
hans Tanke.

I græsk Tænken blev der ofte øvet en Art Kultus overfor det
\guillemotleft Ene\guillemotright , Indbegrebet af de rene Identiteter.
Det fremtræder hos Parmenides og Platon og naar sit Højdepunkt i
Nyplatonismen. For den nyere Tids Tænkning er Identitet ikke Genstand
for Kontemplation, men er et stort Tankemiddel, saaledes som det
træffende er udtrykt i, hvad jeg ovenfor har kaldt Leibniz's
eksperimentale Definition af Identitet. Og ikke mindre væsentligt er
det, at vi gaar et Skridt længere tilbage og spørger, fra hvilket
Synspunkt Identiteterne gælder. Vi gaar dog derfor ikke glip af en
Ideverden. Tværtimod, de Synspunkter, ud fra hvilke Identiteter bliver
mulige, er for os den sande Ideverden, og det er en Betragtning,
Erkendelsesteorien afgiver til Drøftelsen af Tilværelsesproblemet (det
kosmologiske Problem), at den Kendsgerning, at der gives saadanne
Synspunkter, ved hvilke den strengeste Erkendelse, Mennesker kan op%
% --- p. 43 ---
\setcounter{footnote}{0}
naa, er mulig, nødvendigvis maa blive af stor Betydning for en sidste,
afrundende Verdensanskuelse. ---

For det Tredie gør netop Dannelsen af de kvalitative Identitetsrækker,
at Relationerne træder skarpere frem, end Tilfældet er før Udrensningen.
I de nævnte Rækker er hvert Led udelukkende bestemt ved sin Plads i
Forhold til de andre Led. Et Øjeblik faar, ganske bortset fra sit
Indhold, en individuel Karakter ved at følge efter visse Øjeblikke og
gaa forud for visse andre Øjeblikke. Øjeblikkene tjener gensidig til
Udmaaling af hverandres Individualitet. Det samme gælder Tallet: hvert
Tal i Rækken er en Individualitet, netop idet det udtrykker et bestemt
Forhold til andre Tal. Allerede Malebranche bemærker, at de hele Tal er
Forhold, ligesaavel som Brøker er det, skønt man ikke altid lægger
Mærke til det, fordi hele Tal kan udtrykkes ved et enkelt Tegn. Som
Gauss senere har udtalt: \guillemotleft Ethvert reelt helt Tal
repræsenterer Relationen af et vilkaarlig valgt Begyndelsesled til et
bestemt Tal i Rækken\guillemotright .%
% PRINTER'S DEFECT (3): »henægtes« as printed, for benægtes.
% »Russel« / »Russel's« with one l as printed, against »Bertrand Russell«
% earlier in the same note. Both readings kept.
% The printing sets the minus as a division sign (÷) and the fraction as
% m over 1, both in the roman face; given here in math mode.
\footnote{\textsc{Malebranche}: \textit{Recherche de la vérité.} VI, 1,
5. --- \textsc{Gauss} citeret hos \textsc{Cassirer}: \textit{Substanz und
Funktionsbegriff}. p.~72. --- Dette henægtes dog udtrykkelig af
\textsc{Bertrand Russell} (\textit{Introduction to Mathematical
Philosophy.} p.~64), for hvem Tallet m ikke er nogen Relation;
anderledes forholder det sig med $+$~m og $\div$~m. Et Tal er for
Russel's platoniserende Opfattelse en Eksistens i og for sig, bortset fra
de Relationer, i hvilke det ses. Derfor skelner Russel ogsaa bestemt
mellem $\frac{m}{1}$, der er en Relation, og selve m.} Den Kultus af
visse Tal, der er saa hyppig i Folkereligioner, blev kun mulig ved, at
man tog Tal ud af den Række, i hvilken de ene har deres Individualitet.
--- Hvad der gælder for Øjeblikke og Tal, gælder analogt for Grader og
Steder. Dette behøver i denne Sammenhæng ingen nærmere Paavisning. For
Stedets (Rummets) Vedkommende vil en anden, speciellere Sammenhæng give
Anledning til at
% --- p. 44 ---
\setcounter{footnote}{0}
komme tilbage til Spørgsmaalet. Af Interesse er dog netop her
Poincaré's Bemærkning%
% PRINTER'S DEFECT (4): the in-text reference is a superscript 1, but the
% note at the foot of the page is numbered 2 — and it is the only note on
% the page. Set here as note 1, matching the in-text mark; the printed 2
% is recorded here rather than reproduced.
\footnote{\textit{Science et Méthode}, p.~113.}, at Rummets Relativitet
og dets Ensartethed er en og samme Sag. ---

% Sub-subsection mark read off the page image at 600 dpi: GREEK GAMMA
% (γ), upright, roman face, followed by a full stop. Bare on the page —
% no printed title. Title below from the Indhold (p. 109), ToC only.
\addcontentsline{toc}{subsubsection}{γ.~Negationens og Rationalitetens
Relationer}%
\label{sec:3b-gamma}%
γ.~Til de formale Kategorier har vi foruden Identitet og Analogi ogsaa
regnet Negation og Rationalitet. Heller ikke for deres Vedkommende
behøves nogen udførlig Paavisning af, at Relation stadig ligger til
Grund. --- Ethvert Nej forudsætter et (muligt eller virkeligt) Ja. En
Negation er rent logisk Tegn paa en Vej, der ikke kan befares, og
forudsætter altsaa Forsøg paa at gaa en saadan Vej. Eller den
forudsætter en Sammenligning mellem to Anskuelsesbilleder, af hvilke det
ene har Elementer, der savnes i det andet. Matematisk kan Negation
betyde Opfordring til Rækkedannelse i modsat Retning af den tidligere.
--- Rationalitet betyder i sig selv en Relation, et Forhold mellem
Præmisser og Konklusion. Og saa langt som Fordringen om Begrundelse
gaar, saa langt maa der spørges om Forudsætningerne for vore Domme. En
ubetinget Nødvendighed gives ikke. Nødvendighed udtrykker kun Forholdet
mellem Grund og Følge, et Forhold, der stadig maa komme igen, saa længe
der tænkes, hvad enten det er muligt at finde det i det enkelte Tilfælde
eller ikke. Ligesom vi ved Klassifikationen stødte paa Analogibegrebet,
idet vi, hvor Artsbegreber ikke mere kan findes, standser ved Analogier
mellem de enkelte, individuelle Emner (forsaavidt saadanne Analogier
overhovedet endnu kan findes), saaledes stanser vi, hvad Bevisførelsen
angaar, ved den saakaldte Analogislutning. Aristoteles har træffende
kaldt den Slutning ud fra et Eksempel. Fra Forholdene indenfor et Emne
føres vi til at formode tilsvarende Forhold indenfor et andet Emne. Jo
flere Eksempler vi herved har at gaa ud fra,
% --- p. 45 ---
des mere nærmer Analogislutningen sig til Induktion. Men den faar stedse
kun (ligesom iøvrigt Induktionen) Betydning som Opdagelsesmetode,
d.~v.~s. som Anledning til at opstille Sætninger, der derefter nærmere
kan prøves; den kan aldrig blive Bevis. Derfor har man træffende kaldt
den en god Tjener, men en slet Herre. At ogsaa Analogien stedse beror
paa bestemte Forudsætninger, behøver ikke nærmere Paavisning.

% Head read off the page image at 600 dpi. The printing sets it
% LETTERSPACED (spærret) apart from the leading »c.«, which is not
% spaced; the full stop at the end is as printed. The head agrees with
% the Indhold word for word.
\subsection{Relationsbegrebet og de reale Kategorier.}
\label{sec:3c}

% Sub-subsection mark read off the page image at 600 dpi: GREEK ALPHA
% (α), upright, roman face, followed by a full stop, opening the first
% paragraph under the head. Bare on the page — no printed title. Title
% below from the Indhold (p. 109), ToC only.
\addcontentsline{toc}{subsubsection}{α.~Kausalitet, Realitet og
Relation}%
\label{sec:3c-alpha}%
α.~Det Spørgsmaal rejser sig naturlig, om det er muligt at anvende de
Rækker, som de formale Kategorier viser Muligheden af, paa Emner, der er
os givne aldeles bortset fra vore egne Konstruktioner, og det vil igen
sige, om der af saadanne Emner kan dannes Rækker, der i en eller anden
Grad eller paa en eller anden Maade svarer til hine formale Rækker og
derved kan faa Del i den Nødvendighed, som udmærker disse, altsaa om en
Rationalisering af givne Emner er mulig. Forholdet mellem Grund og Følge
er nu engang Forbilledet for al Nødvendighed.

Dette Spørgsmaal giver Anledning til en Drøftelse af de reale
Kategorier, nemlig Kausalitet, Totalitet og Udvikling, som jeg har
fremstillet i \guillemotleft Den menneskelige Tanke\guillemotright{}
p.~207---237, og her kun vender tilbage til for at se, hvilken Rolle
Relationsbegrebet paa deres Omraade spiller.

De reale Kategorier har historisk ikke ventet paa de formale. Snarere er
det dem, der har givet Stødet til Udformning af disse. Videnskabens
Udvikling ud af det praktiske Liv og den sunde Menneskeforstand
fremtræder her slaaende. De reale Kategorier, særlig Aarsagsbegrebet,
har deres faktiske Udspring i den Trang til Orientering i Tilværelsen,
som er en Livsbetingelse for Dyr og for Menne%
% --- p. 46 ---
% One unbroken paragraph; the page carries no footnote and no paragraph
% break at all.
% PRINTER'S DEFECT (5): »Forholdet« is broken across the fifth line from
% the foot with NO hyphen — »For« / »holdet«. Set here as one word.
sker. Tankens Udvikling af Tvivlen viser sig her slaaende. Naar
Skuffelse, Modstand eller Modsigelse gør Ende paa den uvilkaarlige
Tillid til alle opdukkende Fornemmelser og Forestillinger, er der ingen
anden Vej at gaa end den, at søge nye Fornemmelser og Forestillinger,
som kan frembyde en fastere Sammenhæng end den, der blev brudt, og som
kan gøre tryg Videreførelse af Livet mulig. Men dette vil sige, at
Relationernes Kreds udvides, indtil de forskellige Emner viser sig at
staa i faste, uundgaaelige Sammenhængsforhold til hverandre. Det er
Kausalitet eller lovmæssig Sammenhæng, der ligger til Grund ved al
Overbevisning om Realitet, som er mulig, efterat den uvilkaarlige Tillid
til den umiddelbart givne Sammenhæng er brudt. Dette praktiske
Virkelighedskriterium uddybes og præciseres i Videnskaben gennem
nærmere Prøvelse og Bestemmelse af de Relationer, som betinger
Sammenhængen mellem Emnerne. Ud over alle Relationer kommer Tanken ikke.
Om Virkeligheden af Noget, der slet ikke skulde staa i nogen Relation,
kan ingen begrundet Overbevisning være mulig. Virkelighed og Relation
hører nu engang sammen. Enhver Konstatering af, at noget er virkeligt,
bestaar i Paavisningen af en fast Relation mellem dette Noget og den
Kreds af Emner, der hidtil staar som virkelige gennem indbyrdes fast
Sammenhæng; Tankens ene Passerben staar i denne Kreds, og med det andet
Passerben udmaales Forholdet mellem denne Kreds og det, om hvis
Virkelighed der spørges. Rejses der Tvivl om den Virkelighed, til
hvilken man hidtil holdt sig, maa ogsaa det første Passerben flyttes, og
saaledes fremdeles. Der vil stadig blive Brug for begge Passerben ved
Virkelighedserkendelse. Nødes man til at staa paa det ene, uden at kunne
finde sikker Plads til
% --- p. 47 ---
\setcounter{footnote}{0}
det andet, er man i en Tilstand af Tvivl, Forventning eller
Længsel.%
% PRINTER'S DEFECT (6): the note ends »men ogsaa p. 96 f).« — a closing
% parenthesis with no opening one. Transcribed as printed.
\footnote{Thomas af Aqvino forsøgte at befri det teologiske Gudsbegreb
fra al Relation ved at hævde, at Verden vel staar i Relation til Gud,
men Gud ikke til Verden. (Smlgn. herom min \textit{Religionsfilosofi}
§~22). Thomas sluttede sig her til Mystiken, særlig til Hugo a St.
Victore, der igen er paavirket af Nyplatonismen. Ud fra denne Tankegang
har en moderne katolsk Teolog kritiseret Spinoza, fordi han ikke har
vovet at gøre det Skridt at antage, at der til Verdens Afhængighed af
Gud ikke svarer nogen Afhængighed for Gud af Verden.
(\textsc{Dunin-Borkovski}: \textit{Der junge Spinoza}, p.~352). Men en
Relation forudsætter nu altid to Led, og hvorledes man end stiller sig
til Spinoza's Gudsbegreb, maa man anerkende den Konsekvens, med hvilken
han fastholder Relationsbegrebet paa dette Punkt. Smlgn. min Afhandling
\textit{Spinoza's Ethica}, p.~26; 32---35; men ogsaa p.~96 f).}

Paa Grund af den nøje Forbindelse, der er mellem Kausalitet og
Virkelighedskriterium, og som viser sig i, at vi kun har begrundet
Overbevisning om Virkelighed, hvor vi kan forbinde Emner som Aarsag og
Virkning, var det et revolutionært Skridt, Hume gjorde ved at betvivle
Retten til at anvende Aarsagsbegrebet. Han rokkede baade ved den
videnskabelige Erkendelse af Tilværelsen og ved det praktiske
Virkelighedskriterium. Men hans Kritik hvilede paa den Forudsætning, at
Aarsag og Virkning var to helt forskellige Ting, --- Led i en kaotisk
Forskelsrække. Her mangler netop den faste Sammenhæng, da man kan stille
Leddene i en saadan Række i en hvilkensomhelst Orden. Nu er det en
Kendsgerning, at vi dog undertiden staar overfor Rækker, hvis Led ikke
kan byttes om, og det var i denne Kendsgerning, Kant tog sit
Udgangspunkt som Grundlag paa én Gang for Aarsagssætningen og for
Virkelighedsbegrebet. Og Videnskabens Historie viser, hvorledes det
udvortes Forhold mellem Aarsag og Virkning paa mange Omraader under
Undersøgelsens Gang mere og mere giver Plads for et Kontinuitetsforhold
mellem de Emner,
% --- p. 48 ---
\setcounter{footnote}{0}
der staar i Kausalitetsforhold til hverandre, og dette igen maaske for
et Rationalitetsforhold, idet senere Led i Begivenhedernes Række kan
udledes af Loven for de foregaaende Leds Rækkefølge.

Om de Rækker, der kan dannes paa denne Maade, end bliver nok saa
kontinuerlige og rationelle, bliver Virkelighedsbegrebet dog ved at staa
som et Ideal, der aldrig fuldt ud kan belægges med Iagttagelse. Der vil
stadig kunne findes nye Mellemled, og der vil kunne opdukke nye Emner,
hvis kontinuerlige og rationelle Forhold til de tidligere bliver at
paavise. Virkelighed betyder stadig en Relation til hidtil dannede
kontinuerlige og rationelle Rækker. Og der kan aldrig gives nogen
Garanti imod, at disse Rækker ikke maa erstattes ved helt andre Rækker,
eller for at blive ved et ovenfor brugt Billede, at begge Passerben maa
flyttes.

Herbert Bradley har med Rette benyttet Erfaringsbegrebet ved
Tilværelsens for os uafsluttelige Karakteristik. Erfaring er for Bradley
Maalestok for Realitet, men da Erfaring stadig vindes ved Paavisning af
Relationer, og nye Relationer stedse kan gøre sig gældende, er en
fuldstændig Fyldestgørelse af hin Maalestoks Krav umulig. Afsluttet
Harmoni og Relation er uforenelige Begreber. Bradley er paa én Gang
overbevist om, at vi altid tænker i Relationer, og at en fuldstændig
Løsning af Virkelighedsproblemet er umulig ud fra Relationsbegrebet (the
relational point of view).%
% The English tag and the English matter in the note are both set ROMAN
% in the printing, not italic; only the book title is italic. Left roman.
\footnote{\textit{Appearance and Reality}\textsuperscript{1}, p.~171
(When thought begins to be more than relational, it ceases to be mere
thinking); cfr. p.~445; 519.}

Løsningen kan, deri beholder Bradley Ret, aldrig naas fuldt ud. Men
Erkendelsesarbejdet gaar stadig ud paa at finde nye Relationer og søge
at bringe dem i Sammenhæng med tidligere fundne Relationer, maaske bruge
dem til Be%
% --- p. 49 ---
\setcounter{footnote}{0}
% The page foot also carries the printer's signature line »Vidensk.
% Selsk. Filosof. Medd. I. 3.« with the signature numeral 4; omitted as
% running matter.
rigtigelse af disse. Uden her at komme ind paa, hvad der hører under det
kosmologiske Problem, maa det hævdes, at Erkendelsesarbejdet selv er en
Virkelighed, et Led i den store Virkelighed, som maaske er langt mere
omfattende, end hvad noget Tankearbejde kan omspænde. Og naar den Del af
Virkeligheden, Erkendelsesarbejdet udgør, aldrig kan afsluttes, kunde
man maaske deraf uddrage den Formodning, at Tilværelsen idethele ikke er
færdig eller kan blive færdig. En saadan
Formodning\footnote{En saadan Formodning udtalte jeg allerede i mit
Universitetsprogram af 1902 (\textit{Filosofiske Problemer}) p.~61 og
kort efter i Foredrag i Amerika (Mindre Arbejder. Tredie Række:
\textit{En filosofisk Bekendelse}, p.~26 f.).} forudsætter ganske vist,
at Tidens Form gælder for Tilværelsen, ikke blot for vore Tanker, skønt
der til alle Tider, ogsaa nu, er Forskere, der ikke antager dette. Hvad
Bradley angaar, nægter han bestemt, at Tanken er en Del af Realiteten:
The process which moves within Reality, is not Reality itself.
Philosophy is itself but
appearance.\footnote{\textit{Appearance and
Reality}\textsuperscript{1}, p.~448; 454.}

Bradley's dybsindige Hovedværk er Kosmologi (\guillemotleft
Metafysik\guillemotright ), ikke egentlig Erkendelsesteori. Han stiller
Realitet i Modsætning til Sandhed: naar Sandheden bliver fuldkommen,
hører den op at være Sandhed, ti saa træder Realiteten i sin Renhed ind.
Af dette Standpunkt følger, at skønt han lægger saa stor Vægt paa
Relationsbegrebet, at han definerer det at tænke ved at finde
Relationer, saa opfatter han dog Tankens Søgen efter Relationer som en
Ufuldkommenhed, fordi Afslutning og Totalitet derved udelukkes. Og
hermed hænger det igen sammen, at han miskender Betydningen af de
Relationer, med hvilke de specielle Videnskaber hver paa sit Omraade
arbejder. Han har klart set Relationsbegrebets Betydning som Grundlag og
Maale%
% --- p. 50 ---
\setcounter{footnote}{0}
stok for al Søgen efter Sandhed. Men hans eget System er i Grunden et
krystalliseret Sandhedsideal, et Forsøg paa at tænke Erfaringen som
fuldstændig (experience entire, containing all elements in harmony,
p.~172), --- et Forsøg, der umulig kan gennemføres, netop naar man som
Bradley opfatter Tænkningen som relationel.%
% »Haldane« is small-capped at its first occurrence in the note and
% broken HAL-/DANE'S across the line; roman at its second occurrence.
\footnote{I samme Retning som den her givne Kritik af Bradley gaar
\textsc{Haldane}'s kritiske Bemærkninger (\textit{The Reign of
Relativity}, 1921, p.~203---211) overfor denne Tænker, han i andre
Henseender staar saa nær. Haldane's Værk kom mig først i Hænde, efterat
nærværende Afhandling var fuldendt.} ---

% The printing sets a blank line here — extra leading after the dash,
% marking a break within the sub-subsection. No head.
Paavisning af kontinuerlige Emnerækker med uombyttelige Led er endnu
ikke Videnskabens højeste Ideal. Den raader vel foreløbig Bod paa det
udvortes Forhold mellem, hvad man kalder Aarsag, og hvad man kalder
Virkning, som fremkaldte Hume's Skepsis; men den gør ikke fuldstændig
Ende derpaa. Et Skridt videre fører Paavisningen af tvesidig Ækvivalens
mellem forskellige Emner. Der kan da være Mulighed for baade at udlede
det ene Emne af det andet og dette andet af det første. Hvor
Ækvivalensforholdet formuleres som en Ligning, fremtræder dette tydelig.
Det er (i Teorien) ligegyldigt, fra hvilken Side af Lighedstegnet man
begynder, og den kvalitative Forskel, saavel som Tidsforskellen mellem
de to Emner har da ikke længere nogen Betydning; et absolut
Identitetsforhold er naaet. Kvaliteterne (til hvilke ogsaa Tiden,
Successionen kan regnes) staar da som blotte \guillemotleft
Antropomorfismer\guillemotright , subjektive Emner, som den strenge
Videnskab ikke bryder sig om. Man gaar her et Skridt videre end den
mekaniske Naturopfattelse, der ogsaa saa bort fra Kvaliteter i snevrere
Forstand, men endnu anerkendte Tiden som objektiv Form, idet den saa
Bevægelse i Rummet som
% ============================================================
%  BATCH 5 — §3c concludes; §3d Relationen Subjekt—Objekt (head at p. 58)
%  Heads listed in the Indhold at pp. 55 and 58 — find them.
% ============================================================
% --- p. 51 ---
\setcounter{footnote}{0}
det Eneste, der eksisterede, skønt det blev \guillemotleft
forfalsket\guillemotright{} ved Menneskenes subjektive Opfattelse.%
% The note closes with a Latin couplet, set on two lines of its own below
% the prose, indented well in from the note's measure. Only »falsificata«
% is italic; »Thomas Hobbes« is small-capped, »Galilei« is not.
\footnote{\textsc{Thomas Hobbes}, der dels ad sin egen Vej, dels under
personlig Paavirkning af Galilei, blev en af de første Forkæmpere for
strengt mekanisk Naturopfattelse, udtrykker i sin poetiske Selvbiografi
sin vundne Indsigt (Bevægelse som det Eneste) saaledes:\par
\centerline{Et mihi visa qvidem est toto res unica Mundo}
\centerline{Vera, licet multis \emph{falsificata} modis.}}
Det er et betydningsfuldt og paradoxt Skridt, der gaas videre, naar man
antager, at Tiden kan vendes om, saa at Begivenhedernes Rækkefølge
tilsidst er ligegyldig.

I Modsætning til en saadan Opfattelse, der ofte hævdes fra matematisk og
naturvidenskabelig Side, og som i filosofisk Form er gennemført af
Marburgerskolen, maa det hævdes, at den Verden af rene Ligninger, i
hvilken streng Videnskab ender, eller som den i hvert Tilfælde stiler hen
imod, ikke mister sin Betydning, fordi det stadig fastholdes, at et
Forhold mellem før og efter aldrig fuldstændig kan gaa op i et Forhold af
ren Identitet. Fordi der er tvesidigt Ækvivalensforhold mellem to
kvalitativt forskellige Emner (som Varme og Bevægelse), er det dog ikke
ligegyldigt, hvilket Emne der kommer først, og hvilket der følger efter.
Der sker noget forskelligt i de to Tilfælde, hvilket tydelig vil vise
sig, naar Rækken føres videre gennem nye Led: selv om der er
Ækvivalensforhold mellem A og B, vil der fra B kunne være en Overgang
mulig til C, som ikke vilde være mulig fra A. Desuden foregaar (hvad
Carnot's Sætning viser med Hensyn til Forholdet mellem Varme og
Bevægelse) Overgangen faktisk ikke lige let i begge Retninger. Men trods
Alt har hin strengt matematiske Afslutning af Naturerkendelsen sin store
Betydning derved, at der til det rent tidløse, logisk-matematiske Forhold
mellem Ligningernes Led svarer den tidslige Over%
% The signature mark "4*" stands alone at the foot of p. 51, below the
% note; a printer's mark, not transcribed.
% --- p. 52 ---
gang mellem kvalitativt forskellige Emner i vor Oplevelse. Det er gennem
Analogi, at den logisk-matematiske Erkendelse faar sin uvurderlige
Betydning. I Ligningerne kan den kyndige aflæse Fænomenernes Forløb,
ligesom Musikeren i Partituret kan aflæse Tonernes Løb.

Denne Opfattelse antydedes allerede af Kant i hans Lære om, at der mellem
Grund og Følge paa den ene Side, Aarsag og Virkning paa den anden Side,
kun bestod et Analogiforhold, ikke et Identitetsforhold. Den er senere
især paa betydningsfuld Maade gjort gældende af \textsc{Clerk Maxwell} (i
andet Bind af hans \guillemotleft Scientific Papers\guillemotright ). I
mit Universitetsprogram 1902 (\guillemotleft Filosofiske
Problemer\guillemotright ) og i \guillemotleft Den menneskelige
Tanke\guillemotright{} (1910) har jeg fremsat den samme Opfattelse, som
jeg i \guillemotleft Totalitet som Kategori\guillemotright{} (1917) har
udviklet nærmere i Polemik mod Marburgerskolen.

At Forholdet mellem en Række af Ligninger og en Række af kvalitative
Forandringer ikke er et logisk Identitetsforhold, ses deraf, at man af
den rene Logik og Matematik aldrig vilde kunne udlede, at der gives en
tidslig Verden med kvalitative Forskelle. Det gaar med hin Ligningernes
Verden som med Platon's Ideverden (der i Grunden fremtræder her i ny
Skikkelse): man kan maaske stige op til den, men Stigen falder bort, saa
snart man er kommen derop. Og dertil kommer endnu, at et afsluttende
Bevis for de absolute Identitetsrækkers reale Gyldighed ikke kan gives,
da der stedse er Mulighed for nye Relationer, hvis Analogier med de
fundne Identitetsrækker først skal godtgøres. Atter her viser det sig, at
Relationsbegrebet gør Sandhed og Virkelighed til Idealer, der kun under
stadigt Arbejde tilnærmelsesvis kan naas.

\textsc{Émile Meyerson} har i sine erkendelsesteoretiske Værker lagt
meget stor Vægt paa Forskellen mellem Erkendelse af
% --- p. 53 ---
et lovmæssigt Forhold mellem Emner og en rationel Forklaring af et
saadant Forhold. Især i det nylig udkomne store og dybtgaaende Værk
\guillemotleft De l'Explication dans les Sciences\guillemotright{} (1921)
søger han at belyse denne Forskel gennem en Rigdom af Eksempler fra
Naturvidenskaben og dens Historie, en Rigdom, som enhver, der beskæftiger
sig med Erkendelsesteori, maa misunde ham. For mit Vedkommende er jeg i
Hovedpunktet enig med ham, skønt min Terminologi er en noget anden end
hans. Jeg har allerede i det Afsnit i min Psykologi, der handler om
Overgangen fra Psykologi til Erkendelsesteori, hævdet, at Videnskaben
ikke bliver, og ikke kan blive staaende ved den blotte Konstatering af en
uundgaaelig og uombyttelig Rækkefølge af to Emner, men at Opgaven,
efterat en saadan Rækkefølge (og dermed en \guillemotleft
Lov\guillemotright{} i Ordets populære Betydning) er paavist, bliver den
at finde en kontinuerlig Overgang fra det ene Led i Aarsagsforholdet til
% Printer's defect: the printing sets a full stop after »det andet«,
% where the sense (and the lower-case »saa« that follows on the next
% line, within the same paragraph) requires a comma. Read at 600 dpi:
% a round full stop, not a comma. Transcribed as printed.
det andet. saa at det, vi kalder Virkningen, kan komme til at staa som en
Fortsættelse af det, vi kalder Aarsagen, en Fortsættelse, der danner en
Analogi til den Maade, paa hvilken i en Slutning Konklusionen
\guillemotleft følger\guillemotright{} af Præmisserne. Det er det store
Ideal, det 17. Aarhundredes Tænkere opstillede, naar de krævede, at der
ikke maatte være mere i Virkningen, end der var i Aarsagen. Den
Kontinuitet, der findes, naar dette Ideal til en vis Grad naas, er,
saavidt jeg forstaar, det, som Meyerson kalder Rationalitet, og som naas
ved \guillemotleft Forklaring\guillemotright . I et meget lærerigt
Eksempel viser han (I p.~315), hvorledes Videnskaben ikke bliver staaende
ved at paavise et lovmæssigt Forhold mellem visse optiske Fænomener og et
Stofs kemiske Struktur, men at man søger at finde en sammenhængende
Proces, der begynder med Lysets Indvirkning, som er en vis Art af
Bevægelse,
% --- p. 54 ---
og fortsættes gennem de Bevægelser, som denne Indvirkning fremkalder i
Molekulernes Indre, og som saa igen fremkalder visse Bevægelser i de
Sanseorganer, der paavirkes af vedkommende Stof. Det er ifølge dette
Skema en kontinuerlig Række af Bevægelser, der søges konstaterede. Men
forsaavidt det er muligt at gennemføre dette Skema (og Meyerson
fremhæver meget stærkt Vanskelighederne derved), saa lader det sig kun
gøre gennem Paavisning af en Række af Love (for Lysets Refleksion, for
Molekulers indre Bygning o. s. v.), og jeg ser ikke, hvilken Modsætning,
der saa tilsidst bliver mellem \guillemotleft
Lovmæssighed\guillemotright{} (légalité) og \guillemotleft
Rationalitet\guillemotright . Meyerson ser da ogsaa selv klart, at enhver
Loverkendelse, ofte ubevidst for Opdageren, forudsætter et indre Baand
mellem de Emner, der knyttes sammen i Kraft af Loven (II p.~289), og han
indrømmer (II p.~337), at Legalitet og Rationalitet i Virkeligheden er
saa nøje (inextricablement) forbundne, at det kun er ved Eftertanke over
Videnskabens Grundlag, at man kan holde de to Principer ude fra
hinanden. For mit Vedkommende antager jeg ikke to Principer her; det er
et og samme Princip, der gennemføres med stigende Nøjagtighed og
Konsekvens.

Jeg tror derfor heller ikke, at Meyerson har ganske Ret i sin Kritik af
Comte, naar han formulerer denne Kritik saaledes, at Comte vel anerkender
Lovmæssighed, men ikke Rationalitet som Videnskabens Opgave. Comte selv
udtaler bestemt (i Discours sur l'esprit positif, p.~20---21), at
Videnskabens Opgave baade er Forklaring og Forudseen, og at den
\guillemotleft ved sine systematiske Spekulationer imødekommer vor
Forstands uvilkaarlige Enhedstrang ved at gennemføre vore forskellige
Begrebers Kontinuitet og Ensartethed\guillemotright . En anden Sag er
det, at han ikke mente, at der var Ud%
% --- p. 55 ---
% The sign before »til 1830« is the Danish »det er« mark: a turned c
% (ɔ) followed by a colon. Read at 600 dpi; the text layer gives it as
% a box. Transcribed with U+0254.
sigt til Virkeliggørelse af dette Ideal, dels fordi de faktiske
Diskontinuiteter forekom ham uovervindelige, dels fordi han vilde være
Religionsstifter og gøre den hidtil (ɔ: til 1830) vundne Videnskabs
Resultater til den nye Religions Dogmer. Kritiken af Comte maa efter min
Opfattelse formuleres saaledes: han krævede en Afslutning af
Loverkendelsen og saa ikke Muligheden af at gennemføre den i endnu
speciellere Former og Grader, end det hidtil var gjort. ---

Der er endnu et Punkt, i hvilket jeg ikke er enig med den udmærkede
franske Tænker, af hvis Værker jeg har lært saa meget. Jeg tror ikke, at
Videnskabens sidste Maal er Opstilling af almindelige Sætninger,
formulerede saa eksakt som muligt. Der er endnu et Arbejde at gøre,
nemlig at benytte de almindelige Sætninger, de rationelle Synspunkter til
Forstaaelse af de individuelle Eksistenser, der kun bestaar formedelst et
Sammenspil, en Samfoldighed af Lovmæssigheder. Aristoteles havde Uret i,
at det Individuelle ikke kunde være Genstand for Videnskab. Snarere
forholder Sagen sig saaledes, at det uhyre Arbejde, Videnskaben gør for
at finde Lovmæssighed og Rationalitet, selv igen er Betingelse for at
vinde Forstaaelse af de individuelle Helheder, som er det faktisk
Eksisterende. Først ved ikke at glemme denne Opgave aabnes Blikket for
Videnskabens store, uendelige Maal.

% Sub-subsection mark read off the page image at 600 dpi: GREEK BETA
% (β), upright, roman face, followed by a full stop, at the head of an
% ordinary indented paragraph. Bare on the page — p. 55 carries NO
% printed title. The title below is the Indhold's (p. 109), recorded in
% the ToC only, never printed into the text.
\addcontentsline{toc}{subsubsection}{β.~Totalitetens og Udviklingens
Relationer}%
\label{sec:3c-beta}%
β.~Dette minder om, at foruden Kausalitet maa ogsaa Totalitet og
Udvikling regnes til de reale Kategorier --- netop i Kraft af
Relationsbegrebet. Thi de fundne rationelle Synspunkter maa forbindes
indbyrdes og stilles derved i bestemte Relationer til hverandre, og
gennem Arbejdet herpaa finder Tanken nye Totaliteter, hvor der forhen kun
syntes at foreligge sporadiske Iagttagelser. Men enhver
% --- p. 56 ---
\setcounter{footnote}{0}
saadan Totalitet staar i Forhold til Emner og Elementer udenfor sig, og
der stilles derfor --- i Kraft af Relationsbegrebet --- stedse nye
Opgaver. Naar derimod givne Emner umiddelbart optræder som Helheder,
bliver Opgaven ikke syntetisk, men analytisk; det gælder at finde den
Sammenhæng, det Sammenspil, der gør saadanne Helheder mulige. Hvad man
kalder \guillemotleft Ting\guillemotright , umiddelbart givne Helheder,
forstaas kun gennem deres Egenskaber, og Egenskaberne betegner lige saa
mange Relationer til andre \guillemotleft Ting\guillemotright . Molekuler,
Atomer og Elektroner er stadig \guillemotleft Ting\guillemotright ,
Helheder, der kun kendes og forstaas gennem deres Relationer.%
\footnote{En nærmere Udredning af det her Antydede er given i min
Afhandling \textit{Totalitet som Kategori.}}

Udviklingsbegrebet kommer naturlig frem, naar man lægger Vægt paa, at
Kausalitet betyder kontinuerlig Proces; der kan da her skelnes
forskellige Stadier, indtil en vis Afslutning foreligger, og forsøges en
Karakteristik af disse forskellige Stadier. Men dernæst hænger
Udviklingsbegrebet nøje sammen med Totalitetsbegrebet, idet enhver Helhed
har sin Historie. --- Hvad der i det enkelte Tilfælde kaldes en Helhed
eller et Udviklingsstadium, vil bero paa Synspunktet. Hvad der for én
Betragtning staar som Helhed, staar maaske for en anden Betragtning som
en Del; hvad der for én Betragtning staar som et Udviklingsstadium, staar
maaske for en anden Betragtning som et Opløsningsstadium. En nærmere
Indgaaen paa de her opstigende Spørgsmaal vilde høre ind under Drøftelsen
af Relationen Subjekt---Objekt, eller ogsaa under Drøftelsen af
Relationsbegrebet paa Værdibegrebernes Omraade og maa opsættes til senere
Afsnit. ---

Karneades har faaet Ret. Vi kender kun Virkeligheden
% --- p. 57 ---
\setcounter{footnote}{0}
gennem de faste og mere eller mindre gennemskuelige Relationer mellem de
Emner, Iagttagelsen frembyder. Den moderne Videnskabs Fremskridt i
Forhold til de gamle Akademikere og Skeptikere beror paa den Nøjagtighed,
hvormed de enkelte Relationer bestemmes, paa det stigende Antal
Relationer, der paavises, og paa det Sammenspil af disse Relationer, der
opdages. Derved er Virkelighedsbegrebet Skridt for Skridt blevet udviklet
og udvidet. Akademikerne saa ikke de Udviklingsmuligheder,
Relationsbegrebet frembød. De kunde ikke se med Humor paa Erkendelsens
Ufuldkommenheder, fordi de ikke saa Betydningen af fremskridende
Anvendelse af Relationsbegrebet. En særegen Form for Humor er derimod i
nyere Tid bleven mulig, idet Erkendelsens store Ideal holdes fast uden
Overvurdering af hvert Skridt, der gøres.%
\footnote{Smlgn. \textit{Den store Humor.} Kap. 6 (Forstaaelse og Humor)
og Kap. 9 § 54 (Humor og kritisk Realisme).}

Men dertil kommer endnu, at bag Akademikernes skeptiske Holdning laa i
Grunden stedse det populære Virkelighedsbegreb, ifølge hvilket
Virkeligheden en Gang for alle staar fast som en absolut Tingenes Orden,
med hvilken vore Tanker maa stemme, for at vor Erkendelse kan blive sand.
Det er en umulig Relation, der her antages. Virkeligheden kan ikke selv
være Led i en Relation; det vilde jo netop forudsætte, at vi allerede
kendte den. Virkelighed beror derimod selv paa, eller kundgør sig i
Relationerne mellem vore Iagttagelser. Vor Tanke kan ikke sætte sit ene
Passerben paa sig selv og det andet paa noget, der endnu slet ikke er
tilgængeligt for den. Sandhed er ikke Afsløring, men Frembringen, vindes
kun gennem stadigt Arbejde.

% --- p. 58 ---
% Head read off the page image at 600 dpi. The printing sets it
% LETTERSPACED (spærret) apart from the leading »d.«, which is not
% spaced; the full stop at the end is as printed. The dash between
% Subjekt and Objekt is an EM DASH, set tight against both words (the
% white around it is the letterspacing of the head, not word space).
% The Indhold (p. 109) sets the same head with the dash spaced,
% »Relationen Subjekt --- Objekt«; the printed head is followed here.
\subsection{Relationen Subjekt---Objekt.}\label{sec:3d}
% Sub-subsection mark read off the page image at 600 dpi: GREEK ALPHA
% (α), upright, roman face, followed by a full stop. UNLIKE the marks
% on pp. 33, 40, 44 and 45, this one is NOT bare: the page prints the
% title »Indledning.« after it, run in with the text, which then
% continues after an em dash on the same line. The head is set as a
% heading here and the run-in dash dropped. The Indhold's wording
% (»Indledning«, no full stop) is what goes into the ToC.
\subsubsection*{α.~Indledning.}
\addcontentsline{toc}{subsubsection}{α.~Indledning}%
\label{sec:3d-alpha}%
\noindent De hidtil omtalte Relationer bestaar mellem Emner indbyrdes og
udtrykker fra forskellige Sider den Kendsgerning, at den menneskelige
Tanke arbejder ved at stille et Emne i Forhold til andre Emner og ad den
Vej finder den Forstaaelse, den kan vinde. Det gælder tilsidst om at føre
alle Emner tilbage til visse Grundemner, der foreløbig staar fast, selv
om en senere Tids Forsken maaske søger at naa længere tilbage. Som
\textsc{Ernst Mach} (Erhaltung der Arbeit. 1872) har udtrykt det: at man
netop standser ved visse bestemte Emner, kan kun forklares historisk. Han
anfører et stort Eksempel (S. 32): \guillemotleft Naar vi nu til Dags
tror, at mekaniske Kendsgerninger er forstaaeligere end andre
% The two ellipses on this page are printed as three spaced full stops,
% both read at 600 dpi; transcribed as printed.
Kendsgerninger, . . . er det en Illusion. Det beror derpaa, at Mekanikens
Historie er ældre og rigere end Fysikens, og at vi i længere Tid har
staaet paa en fortrolig Fod med mekaniske Kendsgerninger. Hvem tør
paastaa, at elektriske og termiske Fænomener ikke en Gang vil staa for os
paa samme Vis, naar vi har lært deres simpleste Regler at kende og er
blevne fortrolige med dem? . . . En Grundkendsgerning er aldrig
forstaaeligere end en anden\guillemotright . --- At Grundkendsgerningernes
(eller, som jeg kalder dem, Grundemnernes) Plads i vor Erkendelse kun kan
forklares historisk, vil sige, at det er gennem Videnskabens Historie, vi
kan komme til at forstaa, hvorfor man til given Tid netop maatte stanse
ved visse bestemte Kendsgerninger og finde fuld Forklaring ved at føre
andre Kendsgerninger tilbage til dem. Mach har i den anførte Udtalelse
forudsagt, hvad der virkelig er sket. I det tyvende Aarhundrede betragter
Naturvidenskaben ikke mere Inerti, Tyngde og Energibestaaen som
Grundkendsgerninger. Den stiller sig i hvert Tilfælde den Opgave at føre
dem tilbage
% --- p. 59 ---
\setcounter{footnote}{0}
til andre Kendsgerninger. --- I Filosofien har man ofte villet slaa sig
til Ro ved et vist Antal Grundbegreber (Kategorier), og selv Kant mente
her at være naaet til Enden. Men det viser sig, at de Grundbegreber,
Erkendelsen benytter under sit Arbejde ligesaa lidt som
Naturvidenskabens Grundkendsgerninger er de samme til alle Tider, selv om
de vel nok frembyder en vis gennemgaaende Type gennem Tiderne.

Det er forstaaeligt, at Forskerne midt i deres energiske Arbejde ikke
samtidig altid har den historisk betingede Karakter af de Grænser, til
hvilke de naar, klart for Bevidstheden. Meyerson udtrykker dette
saaledes: Videnskaben er ontologisk, kun at den sætter en anden Ontologi
i Stedet for den sunde Menneskeforstands, hvad enten nu den nye Ontologi
fremtræder som Atomistik, som Energetik eller som ren Matematik.%
% »dévanciers« is printed with the acute, for standard French
% »devanciers«; left as printed.
\footnote{La science de nos jours est saturée d'ontologie, et les
savants, en dépit de ce qu'ils affirment expressément eux-mêmes, font de
la métaphysique comme ont fait, de tout temps, leurs dévanciers.
\textit{De l'Explication dans les Sciences.} II, p.~174.}
 --- Det er dog ikke nødvendigt, at naiv Ontologi stedse skal
karakterisere videnskabeligt Arbejde. For det første opløser, som
Meyerson selv viser%
\footnote{La science, en menant ses explications jusqu'au bout, finit par
détruire cette ontologie qui, d'abord, paraissait lui être indispensable.
I, p.~50.}, Videnskaben selv den Ontologi, der foreløbig syntes
nødvendig. For det andet maa jo dog naturlig det Spørgsmaal dukke op,
selv for den mest begejstrede Forsker, med hvilken Ret han nu mener at
have naaet \guillemotleft Tingen i sig selv\guillemotright ; det er ikke
noget fremmed, udenfor Videnskaben selv liggende Spørgsmaal; derom vidner
selve Videnskabens Historie tilstrækkelig. Og for det Tredie maa det
Spørgsmaal uundgaaelig komme frem, hvorledes man ud fra den Ontologi, man
er naaet til, tilbagegaaende kan udlede den
% --- p. 60 ---
Verden af Kvaliteter og Successioner, fra hvilken man steg op til, hvad
man mente at have naaet som den absolute Tingenes Orden. Dette Spørgsmaal
har samme Berettigelse nu, som det i Oldtiden havde overfor Platon og
Demokritos. ---

Men foruden de objektive Relationer, gennem hvilke Emner belyses og
bestemmes i Forhold til hverandre, gør der sig stadig, skønt mere skjult,
en Relation gældende mellem det erkendende Subjekt og de Emner, der er
dets Objekter. Allerede de antike Skeptikere indsaa, at her forelaa en
Grundrelation. (Se ovenfor p.~9). Alle objektive Relationer opfattes og
gennemtænkes af et menneskeligt Subjekt og gælder, foreløbig, i Relation
til det. Med en Forandring af dette Subjekt, dets Organisation og dets
Kaar, maa ogsaa dets Erkendelse kunne blive en anden. Her, ligesom ved de
objektive Grundemner, er der en historisk Opgave, den nemlig, at paavise,
hvorfor Erkendelsessubjektet netop møder med visse Forudsætninger og
stiller visse Spørgsmaal. Denne Opgave vedkommer den sammenlignende
Psykologi og Videnskabens Historie. Her har vi kun at gøre med det
almindelige Synspunkt: uden Subjekt intet Objekt.

Det omvendte Synspunkt gælder dog ogsaa: uden Objekt intet Subjekt. Vi
har nemlig aldrig et rent Subjekt uden objektivt Indhold og uden
Bestemthed og Betingethed ved dette Indhold, ligesaa lidt som vi
nogensinde har et rent Objekt, men stedse (om ikke med fuld Bevidsthed)
forudsætter en mulig Iagttager eller Tænker. Begreberne Subjekt og Objekt
er korrelative, ligesom Syntese og Relation, Kontinuitet og
Diskontinuitet, Lighed og Forskel, o. s. v. --- Enhver Fremstilling af en
Verdens Tilblivelse, ligefra den mosaiske Skabelseshistorie til moderne
kosmogoniske For%
% --- p. 61 ---
\setcounter{footnote}{0}
søg, giver en Beskrivelse af, hvorledes det Hele vilde udfolde sig for en
Tilskuer med visse Sanser og visse Forudsætninger. Det erkendende
Subjekt, Tilskueren, Tænkeren, har selv en Historie, som er bestemt ved
de Erfaringer, han har gjort, og den Maade, paa hvilken han har
bearbejdet disse Erfaringer. Den moderne Erkendelsesteori er klar over,
at der gennem videnskabeligt Arbejde kan udvikle sig en Subjektivitet med
Evne og Trang til at se Emnerne i deres Relationer og med Forstaaelse af,
hvorledes den selv er naaet til det Stade, fra hvilket den stiller sine
Opgaver og fordrer Svar paa sine Spørgsmaal.%
\footnote{Om Vekselvirkningen mellem den psykologiske og den
erkendelsesteoretiske Betragtningsmaade paa dette Punkt se \textit{Den
menneskelige Tanke}, p.~297---303.}
 Man kan kalde en saadan Subjektivitet det erkendelsesteoretiske Subjekt.
% »Heinrich Rickert« is set in ordinary roman here, not small caps —
% checked at 600 dpi against the small-capped »Ernst Mach« on p. 58.
Heinrich Rickert, der klart har belyst dette Punkt, kalder træffende det
erkendelsesteoretiske Subjekt et Grænsebegreb i Forhold til det
psykologiske Subjekt.%
% The raised 2 is an edition mark on the title, not a note reference.
\footnote{\textit{Grenzen der naturwissenschaftlichen
Begriffsbildung}\textsuperscript{2}, p.~134---138; 308.}
 At ingen Kategorilære kan være definitiv, følger netop af, at dette
Grænsebegreb stedse kan og maa bestemmes paany, fordi det udvikler sig i
Forhold til det objektive Indhold, der til hver Tid er tilstede. Verden
forstaas ud fra Menneskets Natur og Stade, men Menneskets Natur og Stade
forstaas igen ud fra Verdenslovene og Verdensudviklingen. Ud over dette
Vekselspil kommer man ikke, selv om de to Synspunkter til forskellige
Tider kan optræde med forskellig Vægt i den videnskabelige Diskussion.

Den Vægt, der lægges paa Relationen Subjekt---Objekt, fører ingenlunde
til Subjektivisme, da der, som sagt, ikke gives noget rent Subjekt, som
skulde kunne frembringe alt ud af sig selv. Der gives stedse kun et
objektivt be%
% ============================================================
%  BATCH 6 — §3d continues
%  Heads listed in the Indhold at pp. 62 and 67 — find them.
% ============================================================
% --- p. 62 ---
% p. 62 CONTINUES the paragraph running from p. 61 — its first line is set
% flush to the left margin, not indented. No blank line before this marker.
\setcounter{footnote}{0}
stemt S (et S\textsubscript{o}), ligesom der stedse kun gives et subjektivt
bestemt O (et O\textsubscript{s}). At denne Relation ikke uden videre kommer
til Bevidsthed, er en Følge af den uvilkaarlige Tillid, med hvilken vi fra
først af møder enhver Sansning, enhver Erindring, enhver Fantasi, enhver
Tanke. Derpaa grunder den Tanke sig, som Grækerne havde saa vanskeligt ved at
komme fra, og som stadig spøger selv i klare Hjerner: hvad jeg erkender, maa
dog i sig selv være Noget; var det Intet, var der heller ingen Erkendelse!
--- Der er jo heller ikke straks Grund til Tvivl om selve de Tankeformer, i
hvilke vi drøfter vore Problemer. Og foreløbig tænker vi ikke paa, at vi selv
er et Passerben, der staar paa et bestemt Sted, ud fra hvilket det andet
Passerben kan anbringes forskellige Steder, men stedse i Relation til det
første Passerben. Vi opdager ikke straks, at vi er et Forholdsled. Den antike
Skepsis opdagede det. Og vi vil fremdrage nogle Eksempler paa, hvorledes man
i nyere Tid mere og mere har faaet Øje for denne Relations Betydning for al
vor Erkendelse. Hvilken særlig Rolle den subjektive Relation (som vi for
Kortheds Skyld kan kalde den) spiller i det enkelte Tilfælde, kan naturligvis
ikke paa Forhaand siges. Det er gennem Analyse af den fremskridende
Erkendelse af Verden, at den Andel, Mennesket (Subjektet) har i denne
Erkendelse, efterhaanden kan blive klar. Det er nogle Eksempler af den nyere
Tids Naturvidenskabs og Erkendelsesteoris Historie, som nu skal fremdrages.

% Sub-subsection mark read off the page image at 600 dpi: GREEK BETA (β),
% upright, roman face, followed by a full stop, at the ordinary paragraph
% indent. UNLIKE the marks of §3b and §3c (pp. 33, 40, 44, 45), which stand
% bare, this one CARRIES A PRINTED TITLE — »Omkring Kopernikanismen.« — set
% in the ordinary roman body face, on the same line, with the footnote
% reference hanging on the title and the text running on after an em dash.
% The Indhold (p. 109) gives the same wording, so head and Indhold agree.
% The run-in em dash after »Kopernikanismen.¹« is dropped here, the head set
% as a heading, and the following paragraph begun \noindent.
\subsubsection*{β.~Omkring Kopernikanismen.\footnote{I dette og det næste
Afsnit bygger jeg paa min Fremstilling af den nyere Filosofis Historie.}}%
\addcontentsline{toc}{subsubsection}{β.~Omkring Kopernikanismen}%
\label{sec:3d-beta}%
\noindent En af de Grunde, paa hvilke Kopernikus byggede sin Hypotese, var
det subjektive Relationsbegreb. Sansning kan, siger han, ikke
% --- p. 63 ---
umiddelbart sige os, hvad det er, der bevæger sig, --- om det er den sansede
Ting, eller det sansende Subjekt, eller maaske begge med forskellig Hastighed
eller i forskellig Retning. Hvis nu Jorden bevægede sig, vilde de sansende
Subjekter ogsaa bevæge sig, og der vilde ikke være nogen Grund til at tillægge
Solen Bevægelse. Allerede et Aarhundrede før Kopernikus var en lignende
Tankegang bleven gjort gældende af den dybsindige Tænker Nicolaus Cusanus,
som gjorde opmærksom paa, at hvor et Menneske saa tænkes at befinde sig, vil
det mene at være i Centrum, ligegyldigt om det befinder sig et Sted paa Jorden
eller paa Solen eller paa andre Stjerner. Man har derfor ikke Ret til at tale
om et Midtpunkt i Verden, altsaa heller ikke Ret til at paastaa, at Jorden er
dette Midtpunkt. Og deraf følger igen, at man ikke har Ret til at paastaa, at
Jorden er i Hvile. Naar vi ikke mærker nogen Bevægelse, kan det jo komme af,
at vi ikke har noget absolut hvilende Punkt, i Forhold til hvilket den kunde
træde frem. Denne Tankegang er hos Cusanus Led af en hel Synsmaade, der
indskærper alle Tankers Forholdsbestemthed (Relativitet), og denne Synsmaade
begrundes igen ved, at al Erkendelse bliver til ved Forbindelse af en Flerhed
af Emner, hvilken Forbindelse nødvendigvis bringer disse Emner i Forhold til
hverandre og til gensidig Bestemmelse ved dette Forhold. Cusanus har i det
hele en mærkelig Evne til at tænke ethvert Emne i dets bestemte Relationer.
Han staar derved som Forløber for Kopernikus, der paaviste, hvilket nyt
Verdensbillede der kunde blive Følgen af denne Synsmaade, og for Giordano
Bruno, der ligesom Cusanus udleder Relationsbegrebet af vor Erkendelses Natur
% "del modo nostro de intendere" is set in the ordinary roman face, not in
% italic; italics.py finds no italic run on p. 63 and the image confirms it.
(del modo nostro de intendere). Vor Tanke arbejder stedse ved at bestemme et
Emne ved et andet, ved at sætte Grænser og
% --- p. 64 ---
\setcounter{footnote}{0}
ophæve Grænser --- i Kraft af et og samme Princip. Saaledes bestemmes det ene
Sted i Forhold til det andet, og Bevægelse i Forhold til et forudsat fast
Punkt. Søger jeg et saadant Punkt paa Solen, vil en Bevægelse tage sig
anderledes ud, end naar jeg søger det paa Jorden. Det gamle Verdensbillede
forudsatte det, der netop skulde bevises, at Jorden var det ene faste Punkt,
i Forhold til hvilket al Bevægelse skulde bestemmes for at være sand. Om
Centrum, Pol, Zenith og Nadir kan vi kun tale absolut, naar vi antager Jorden
for hvilende. Af særlig Interesse er det, at for Bruno følger Tidens
Relativitet med Stedets og Bevægelsens. Aristoteles havde (Physica IV, 14)
lært, at Begreberne Tid og Bevægelse hænger nøje sammen og gensidig maales ved
hinanden, og for ham var Maalestokken given i Himmellegemernes regelmæssige
% p. 64 note 1: only the title »Acrotismus« is italic. »Opera latina« and the
% editor's name »Ed. Fiorentino.« are in the ordinary roman upper-and-lower
% case, NOT small caps — checked at 600 dpi. So the small-caps convention
% round 1 found in the notes of pp. 3--50 does not hold on this page.
Bevægelser. Hertil bemærker Bruno\footnote{\textit{Acrotismus}, Art. 38
(Opera latina. Ed. Fiorentino. I, p. 143---146).}, at da en og samme Bevægelse
maa tage sig forskellig ud fra de forskellige Kloder i Rummet, maa der være
ligesaa mange forskellige Tider, som der er Stjerner! --- Interessant hos
Bruno er idethele den Maade, paa hvilken Relationsbegrebet fører ham til at
kræve nye Iagttagelser og nye Forsøg, hvor man hidtil havde slaaet sig til Ro
ved det umiddelbart Givne som indeholdende sin Forklaring i sig selv.

Hos Galilei indtræder den Vending, at der sondres skarpt mellem den
matematiske og den sanselige Opfattelse. I den matematiske Erkendelse falder
for Galilei Forskellen mellem guddommelig og menneskelig Tanke bort. Det er
hans Platonstudier, der virker, naar han antager, at i Matematiken gælder
Forskellen mellem selve Tilværelsen og vor Opfattelse ikke: her har
menneskelig Erkendelse Del i den Nødvendighed, med hvilken Guddommen
% --- p. 65 ---
tænker de Sandheder, der bærer Tilværelsens Sammenhæng; kun den Forskel
bliver tilbage, at hvad Mennesker møjsommelig og successivt maa arbejde paa
at forstaa, det træder for Guddommen frem i umiddelbar Skuen, og at den
guddommelige Erkendelses Omfang langt overgaar den menneskeliges. Kun for den
matematiske (guddommelige) Opfattelse eksisterer der absolut Rum og absolut
Bevægelse. For vor sanselige Opfattelse er al Bevægelse relativ, da vi efter
Omstyrtelsen af det antike, begrænsede Verdensbillede ikke har noget absolut
Ubevæget, i Forhold til hvilket Bevægelsen kan konstateres. Vor sanselige
Opfattelse er betinget ved vort Standpunkt paa Jorden: \guillemotleft Tænk
Jorden borte, saa eksisterer der ingen Solopgang eller Solnedgang, ingen
Meridian, ingen Dag og ingen Nat!\guillemotright{} Og her paa Jorden er det
igen vor legemlige Organisme, der betinger de Egenskaber, vi --- bortset fra
de rent matematiske (primi e reali accidenti) --- tillægger Tingene, altsaa
alle de egentlige Sansekvaliteter (Farve, Lugt, Smag o. s. v.). Tag Legemet
bort, saa falder ogsaa alle Sansekvaliteter bort, og saa vil vor Erkendelse
kunne gaa helt i et med den guddommelige Erkendelse.

Den skarpe Modsætning mellem matematisk (guddommelig) Erkendelse, for hvilken
der eksisterer absolut Rum og absolut Bevægelse, og den sanselige
(menneskelige) Erkendelse er det saa meget des mærkeligere at finde hos
Galilei, som hans Stordaad netop var den metodiske Samvirken af Deduktion og
Induktion, af Matematik og Eksperiment. Men denne Modsætning fik stor
Indflydelse i de følgende Aarhundreder, især da den støttedes ved Newton's
store Autoritet.

Paa Grund af Sætningen om Bevægelsens Relativitet kunde rent formelt
% The foot of p. 65 carries the printer's signature line »Vidensk. Selsk.
% Filosof. Medd. I, 3.« with the signature numeral 5; omitted as printer's
% apparatus. Note the comma in »I, 3.« — the sheets signed on pp. 33 and 45
% were recorded with a full stop, »I. 3.«
(matematisk) Himmelfænomenerne lige%
% --- p. 66 ---
saa godt forklares ved, at Solen gik rundt om Jorden, som ved, at Jorden gik
rundt om Solen. Og Tycho Brahe havde opstillet en Mellemform mellem de to
Hypoteser, ifølge hvilken Jorden stod fast, medens Solen og Planeterne gik
rundt om den. Gassendi anstillede en Sammenligning mellem de tre saaledes
opstillede Hypoteser og søgte at vise, at man rent videnskabelig ikke kunde
vælge mellem dem. Men saa paaviste Newton i sit berømte Værk, at de Love,
Kepler havde opdaget for Planetbevægelse, kunde udledes, hvis man udstrakte
Tyngdeloven til at gælde for Himmellegemerne, og da maatte det nødvendigvis
være Jorden og Planeterne, der bevægede sig om Solen. I Kraft af
Faldbevægelsens Lov blev altsaa tilsidst kun én Hypotese mulig. Paa analog
Maade førte omtrent samtidig Leibniz ud over den Skepsis, Bevægelsens
Relativitet tilsyneladende motiverede, idet han hævdede, at det afgørende for
Erkendelsen ikke var selve Bevægelsen, men Bevægelsens Love, og her opstillede
han saa --- under Kritik af Descartes' Lære om Bevægelsens Bestaaen ---
Sætningen om Kraftens Bestaaen. Her var da en Realitet, der ikke rettede sig
efter Iagttagerens Stade i Rummet. De to store Forskere førte, hver ad sin
Vej, ud over den tilsyneladende Subjektivisme og tog det afgørende Stade i
Lovene for Bevægelsen.

Alligevel hævdede Newton (her vistnok under Indflydelse af Platonikeren Henry
More), ligesom Galilei en skarp Modsætning mellem det sanselige Rum, som den
populære Opfattelse holder sig til, men som ikke gør nogen absolut
Stedsbestemmelse mulig og kun frembyder relativ Bevægelse, og det absolute
% p. 66: »Stedbestemmelse« here, without the s of »Stedsbestemmelse« three
% lines above and of p. 67; printed so, transcribed so.
Rum, som ikke staar i nogensomhelst ydre Relation, og som tilsteder absolut
Stedbestemmelse og absolut Bevægelse. Det er ved videnskabelig Orientering i
Verden (in rebus philosophicis), at vi benytter
% --- p. 67 ---
os af de absolute Steder (loca primaria), og det er da det matematiske Rum,
der ligger til Grund. I det praktiske Liv (in rebus humanis), derimod lader vi
os nøje med det sanselige Rum, der ikke gør absolut Bestemmelse mulig. Denne
Opfattelse, der hænger sammen med Newtons ejendommelige Metafysik, blev
foreløbig den herskende indenfor Naturvidenskaben. Det var en forsvarlig Bom,
der ved den store Naturforskers Lære paa dette Punkt blev slaaet for
Relationsbegrebets fortsatte Anvendelse. Da man ikke længere, som endnu
Kopernikus, antog Fiksstjernesfæren (\guillemotleft den ottende
Sfære\guillemotright) for at staa fast, saa at den --- som communis
universorum locus (for at bruge Kopernikus' Udtryk) --- kunde tjene til
% "point de repère" and "communis universorum locus" are both roman here,
% not italic; italics.py finds no italic run on p. 67.
Grundlag (point de repère) for absolut Stedsbestemmelse, var det egentlig en
rent mystisk Henvisning at appellere til det absolute Rum, der ikke var
tilgængeligt for nogen menneskelig Erfaring. --- Paa dette Punkt af Drøftelsen
griber nu to filosofiske Tænkere ind paa ejendommelig Maade.

% Sub-subsection mark read off the page image at 600 dpi: GREEK GAMMA (γ),
% upright, roman face, followed by a full stop, at the ordinary paragraph
% indent. Like β on p. 62, and unlike the bare marks of §3b and §3c, it
% CARRIES A PRINTED TITLE — »Filosofiske Synspunkter.« — run in on the same
% line in the roman body face, the text continuing after an em dash. The
% Indhold (p. 109) gives the same wording. The printing sets rather more
% than the normal leading above this head. The run-in dash is dropped here.
\subsubsection*{γ.~Filosofiske Synspunkter.}%
\addcontentsline{toc}{subsubsection}{γ.~Filosofiske Synspunkter}%
\label{sec:3d-gamma}%
\noindent Det var en ny Platonisme, der sejrede med Galilei og Newton, i
ejendommelig Modsætning til den Maade, paa hvilken disse store Forskere i
deres Naturvidenskab anvendte Vekselvirkningen mellem Tænkning og Iagttagelse.
Men der rejser sig nu en Opposition, der kan sammenlignes med den, der i
Oldtiden rejstes af Karneades og Ainesidemos mod den platoniske Idelære. Det
var Berkeley og Kant, der rejste denne Opposition. De kritiserer begge, hver
paa sin Vis, det absolute Rum og hævder, at hvis man lader det erkendende
Subjekt ude af Betragtning, falder med det samme alle Sandheder og alle Gaader
bort.

% The foot of p. 67 carries the signature mark »5*«; omitted as printer's
% apparatus.
Berkeley's Hovedtanke var, at vi tænker og maa tænke
% --- p. 68 ---
\setcounter{footnote}{0}
i Eksempler. Almindelige Begreber, Sætninger og Love kender vi kun fra de
specielle Tilfælde, i hvilke de lægger sig for Dagen, og det er uberettiget at
tillægge dem og de Forhold, de udtrykker, Eksistens som en Verden for sig bag
de for Iagttagelsen fremtrædende Kvaliteter og Kvalitetsforskelligheders
Verden. De \guillemotleft primære og reale\guillemotright{} Egenskaber, som
Galilei tilskrev absolut Eksistens, opfatter vi dog stedse kun under bestemte
Betingelser: i en eller anden Afstand, makroskopisk eller mikroskopisk
o. s. v. Udstrækning i Rummet maa dele Skæbne med alle andre Egenskaber.
Enhver Bevægelse maa have en bestemt Retning og en bestemt Hastighed. Der
gives hverken Udstrækning eller Bevægelse i al Almindelighed. Naar
Spørgsmaalet opkastes om det rigtige Sted, den rigtige Størrelse, Retning
eller Hastighed, da kan et saadant Spørgsmaal kun afgøres under Hensyntagen
til det opfattende og tænkende Subjekt. Hvis jeg, siger Berkeley, som den
barnlige Troende, han var, havde været tilstede ved Skabelsen, vilde jeg have
set Tingene blive til i den Orden, Bibelen fortæller. Men i daglig Tale ser vi
bort fra Subjektet, skønt det stedse maa underforstaas. --- Her fremtræder
klart Modsætningen til Newton, der mente, at vi i streng Videnskab ser bort
fra Subjektet. Berkeley hævder, at der i Videnskaben (in rebus philosophicis,
for at bruge Newton's Udtryk) ligesaavel maa underforstaas et Subjekt som i
det praktiske Liv (in rebus humanis).

% p. 68 note 1 runs over onto p. 69, where it is continued above that page's
% own note 1; set here as a single \footnote at the reference. »Thomas
% Aqvinas« is in ordinary roman upper-and-lower case, NOT small caps —
% checked at 600 dpi.
Kant har, som allerede omtalt, set, at Relation er et for al Erkendelse
gældende Grundbegreb, kun at han, med afgjort Inkonsekvens, forbeholdt sig en
relationsløs \guillemotleft Ding an sich\guillemotright{} til Forklaring af,
at der eksisterer for os noget, som vi ikke selv har
frembragt.\footnote{Naar \guillemotleft Ding an sich\guillemotright{} skal
indeholde Grunden til Stof (og Form) i vor Erkendelse, staar den alligevel i
Relation til Erkendelsen, uagtet den skulde være relationsløs. Eller vor
Erfaring forholder sig til den, men den ikke til vor Erfaring. Det er samme
ensidige Relation, som Thomas Aqvinas prøvede at gennemføre i Teologien. Se
ovenfor p. 47 Note.} Blandt Relationerne er
% --- p. 69 ---
\setcounter{footnote}{0}
det for ham, ligesom for Berkeley, Relationen Subjekt---Objekt, der optager
Interessen, og han hævder Subjektets fundamentale Stilling i denne Relation.
Selv i den rene Matematik og den eksakte Erfaringsvidenskab forudsættes et
Erkendelsessubjekt. Kun for et saadant Subjekt eksisterer der Erfaring i
Ordets strenge Betydning, en Erfaring, i hvilken matematiske Begreber faar
deres Anvendelse og har deres reale Betydning.

Kant begyndte som ivrig Newtonianer og hævdede foreløbig, ligesom sin Mester,
det absolute Rum som objektiv Forudsætning for streng Erfaringsvidenskab
saavel som for den rene Geometri (der for ham havde naaet sit definitive
Grundlag med Euklid). Endnu i \guillemotleft Kritik der reinen
Vernunft\guillemotright{} hævder han et absolut Rum, kun at det newtonske Rum
saa at sige er \guillemotleft slaaet ind\guillemotright{} og blevet til en
menneskelig Anskuelsesform i Stedet for en guddommelig Anskuelsesform, et
sensorium hominis i Stedet for, som for Newton, et sensorium
% p. 69 note 1: the italic run is the title of Høffding's own earlier paper;
% \textit, not \emph. »Smlgn.« — the ABBYY layer reads it »Smign.«
dei.\footnote{Smlgn. \textit{Om Kontinuiteten i Kants filosofiske Udvikling},
p. 31.} Af særlig Interesse er det dog, at han fremhæver, at Begrebet Rum (og
det Samme gælder Begrebet Tid) ikke blot er et Almenbegreb, men hvad man kunde
kalde et typisk Individualbegreb: ethvert bestemt Rum er en Del af Rummet,
ikke blot et Eksempel paa Begrebet Rum. Og tillige maa mærkes, at naar han
kalder Rummet en Anskuelsesform, mener han ved Anskuelse ikke et én Gang for
alle færdigt Resultat, men en Funktion, en Anskuen, der fører til at danne
Anskuelsesbilleder eller konstruere Skemaer. Overhovedet mener Kant ved
\guillemotleft Former\guillemotright{} Maader, paa hvilke den forbindende
Funktion, den Syn%
% --- p. 70 ---
tese, hvori al Bevidsthed og al Erkendelse bestaar, foregaar i de enkelte
Tilfælde. Vi konstruerer Rummet ved at føje Punkt til Punkt, Stykke til Stykke
efter en bestemt Regel. Det er i Identiteten af en saadan Regel, at
Bevidsthedens Enhed ytrer sig i det enkelte Tilfælde.

Ved nærmere Gennemførelse af denne Tankegang kom Kant senere til at nægte det
absolute Rums subjektive Realitet, ligesom han allerede i \guillemotleft Kritik
der reinen Vernunft\guillemotright{} havde nægtet dets Realitet som i sig
eksisterende. Allerede i første Udgave af sit Hovedværk havde han udtalt, at
bortset fra alle Ting, der udfylder og begrænser, er det absolute Rum ikke
andet end \guillemotleft die blosse Möglichkeit äusserer
Erscheinungen\guillemotright. Nogle Aar efter (i \guillemotleft Metaphysische
Anfangsgründe der Naturwissenschaft\guillemotright) gaar han et Skridt videre,
støttende sig til ethvert Steds og enhver Bevægelses Relativitet. I det
absolute Rum skulde der jo kunne ses bort fra alle Relationer. Men naar bliver
vi færdige med at sætte alle Punkter i Forhold til hverandre, med at drage
alle Linier o. s. v.? Hvor er det Punkt eller den Grænse, der selv
relationsløs bestemmer alle Punkter og Grænser? Kant erklærer derfor, at det
rene, absolute Rum \guillemotleft für alle mögliche Erfahrung nichts
% p. 70: the printing has »äusserer Erscheinungen« with the umlaut but
% »ausser dem gegebenen« without it — both checked at 600 dpi, and both are
% correct German, so neither is a compositor's slip. The ABBYY layer reads
% the second one »äusser«; that is the text layer's error, not the print's.
ist\guillemotright: \guillemotleft Der absolute Raum ist an sich nichts und
gar kein Objekt, sondern bedeutet nur einen jeden relativen Raum, den ich mir
ausser dem gegebenen jederzeit denken kann.\guillemotright{} (Anf. Skr. 1786,
p. 16 smlgn. p. 3). Kant hylder her faktisk Berkeley's Lære, at vi tænker i
Eksempler. Det absolute Rum betyder blot det Negative, at intet givet Rum er
afsluttende. Der kan stedse forestilles nye Rum. Det er selve Relationens,
Forholdssættelsens Tanke, der fører Kant bort fra hans absolute
Anskuelsesform. Rummet bliver en \guillemotleft Ide\guillemotright{} i Kants
Betydning, et Grænse%
% --- p. 71 ---
begreb, der afviser Afslutning og hævder Muligheden af fortsat Rækkedannelse.
Rummets Begreb faar kun regulativ Betydning, ligesom de af Kant selv opstillede
\guillemotleft Ideer\guillemotright, og i Grunden hører det ind under
\guillemotleft den kosmologiske Ide\guillemotright{} (Begrebet Verden).

Kant har selv sammenlignet sin Filosofi med Kopernikanisme. Og paa en Maade
kunde han og Berkeley være omtalte i det Afsnit, som handlede om de Tanker,
der udviklede sig \guillemotleft omkring Kopernikanismen\guillemotright. Men
de to Filosofer var dog mest optagne af at bekæmpe den Dogmatisme, der, med
formentlig Støtte i matematisk Videnskab, havde holdt sig trods
Kopernikanismens Sejr. Og hvad særlig Kant angaar, saa var det hans
Hovedbestræbelse paa én Gang at hævde den store Betydning af det Newton'ske
Videnskabsbegreb og at føre de absolute Former, Newton endnu havde ladet staa,
tilbage til at være Former for den arbejdende Tanke.

Skønt Kant stadig satte Tiden parallelt med Rummet, drager han dog for dens
Vedkommende ikke en lignende Konsekvens som med Hensyn til Rummet. Og dog
vilde det --- særlig ved hans mærkelige Lære om \guillemotleft
Skematismerne\guillemotright, der fremhæver Aktiviteten i Tidsanskuelsen ---
have ligget ham nært at se, at vi ligesaa lidt bliver færdige med Tiden som
med Rummet, og at den absolute Tid kun kan betyde, at nye Tidsrelationer
stedse er mulige. ---

Jeg behøver ikke her at gaa ind paa en Kritik af Kant. En saadan vilde særlig
gaa ud paa at vise, at selve Begrebet Erfaring i Kant's ovenfor angivne
strenge Betydning er en \guillemotleft Ide\guillemotright, ligesom Rummet og
Tiden, og at Kant kun ved en forunderlig Inkonsekvens kan antage en
\guillemotleft Ding an sich\guillemotright{} uden Relationer.


% ============================================================
%  BATCH 7 — §3d, Den nye Relativitetsteori (head at p. 72)
%  COLLATION BATCH. A transcription of these five pages already existed in
%  this file before scan.pdf was obtained; it is preserved at
%  .parts/legacy-pp72-76.texfrag and is to be collated, not re-derived.
% ============================================================
% --- p. 72 ---
% Sub-subsection mark read off the page image at 600 dpi: GREEK DELTA (δ),
% unambiguous — the upright loop-and-hook form, not a Latin "b" (no straight
% ascender, and the bowl closes at the top left). The text layer reports it as
% a Latin "b."; the earlier skeleton, built from the Indhold's OCR, left it
% open. Confirmed also by round 1's finding that the Indhold's OCR renders β
% as "ß" and γ as "Y-", so a plain "b." cannot be a β.
%
% Unlike the bare Greek marks at pp. 33, 45, this one DOES carry a printed
% title, and the head is RUN-IN: the printed page reads
%     δ. Den nye Relativitetsteori. — Trods den store Udvi-
% at the ordinary paragraph indent, in the same roman as the body (not
% letterspaced — spacing.py finds no spærret run on p. 72). The run-in em dash
% is dropped here and the head is set as a heading; the following paragraph
% takes \noindent. The head agrees with the Indhold word for word.
\subsubsection*{δ.~Den nye Relativitetsteori.}%
\addcontentsline{toc}{subsubsection}{δ.~Den nye Relativitetsteori}%
\label{sec:3d-delta}%
\noindent
Trods den store Udvidelse af Horizonten, som Kopernikanismen
gav Stødet til, havde Naturvidenskaben fastholdt Begreberne absolut Sted
og absolut Bevægelse, der dog kun havde deres Berettigelse indenfor det
gamle Verdensbillede med dets begrænsede og faste Ramme. Galilei's og
Newton's absolute Rum afløste det ved den ottende Sfære begrænsede Rum.
Men der kom den Tid, da Relationsbegrebet ogsaa her gjorde sin Konsekvens
gældende, og det ikke blot hvad Forholdet mellem Emnerne indbyrdes angik,
men ogsaa med Hensyn til Forholdet mellem Emnerne og det erkendende
Subjekt.

Det var Bevægelseslærens, og dermed Naturvidenskabens første Sætning,
der først blev drøftet fra Relationsbegrebets Synspunkt. \textsc{Neumann}
(Ueber die Prinzipien der Galilei-Newton'schen Theorie, 1870) gjorde
% The title is set in ORDINARY ROMAN on p. 72 — not italic, not small caps,
% not letterspaced. italics.py finds no italic run anywhere on p. 72 and the
% 600 dpi image confirms it. Only the name NEUMANN is distinguished (small
% caps, first mention). Same for MACH and the Mach title below.
opmærksom paa, at Inertisætningen, ifølge hvilken et Legemes Hastighed og
Retning kun ændres ved Paavirkning udefra, nødvendigvis i hvert enkelt
Tilfælde forudsætter Angivelse af det bestemte Sted, i Forhold til hvilket
Hastigheden og Retningen bevares, naar der ingen ydre Paavirkning kommer.
Neumann fandt nu intet Sted, der kunde gøre Fordring paa at tages til
Udgangspunkt for et absolut Koordinatsystem, men mente, at man hypotetisk
maatte antage et Centrum, for at Begrebet absolut Bevægelse kunde bevare
sin Mening. \textsc{Mach}, der omtrent samtidig ad sin egen Vej var
kommen til at indse, at Inertisætningen hvilede paa Forudsætningen om et
bestemt Koordinatsystem, fandt dog ikke Neumann's Udvej forløsende. Han
fordrede Relationsbegrebet gennemført. Den dogmatiske Opfattelse af
Inertisætningen falder efter hans Mening først bort, naar man bliver klar
over, at i og for sig kan ethvert Legeme i Verden bruges som Udgangspunkt,
og at vi i Praksis i hvert
% --- p. 73 ---
\setcounter{footnote}{0}%
enkelt Tilfælde benytter et eller andet Legeme som forudsat fast
Udgangspunkt for de Koordinater, i Forhold til hvilke en Bevægelses
Hastighed og Retning bestemmes. (Die Erhaltung der Arbeit, 1872). Denne
% As printed: the full stop stands BEFORE the parenthesis, and the title is
% in ordinary roman, not italic (italics.py: no italic on p. 73 except the two
% runs in the footnotes; confirmed on the image).
Vej er man faktisk altid gaaet, kun at man indenfor
det gamle Verdensbillede, der netop hvilede paa Forudsætningen om et
absolut Midtpunkt og en absolut fast Grænsesfære for Verden, intet Valg
behøvede. Mach lægger overhovedet stor Vægt paa de historiske
Forudsætninger, som hver enkelt Tids Forskning beror paa, men viser saa,
hvorledes disse Forudsætninger varierer fra Tid til Tid, og at ingen af
dem kan gøre Fordring paa absolut Nødvendighed og Forstaaelighed. Det
kommer her stedse an paa, hvad franske Fysikere kalder
points de repère.\footnote{Smlgn.\ f.\ Eks.\ \textsc{Painlevé}
  (\textit{Mécanique}, p.~399) [i Sammelværket \guillemotleft De la Méthode
  dans les Sciences\guillemotright , 1909]: \guillemotleft\ldots\ les
  mouvements absolus, c'est-à-dire au fond les mouvements convenablement
  repérés.\guillemotright}
% »points de repère« is set in ORDINARY ROMAN on p. 73, not italic — so no
% \textit here, despite the phrase being one the edition italicises elsewhere.
% In the note, »Sammelværket« is as printed (not »Samleværket«/»Samlerværket«);
% the Painlevé volume title is in guillemets, not italic; only »Mécanique« is
% italic.

En endnu mere radikal Anvendelse af Relationsbegrebet er fremkommen i
Begyndelsen af det tyvende Aarhundrede. Hvad der i denne Sammenhæng har
Interesse, er, at det har vist sig nødvendigt at tage særligt Hensyn til
det opfattende Subjekt og dets Stade. Kopernikus havde allerede indskærpet
dette, hvad Opfattelsen af Himmellegemernes Bevægelse angaar, og Neumann
og Mach, hvad Bevægelsens Hastighed og Retning angaar.
\textsc{Einstein}\footnote{\textit{Über die spezielle und die allgemeine
  Relativitätstheorie,} (1916)}
% As printed: a comma stands inside the italic before the parenthesis, and
% the note ends WITHOUT a full stop. Both as set; not repaired.
har nu søgt at vise, at hvad der gælder Rum og Bevægelse ogsaa gælder
Tiden. Naar man ikke har opdaget dette før, er det ifølge Einstein at
udlede af, at man kun tog Hensyn til Hastigheder, der var meget smaa i
Forhold til Lysets Hastighed. Den sidste Tidsmaalestok bliver for den
moderne Fysik Lysets Hastighed, medens den for Aristoteles, og endnu for
Kopernikanerne, hentedes fra Himmellegemernes
% --- p. 74 ---
\setcounter{footnote}{0}%
Bevægelser. (Bruno's Lære
om Tidsbegrebets Relativitet hviler paa denne Maalestok. Se ovenfor). Paa
Grundlag af den elektromagnetiske Teori om de fysiske Fænomener, for
hvilken den mekaniske Naturopfattelse kun staar som et specielt Tilfælde,
søger Einstein at vise, at der kan paavises Forhold, under hvilke to
Iagttagere, af hvilke den ene er knyttet til faste, den anden til
bevægelige Akser, men begge i Hvile i Forhold til deres Akser, ikke vil
faa samme Tidsopfattelse: hvad der for den ene Iagttager staar som et
Samtidighedsforhold mellem to Begivenheder, vil for den anden staa som et
Successionsforhold. En Tidsangivelse faar i det hele kun Mening, naar der
angives, i Forhold til hvilket Koordinationssystem det gælder, og særlig,
om dette System selv forholdsvis er i Hvile eller Bevægelse.

Hvis man vilde mene, at der her hos den ene eller den anden Iagttager maa
% spacing.py flags »vilde mene,« and »anden Iagttager« here as letterspacing
% candidates. The image shows ordinary justification stretch, not spærret;
% no emphasis is set.
være en Illusion tilstede, svares der fra det Einsteinske Standpunkt, at
enhver Iagttager her har sit Maalsystem, der er ligesaa berettiget som en
andens, men at det naturligvis meget vel er muligt for den ene Iagttager
at prøve paa at sætte sig paa den andens Standpunkt og ræsonnere ud fra
det.\footnote{\textsc{Langevin} i \guillemotleft Bulletin de la Société
  française de Philosophie\guillemotright\ 11.\ Oktbr.\ 1911. --- Der maa vel
  forudsættes, at den Iagttager, der kan sætte sig paa en anden Iagttagers
  Standpunkt, kan tilstrækkelig Fysik.}

De Kendsgerninger og Slutninger, paa hvilke den Einsteinske Teori bygger,
er endnu Genstand for Drøftelse, og Anskuelserne staar skarpt imod
hinanden. I en saadan Strid kan Filosofen som saadan ikke gribe ind, og
Filosofien har da ogsaa Tid til at vente. Om Fysiken virkelig kan paavise
Situationer, hvor det, der staar som Samtidighed for en Iagttager, staar
for en anden Iagttager som suc%
% --- p. 75 ---
\setcounter{footnote}{0}%
cessivt, maa Fysikere afgøre mellem
hverandre indbyrdes. Hvis det maa besvares med Ja, staar vi overfor det
betydningsfulde Resultat, at Naturvidenskaben nu igen, ligesom ved den
kopernikanske Hypotese, ad sin egen Vej er kommen til at se Nødvendigheden
af ikke blot at holde sig til Iagttagelsernes Indhold, men at tage Hensyn
til Iagttagerens Forudsætninger. I en interessant Diskussion i
\guillemotleft Aristotelian Society\guillemotright\ (1919) indskærpede
\textsc{Whitehead} Betydningen af, hvad han træffende kaldte
\guillemotleft the percipient event\guillemotright\footnote{Se ogsaa hans
  \textit{Principles of Natural Knowledge} (1919), p.~68. --- Whitehead's
  Opfattelse er dog forskellig fra Einstein's derved, at \guillemotleft the
  percipient event\guillemotright\ for ham kun er en Faktor ved Siden af mange
  andre, og ikke, som hos Einstein, af væsentlig Belydning for al Maaling af
  Tid og Rum. Smlgn.\ herom \textsc{Haldane}: \textit{The Reign of Relativity}
  (1921), p.~69\,f.; 107\,ff.},
% PRINTER'S ERROR, transcribed as printed: »Belydning« for »Betydning«
% (600 dpi: a clear l, no crossbar; the text layer silently normalises it).
% Note the reference numeral: it stands BEFORE the comma on p. 75, at all
% three sites in this paragraph — »event«¹,  Betydning²,  Ret³.
og \textsc{Nicholson} erklærede, at til de Data, som er af
afgørende Betydning (ultimate character), maa ogsaa regnes
% »ultimate character« is in ORDINARY ROMAN inside the parentheses, not
% italic: italics.py reports only three italic runs on p. 75, all of them
% work titles in the footnotes, and the image agrees.
\guillemotleft the essential framework of our modes of
perception.\guillemotright

Man har overfor den Einsteinske Hypotese, ligesom overfor tidligere
Paavisninger af Relationsbegrebets erkendelsesteoretiske
Betydning\footnote{Smlgn.\ \textit{Den menneskelige Tanke}, p.~304.},
gjort gældende, at der dog kun kan være én Sandhed i hvert enkelt
Tilfælde, og at det vilde stride mod Modsigelsens Grundsætning, hvis to
Iagttagere, af hvilke den ene erklærer for samtidigt, hvad den anden
erklærer for successivt, begge skulde have Ret\footnote{\textsc{Kroman} i
  Tidsskrift for Fysik, XIV, p.~15.}.
% The journal title is in ordinary roman in the print, not italic.
Men der er ingen Grund til at
frygte for Sandhedens Enhed, selv om Einstein skulde faa Ret. Den
egentlige Sandhed ligger nemlig da i den Nødvendighed, der fysisk paavises
af, at forskellige Iagttagere under visse bestemte Betingelser faar
forskellig Tidsopfattelse. Hver enkelt Iagttager opfatter Tiden, som han
under sine Betingelser maa. Men hvis en Iagttager tillige er Fysiker, vil
han (Rigtigheden af
% --- p. 76 ---
Einsteins Teori stadig forudsat) kunne forklare os,
hvorfor en anden Iagttager faar en anden Opfattelse end han selv. (Smlgn.\
Newton's og Leibniz' Fremhæven af Bevægelsens Lovmæssighed som Udtryk for
dens Objektivitet. Ovenfor p.~66). Naturligvis beror den fysiske Forklaring
af de forskellige Iagttageres forskellige Tidsopfattelse selv paa visse
Forudsætninger, forudsætter et Koordinationssystem, ud fra hvilket de
forskellige Iagttageres Koordinationssystemer bestemmes, og her er da
Mulighed for en ny Diskussion. Hvis denne Diskussion afsluttes, kommer det
af, at der ikke kan paavises Standpunkter, som er endnu mere fundamentale.
Kant har forlængst hævdet, at al Nødvendighed er hypotetisk. Men derfor
er den alligevel Nødvendighed, eller rettere: netop derfor er den
Nødvendighed. En Nødvendighed uden Forudsætninger kender vi ikke.

Hvis den nye Relativitetsteori vinder fortsat Bekræftelse, er det ikke
blot et i og for sig interessant Resultat af fysisk Forsken, der kommer
til at foreligge. Denne Teori indeholder desuden, som Einstein selv har
udtalt, en Spore til fortsat Forsken. Hvis det er saa, at der er Mulighed
for forskellige Relationssystemer ved Opfattelsen af et Emne, saa ligger
netop heri en Opfordring til at finde alle de bestemte Relationer, der
ligger til Grund i hvert enkelt Tilfælde, og Forskningen leder saaledes
stedse videre i Stedet for at slutte af paa dogmatisk Vis, som om de
Forudsætninger, man hidtil havde arbejdet med, var de eneste mulige. Der
knytter sig, erkendelsesteoretisk set, to Interesser til Relationsbegrebet.
Den ene er, at der kastes Lys over vor Erkendelses Natur og Betingelser;
den anden er, at der anvises nye Opgaver, viser sig et plus ultra,
% »plus ultra« is in ORDINARY ROMAN on p. 76, not italic (italics.py finds no
% italic anywhere on p. 76; confirmed on the image).
idet enhver Tanke peger ud over sig selv. ---

Campanella sagde i sin Apologi for Galilei, at hvis
% ============================================================
%  BATCH 8 — §4 Relationsbegrebet i Etiken (head at p. 77), §4a–§4c
% ============================================================
% --- p. 77 ---
\setcounter{footnote}{0}
Galilei havde Ret, maatte man filosofere paa en ny Maade.
Men hvis Einstein har Ret, behøver man derfor ikke at
filosofere paa en ny Maade.\footnote{\textsc{Einstein} selv er enig heri. I
sin Fortale til \textsc{Lucien Fabre}'s Bog \guillemotleft Les théories
d'Einstein\guillemotright{} (1921) siger han: \guillemotleft
Relativitetsteorien hverken kan eller vil give noget Verdenssystem, men
kun angive en Begrænsning, som Naturlovene
underligger\guillemotright. At der ikke blot er Tale om en Grænse, har
vi ovenfor set.}
% p. 77 note 1: EINSTEIN and LUCIEN FABRE are set in SMALL CAPS in the
% printing (checked at 400 dpi); »Les théories d'Einstein« is roman inside
% guillemets, not italic. spacing.py flagged »noget« and »ovenfor« as
% letterspaced singles; on the image both are ordinary roman widened only
% by justification. No emphasis there.
Allerede rent filosofisk er man, som den foregaaende Fremstilling har
vist, forlængst bleven opmærksom paa Relationsbegrebets Betydning for
Erkendelse og Verdensanskuelse. Men hvis Einstein har Ret, foreligger
der nu en ny, uventet, men betydningsfuld Erfaring i den Henseende. ---

Gennem Nødvendigheden af at tage Hensyn til Erkendelsessubjektet og
gennem Muligheden af en Modsætning mellem forskellige
Erkendelsessubjekters hver for sig nødvendige Opfattelser, frembyder den
nyere Naturvidenskab --- hos Einstein som hos Kopernikus --- Analogi
til, hvad der gør sig stærkt gældende paa Aandsvidenskabens Omraade,
saaledes som jeg nu skal forsøge at vise.

\section{Relationsbegrebet i Etiken.}\label{sec:4}
% p. 77: printed head reads "4. Relationsbegrebet i Etiken." --- bold,
% centred, on its own line, the numeral supplied by the counter. Matches
% the Indhold (p. 109).

\subsection{Historieforskningen og Relationsbegrebet.}\label{sec:4a}
% p. 77: printed head reads "a. Historieforskningen og Relationsbegrebet."
% --- letterspaced (spærret), and RUN IN, unlike §3's lettered heads, which
% stand on their own line: here the head sits at the ordinary paragraph
% indent and the text follows on the same line after an em dash,
% "a. Historieforskningen og Relationsbegrebet. --- Aristoteles lærte...".
% The run-in dash is dropped and the following paragraph set \noindent.
% Matches the Indhold word for word.
\noindent
Aristoteles lærte, som allerede omtalt, at kun det Almene, ikke det
Individuelle som saadant kunde være Genstand for Videnskab. Det hænger
sammen med hans Filosofis logisk-skematiske Karakter. Videnskab var for
ham, som for Platon, Underordning af det Enkelte og Individuelle under
Almenbegreber. Men nu var tillige for Aristoteles alt Virkeligt enkelt
og individuelt; det er paa dette Punkt, han træder i Modsætning til
Platon. Den Vanskelighed, som herved opstaar, og som paa en glimrende
Maade er belyst
% --- p. 78 ---
\setcounter{footnote}{0}
af \textsc{Zeller} (Die Philosophie der Griechen, II, 2), kan kun
overvindes gennem en nærmere Drøftelse af, hvad man forstaar, og
videnskabelig maa forstaa ved Virkelighed, og en saadan Drøftelse blev,
som vi har set, indledet med stor Energi af Oldtidens Akademikere og
Skeptikere.\footnote{En anden Sag er det, hvorledes Vanskeligheden kunde
klares indenfor selve den aristoteliske Filosofi. \textsc{Hans v. Arnim}
(\textit{Die europäische Philosophie des Altertums}. Kultur der Gegenw.
I, 6, p.~173 f.) finder Forklaringen i, at baade Arter og Individer hos
Aristoteles betragtes som \guillemotleft Væsener\guillemotright. Men
derved fremkaldes blot et nyt Spørgsmaal, nemlig, hvorledes og hvorfor
den ene Art \guillemotleft Væsener\guillemotright{} kan være Genstand for
Videnskab, den anden ikke.}
% p. 78: ZELLER in the body and HANS V. ARNIM in the note are SMALL CAPS
% (checked at 400 dpi; the "v." is likewise small-capped). Zeller's title
% "Die Philosophie der Griechen" is set ROMAN in the printing, inside
% parentheses and without quotation marks --- left roman; only Arnim's
% "Die europäische Philosophie des Altertums" is italic. The volume
% reference reads "Kultur der Gegenw. I, 6" with a roman capital I, not
% the arabic 1 the text layer gives. spacing.py's RUN on "overvindes
% gennem nærmere" is justification, not spærret.
Allerede her kom den Tanke frem, at Virkelighedskriteriet kun kan bestaa
i den faste Sammenhæng mellem Iagttagelserne. Ligesom de formale
Videnskaber udspringer af en Stræben efter at bevare Overensstemmelsen
med sig selv, saaledes udspringer de reale Videnskaber af en Stræben
efter at finde Overensstemmelse mellem de for Iagttagelse givne Emner.
Og det heri liggende Virkelighedskriterium maa ligesaa vel gælde for det
Individuelle, det, der maaske kun forekommer en eneste Gang, som for
det, der kan gentage sig atter og atter, fordi det gælder for hele
Rækker af individuelle Emner.

Et Emne erkendes som realt, dels gennem dets Sammenhæng med andre
Emner, overfor hvilke ingen Tvivl er vakt, dels gennem den Sammenhæng,
der kan findes mellem de forskellige Former, under hvilke selve Emnet
til forskellige Tider fremtræder, og mellem de forskellige Egenskaber,
med hvilke det til en given Tid optræder.

Historieforskning gaar nu netop ud paa begge Arter af Sammenhæng. De
Begivenheder, Personligheder eller Institutioner, der er dens Emner,
undersøger den i Forhold til de Omstændigheder, under hvilke de
foreligger, og hvert
% --- p. 79 ---
\setcounter{footnote}{0}
enkelt af hine Emner undersøger den med Hensyn til dets indre
Sammenhæng. I begge Henseender begyndes der med Kritik af Beretninger,
men selve denne Kritik lægger det nævnte Virkelighedskriterium til
Grund. Derefter bliver det Opgaven at finde det Baand, der forbinder
Emnet med andre Emner, og det Baand, der sammenknytter Emnets
forskellige Sider og Stadier.

Videnskabens Genstand er hverken det Almene eller det Individuelle i og
for sig alene, men Virkeligheden, der kun kan kendes gennem de Baand,
der forener Emnerne indbyrdes saa vel som Elementerne indenfor det
enkelte Emne.\footnote{Smlgn. \textit{Totalitet som Kategori}, p.~59---66.}
Hvad der ikke kan belyses gennem Anvendelse af Relationsbegrebet, kan
ikke forstaas.

Historiefilosofer har i den nyeste Tid fremdraget tre Begreber som
særlig karakteristiske for historisk Forsken\footnote{Anf. Skr.
p.~71---72.}, og disse tre Begreber udtrykker netop Relationer mellem et
Emne, dets Betingelser eller dets Virkninger.
% p. 79: both notes carry plain arabic numerals; the text layer's "*" for
% note 2 is an artefact of the superscript 2 (checked at 400 dpi). Note 1
% opens "Smlgn.", not the layer's "Smign.". The two self-citations are
% titles of Høffding's own works -> \textit{}.

Historisk Centrum er den bestemte Begivenhed, Personlighed eller
Institution, hvis Oprindelse, Bestaaen og Udvikling Historikeren vil
undersøge. Om det samler man alt, hvad der kan bidrage til dets
Belysning, men tager intet Hensyn til det, der hverken fremmer eller
hæmmer det. Korrelat til dette Begreb er Begrebet historisk Tærskel, som
betegner Grænsen for, hvad der kan tages med i den enkelte Undersøgelse.
Historisk Storhed betegner den Beskaffenhed af en Begivenhed, en
Personlighed eller en Institution, at der i dem ytrer sig Kræfter, som
virker i Retning af nye, betydningsfulde sociale Opgavers Løsning.
% p. 79: the three technical terms --- historisk Centrum, historisk
% Tærskel, Historisk Storhed --- are set ROMAN in the printing, neither
% italic nor spærret; no emphasis added.

Nu staar der i Historien som i Naturen en Kamp for Tilværelsen.
Historikeren staar overfor Begivenheder og
% --- p. 80 ---
% p. 80: running head "Nr. 3. Harald Høffding:" omitted as running matter.
% The page carries no footnote and no paragraph break: it is one
% continuous paragraph running from p. 79 to p. 81.
Personligheder i den menneskelige Verden, som Biologen staar overfor
Planter og Dyr i deres Kamp for Tilværelsen. Det gælder i begge Tilfælde
at forstaa de kæmpende Parter hver for sig, hvad deres Natur og Kaar
angaar, og derved at forstaa, hvorledes der opstaar Kamp imellem dem.
Biologen som saadan \guillemotleft holder ikke
med\guillemotright{} nogen af de organiske Former, selv om hans
intellektuelle Interesse har gjort en enkelt af disse Former til Centrum
for hans Undersøgelser, eller selv om han kan føle en vis Sympati for
visse Livsformer, som naar det glædede Darwin, naar han mente at kunne
finde Funktioner hos Planterne, der kunde pege paa et højere Livstrin,
end man plejede at tillægge dem. Det gælder for Biologen blot om at
forstaa Kampens Nødvendighed under de givne Forhold. Heller ikke
Historikeren som saadan tager Parti, men søger at finde, hvilke Forhold
der gør sig gældende hos hver af de stridende Parter, for at forstaa,
hvorledes Interesser, Tanker og Følelser har maattet udfolde sig i en
bestemt Grad og i en bestemt Retning hos hver af dem, indtil
Sammenstødet blev uundgaaeligt. Enhver Nation, siger William James i et
af sine Breve, har sine Idealer, der er en død Hemmelighed for andre
Nationer og maa under Indflydelse af disse Idealer udvikle sig paa sin
egen ejendommelige Maade. Hvis dette er saa, kræves der af Historikerne
Forstaaelse baade af, hvorledes selve disse Idealer har udviklet sig, og
af, hvorledes de har bestemt den nationale Udvikling. Derved bliver det
igen muligt at forstaa Spændingen og Striden mellem Nationerne. Det er
blevet sagt om Leopold Ranke, at naar han i sin historiske Forsken har
undgaaet ensidig Partitagen, var det ikke, fordi han var neutral, i
Betydning af Ligegyldighed overfor alle stridende Parter, men fordi han
var universel og kunde sætte sig ind i og følge
% --- p. 81 ---
det Liv, der rørte sig hos enhver af dem. Denne historiske Universalitet
er beslægtet med den kunstneriske Universalitet, der gjorde Shakespeare
i Stand til at lade enhver af hans Skikkelser udfolde sig efter deres
egne indre Love. Hans Upartiskhed overfor Personer og Skæbner udsprang
ikke af Ligegyldighed og Fjernhed, men af Grebethed af Livet, som det
kan røre sig hos forskellige Mennesker og føre dem til Sammenstød, der
maaske bliver tragiske.

Historikeren (og Digteren) staar her i Principet, som Fysikeren staar,
naar han konstaterer visse Forhold, under hvike det, der for en
Iagttager staar som samtidigt, for en anden Iagttager staar som
successivt.
% p. 81: PRINTER'S ERROR --- the printing reads "under hvike det", with
% the l dropped from "hvilke" (verified at 500 dpi). Transcribed as
% printed.

\subsection{Etiken og de psykologiske Relationer.}\label{sec:4b}
% p. 81: printed head reads "b. Etiken og de psykologiske Relationer."
% --- letterspaced (spærret) and RUN IN, like §4a and unlike §3's lettered
% heads: "b. Etiken og de psykologiske Relationer. --- Den historiske
% Vurdering...". The run-in dash is dropped and the following paragraph
% set \noindent. The text layer misreads the first word as "Eliken"; the
% image at 400 dpi shows "Etiken" plainly. Matches the Indhold.
% spacing.py's flags on this page ("Forskel mellem", Stand, Sammenstød,
% Relationer., Domme, Finhed, Sted,, Seele) are the head's own
% letterspacing plus justification; none is emphasis.
\noindent
Den historiske Vurdering er ikke egentlig en etisk Vurdering. Den angaar
Betydningen af en Begivenhed, en Personlighed eller en Institution for
den fortsatte faktiske Udvikling, men fælder ikke sine Domme ud fra et
etisk Ideal. Med stor Finhed og Energi har \textsc{Heinrich Rickert} i
sit Værk \guillemotleft Grenzen der naturwissenschaftlichen
Begriffsbildung\guillemotright{} paavist denne Forskel mellem Historie og
Etik. Og fra historisk Side er det energisk blevet udtalt, at den
historiske Forsken ikke trænger ind til det Sted, hvor Handlinger og
Bestræbelser har deres Kilde: \guillemotleft Das tiefinnerste Geheimnis
einer Seele zu finden, damit deren sittlichen Wert, das will sagen, den
ganzen Wert der Person richtend zu bestimmen, hat die historische
Forschung keine Methode und keine Kompetenz\guillemotright.
(\textsc{Droysen}: Geschichte Alexanders der Grossen, p.~259).
% p. 81: HEINRICH RICKERT and DROYSEN are SMALL CAPS in the BODY text, not
% in a footnote --- the first sites of this kind in the book (checked at
% 400 dpi). Neither the Rickert nor the Droysen title is italic: the first
% stands in guillemets, the second in roman inside parentheses.

Hvis alligevel det etiske Værdistandpunkt uvilkaarlig trænger sig frem
hos Historikeren, naar han prøver at give sin Fremstilling den sidste
Afrunding, vil han staa overfor
% p. 81: the page foot carries the signature line "Vidensk. Selsk.
% Filosof. Medd. I, 3." with the signature numeral 6 at the outer edge;
% omitted as running matter. Likewise the "6*" at the foot of p. 83.
% --- p. 82 ---
de principielle Vanskeligheder, der frembyder sig ved ethvert Forsøg paa
at begrunde etiske Domme, altsaa ved enhver Etik.

Paa dette Punkt optræder en ny Klasse af Kategorier, der ikke falder
sammen med nogen af de Klasser, --- de fundamentale, de formale, de
reale, --- som vi hidtil har omtalt, --- Værdibegreberne eller de ideale
Kategorier. Relationsbegrebet gør sig her straks gældende, idet Værdi
kun kan bestemmes i Forhold til Værdi. Al Villen er en Foretrækken, og
enhver Følelse faar sin Grad og sin Kvalitet bestemt ved Forholdet til
andre Følelser; derfor kommer det ved enhver Vurdering an paa, hvad der
forud eller samtidig staar som værdifuldt. Tilsidst vil al Vurdering i
de enkelte Tilfælde bero paa Forholdet til en Grundværdi. Ligesom et
Sted kun kan bestemmes ved sit Forhold til andre Steder, en Tid kun
gennem sit Forhold til andre Tider, en Virkelighed kun gennem sit
Forhold til andre Virkeligheder, saaledes forudsætter enhver
Værdibestemmelse, at visse andre Værdier forud staar fast. Grundværdien
svarer til den Klode eller det Koordinatsystem, hvorudfra vi orienterer
os i Rummet, --- til den Begivenhed, i Forhold til hvilken vi tidfæste
andre Begivenheder, --- til den Virkelighed, der i det enkelte Tilfælde
afgør, hvad der skal være Virkelighed for os. Og undersøger man nærmere
de forskellige mulige Grundværdier, der gør sig gældende ved
menneskelige Vurderinger, vil man finde, at de alle forudsætter en eller
anden Helhed, hvis Bestaaen og Udvikling kommer til at staa som Formaal
og derigennem bestemmer Vurderingen. En saadan Helhed kan være selve den
vurderende Personlighed eller et Samfund (Familie, Stand, Stat, Kirke)
eller et Indbegreb af Bestræbelser (Kunst, Videnskab o. s. v.). En
Vurdering vil tilsidst være bestemt
% --- p. 83 ---
\setcounter{footnote}{0}
ved en Handlings Forhold til Betingelserne for en saadan Helheds
Bestaaen og Udvikling. Kun naar der forudsættes en i Erfaringen given
Helhed, hvis Eksistens- og Udviklingsbetingelser kan drøftes, er en
rationel Vurdering, altsaa en rationel Etik mulig.\footnote{Se om
Sammenhængen mellem Totalitetsbegrebet og Vurderingsproblemet:
\textit{Den menneskelige Tanke}, p.~237---241; 346---369. \textit{Totalitet
som Kategori}. Kap.~5.}
Etisk Idealisme arbejder ofte sig selv imod ved at søge ud over alle
Relationer. De enkelte Steder, paa hvilke Platon omtaler
\guillemotleft det Godes Idé\guillemotright, hæves den udtrykkelig
(ligesom andre Ideer) ud over Relationernes Verden. Vel hævder han
(især i \guillemotleft Gorgias\guillemotright{} og \guillemotleft
Filebos\guillemotright) en Sammenhæng mellem det Værdifulde og bestemte
Maalforhold, nemlig Proportionalitet (\guillemotleft geometrisk
Lighed\guillemotright) eller Harmoni mellem Fylde og Begrænsning. Men
saadanne Maalforhold svæver i Luften, naar de ikke udledes af en vis
bestemt Helheds Livsbetingelser, og her svinger Platon mellem en
individuel og en social Totalitet.

Begrebet Norm angaar de Midler, ved hvilke en forudsat Helhed kan bestaa
og udvikle sig. Den forudsatte Helhed staar da mere eller mindre bevidst
som Formaal, idet dens Værdi betegner en Grundværdi. Begreberne Værdi,
Formaal og Norm ligger i samme Række, og i en rationel Etik maa de ligge
i denne bestemte Orden. Ingen Etik kan unddrage sig Sammenhængen mellem
disse tre Begreber. Ofte skjuler man det første Led i Rækken og holder
sig enten til Begrebet Formaal eller til Begrebet Norm, som om de var
primære. Vi møder atter her den Mangel paa Bevidsthed om, hvor det ene
Passerben staar, der saa ofte fører til Miskendelse af
Relationsbegrebets Betydning. Man betragter da det andet Passerbens
Stilling som absolut given, selvfølgelig eller \guillemotleft
objektiv\guillemotright. Ikke mindst fremtræder
% --- p. 84 ---
dette hos Etikere, der mener at have sikret sig et \guillemotleft
objektivt\guillemotright{} Stade ved at gaa ud fra \guillemotleft
Samfundet\guillemotright{} (uden engang at sige, hvilket Samfund de
mener). Eller det er, som hos Kant, Begrebet Norm, man lægger til Grund.
Kant, der i sin Erkendelsesteori er klar over, at enhver Nødvendighed er
hypotetisk, hævder en etisk Nødvendighed (et kategorisk Imperativ), der
ikke igen skal hvile paa nogen Forudsætning, og i forskellige Former har
senere Etikere fulgt ham heri. Naar et etisk System kritiseres, vil
Kritiken --- bortset fra en Undersøgelse om Systemets formale Konsekvens
--- væsentlig gaa ud paa at vise, at Systemet forudsætter en Grundværdi,
der ikke selv drages med ind i Undersøgelserne. Kritiken vil søge det
skjulte Passerben. Selv den, der proklamerer en \guillemotleft
Omvurdering af alle Værdier\guillemotright, maa forudsætte en eller
anden Værdi, i Forhold til hvilken han foretager Omvurderingen.

Grundværdien vil tilsidst være bestemt ved Selvopholdelsestrangen hos
det enkelte Individ (forsaavidt det kan tænkes isoleret) eller hos det
større eller mindre Samfund, med hvilket Individet føler sig solidarisk,
saa at dette Samfunds Ve og Vel, og dermed dets Formaal, ogsaa bliver
Individets og bestemmer dettes etiske Opgaver. Der kan dog ogsaa ytre
sig en Stræben efter at gøre det enkelte Øjeblik eller en enkelt Side af
Individets Liv til en Helhed i og for sig. Fra det enkelte Øjeblik til
Menneskeslægten i Nutid og Fremtid gives der en Række af mulige Grundlag
for etiske Systemer, og disse Grundlag er, som de etiske Ideers Historie
viser, tildels benyttede.

Den sammenlignende Etiker har nu (ligesom paa deres Omraade Fysikeren og
Historikeren) den Opgave at forstaa de Opfattelser, som de forskellige
Grundlag (Grundværdier) gør mulig, og som kan komme i afgørende Strid
% --- p. 85 ---
\setcounter{footnote}{0}
med hverandre. Derved forstaas den Kendsgerning, at samme Handling,
samme Personlighed, samme Institution kan være Genstand for modstridende
etiske Domme. Paa Grund af den lovmæssige Sammenhæng, der ad psykologisk
og historisk Vej kan paavises i den aandelige Verden, søger vi en
Forklaring af, hvorfor andre dømmer anderledes end vi, og altsaa har
andre Værdier, Formaal og Normer end vi.
% p. 85: spacing.py's singles (samme, Domme., Domme) are justification,
% not spærret --- checked on the image. No emphasis on this page except
% the note's title.

Og det er ikke blot den sammenlignende Etiker, der har Brug for saadan
Forstaaelse. Det gælder jo dog for ethvert etisk Standpunkt at faa
Mennesker til at kende og anerkende de Værdier, Formaal og Normer, det
selv lægger til Grund. Og der maa herved begyndes med at tage
Menneskene dér, hvor de faktisk staar, sætte sig ind i deres
Forudsætninger for at kunne lede dem til nye Forudsætninger. Den
sokratiske Ironi er, fra en væsentlig Side set\footnote{Smlgn.
\textit{Den store Humor}, p.~63.}, et Middel til at faa Mennesker til at
lukke sig op, til at aabenbare sig, for at det derigennem kan blive
klart, hvorledes Hjælp, Opdragelse og Samarbejde kan blive mulige. Det
er den analytiske Metode, som ene kan føre fra et Standpunkt til et
andet, naar det overhovedet er muligt. Det er naturligvis lettere at
begynde med Konstruktion af den \guillemotleft
rette\guillemotright{} Etik og saa overlade til Menneskene, hvorledes de
kan finde sig til Rette med den. Men enhver Værdi og enhver Vurdering
har i Kraft af Relationsbegrebet sine bestemte psykologiske
Forudsætninger, og det gælder først og fremmest at opdage, hvorvidt
disse er tilstede, saa at en fælles Vurdering kan blive mulig. Kun
fælles Grundværdi gør fælles Vurdering mulig.

Miskendelse af Forholdet mellem, hvad der er en Etikers Grundværdi, og
hvad han ud fra dette Grundlag føres til at lære, har ofte ført til
uretfærdige Domme om etiske
% --- p. 86 ---
\setcounter{footnote}{0}
Tænkere. Saaledes har man ment, at den Vægt, Filosofer som Adam Smith,
Helvetius og Bentham har lagt paa Erhvervstrangen, paa den individuelle
Interesse og paa Lykketrangen, skulde skyldes udelukkende Optagethed af
individuelle eller egoistiske Interesser, saa at deres egen Grundværdi
skulde bestaa i saadan Interesse. En nærmere Undersøgelse viser derimod
(som jeg har vist i vedkommende Afsnit af \guillemotleft Den nyere
Filosofis Historie\guillemotright), at de nævnte Tænkere alle gaar ud
fra universel Sympati og fra bevidst Hensyn til menneskelige Samfunds
Vel. Deres Opfattelse gaar ud paa, at Samfundet vil trives bedst ved, at
hver enkelt bygger paa sine Erfaringer om, hvad der tjener ham bedst, og
at ingen udenfor ham staaende Myndighed kan vide dette bedre end han
selv, naar han faar Øjnene op. Gennem det frie Arbejde og gennem de
selvstændig stræbende Individers Vekselvirkning vil for de nævnte
Tænkere den Grundværdi, de bygger paa, bedst komme til sin Ret.

Den enkelte Etiker maa være sig bevidst, at den Grundværdi, som bærer
hans System, selv kan blive dragen frem til Prøvelse. Hvis Systemet har
et saarbart Punkt, vil det tilsidst ligge her, selv om Systemet ellers
er udviklet med nok saa stor Konsekvens og støttet ved nok saa mange
Erfaringer. Og han maa, naar han er tilstrækkelig kritisk, give Henry
Sidgwick Ret i, at det --- i hvert Tilfælde paa visse Udviklingstrin ---
er godt, at der i et Samfund gøres forskellige, endog modstridende
etiske Opfattelser gældende, selv om han naturligvis kun kan følge sin
egen.\footnote{\textit{Methods of Ethics,}\textsuperscript{2} p.~450.}
% p. 86: the note gives no author (Sidgwick is named in the body, in
% ordinary roman). The raised 2 after the title is an EDITION mark, not a
% note reference --- \textsuperscript{}, as with Wissenschaftslehre² on
% p. 11. The printed italic runs through the comma, so the comma is set
% inside \textit{}.
Ogsaa paa det etiske Omraade kan Sandheden kun arbejde sig frem gennem
Strid. Psykologien og Historien har saa at
% --- p. 87 ---
sørge for, at der bag den brændende Strid kan være en underjordisk
Forstaaelse.

Den hele Opfattelse af Etikens Stilling, som her og i det Følgende
udvikles i Sammenhæng med Relationsbegrebet, har jeg fra 1887 af hævdet
i mine forskellige Skrifter, og den har stedse vundet ny Fasthed hos
mig, ikke mindst paa Grund af den Kritik, den fra forskelligt Hold har
været Genstand for.

\subsection{Etiken og de historiske Relationer.}\label{sec:4c}
% p. 87: printed head reads "c. Etiken og de historiske Relationer." ---
% letterspaced (spærret) and RUN IN, like §4a and §4b: "c. Etiken og de
% historiske Relationer. --- Udformningen og...". The run-in dash is
% dropped and the following paragraph set \noindent. Matches the Indhold.
% spacing.py's RUN "historiske Relationer." is the head's own
% letterspacing; its singles (sociale, sociale) are justification.
\noindent
Udformningen og Gennemførelsen af etiske Normer --- altsaa Bestemmelsen
af de underordnede Formaal, som Hovedformaalet fører til at stille, og i
hvilke de afledede Værdier, Grundværdien afføder, finder deres Udtryk
--- er kun mulig i Vekselvirkning med de givne historiske Betingelser.
Hverken de værdifulde Karakteregenskaber (\guillemotleft
Dyder\guillemotright) eller de sociale Ordninger, en Grundværdi kræver
for sin fulde Virkeliggørelse, kan konstrueres paa Forhaand; de kan ikke
\guillemotleft indføres\guillemotright, som man i gamle Dage mente at
kunne \guillemotleft indføre\guillemotright{} en Religion eller en
Forfatning. I hvert Tilfælde vil de historiske Betingelser til
forskellig Tid og paa forskelligt Sted give baade individuelle
Karakteregenskaber og sociale Ordninger væsentlige Ejendommeligheder,
der ikke kan udledes ene og alene af de almene Normer, Grundværdien
afføder. Allerede en kort Oversigt over Dydernes Historie, som den, jeg
har forsøgt at give i min Etik (X, 4), viser dette. For de sociale
Ordningers Vedkommende ses det af Ægteskabets (Etik XV) og det sociale
Spørgsmaals Historie (Etik XXV). De gamle Læderflasker kan ikke straks
erstattes med nye. Desuden har etiske Normer stedse en almen Karakter,
og Specialiseringen kan kun grundes paa de historiske Betingelser, under
hvilke de skal gennemføres.
% p. 87: the text layer's "mentq" is a speck under the e of "mente"
% (checked at 400 dpi); and its "1 hvert Tilfælde" is "I hvert Tilfælde"
% with a roman capital I. The page carries no footnote.
% --- p. 88 ---
\setcounter{footnote}{0}
Abstraktioner og vage Drømme er ikke nok. Det gælder foreløbig om, at
Retningen for Karakterudvikling og Samfundsordning bestemmes ved
Grundværdien. Et fuldstændigt Program kan kun Erfaringen føre til at
udarbejde. Det gaar med etiske Normer, som det gaar med Aarsagsbegrebet
i Erkendelsesteorien, der heller ikke i sig selv viser nogen Løsning af
specielle Spørgsmaal, men blot motiverer en Spørgen og en Søgen, der
faar en foreløbig Karakter paa forskellige Stadier af Videnskabens
historiske Udvikling.

Der ligger her Mulighed for tragiske Konflikter, idet de, der er
besjælede af en Grundværdi, kan mangle Blik for de historiske Forhold,
medens de, der staar fast paa historisk Virkeligheds Grund, kan mangle
Blik for Værdier, der kræver Virkeliggørelse. I Nutiden staar der mere
end nogensinde en Kamp mellem sociale Idealer og sociale Virkeligheder.
Bertrand Russell, for hvem det sociale Ideal er givet i Kommunismen,
finder, at den Pris, der ved Anvendelsen af bolschevikiske Metoder skal
betales af Menneskeslægten for at naa til Kommunisme, er altfor høj og
endog vil ophæve den Værdi, Kommunismens Gennemførelse i sig selv vilde
have. En Hovedaarsag hertil finder han i, at i Bolschevismen spiller
Hadet til det Gamle en større Rolle end Haabet om det Nye; ogsaa disse
psykologiske Momenter er historisk betingede.\footnote{\textsc{Bertrand
Russell}: \textit{The Practice and Theory of Bolschevism}. (1920),
p.~146---178.}
% p. 88: BERTRAND RUSSELL is SMALL CAPS in the note but ordinary roman in
% the body (both checked at 400 dpi). The body reads "Bolschevismen" with
% a v, not the layer's "Bolschewismen"; the title in the note likewise
% "Bolschevism".
Modstykket hertil findes mere end tilstrækkelig repræsenteret i Europa
som Had til det Nye og Haab om trods alt at kunne bevare det Gamle. Om
begge Modsætninger gælder det, at de intet har lært og intet glemt. Naar
den, der er opfyldt af et Ideal, tillige besidder intellektuel
Redelighed, vil han,
% --- p. 89 ---
\setcounter{footnote}{0}
for Idealets egen Skyld, undersøge alle Forhold, under hvilke det skal
virke, og alle Maader, paa hvilken denne Virken maa gribe ind i
menneskelige Forhold.
% p. 89: the printing reads "alle Maader, paa hvilken denne Virken" ---
% singular "hvilken" after a plural antecedent (verified at 500 dpi).
% Transcribed as printed.

Det forholder sig nu saaledes, at selve Grundværdierne og Normerne kun
bliver til under aldeles bestemte historiske Betingelser. Idealer, der
skulde være opstaaede og udformede uden Indflydelse af historiske
Virkeligheder, har aldrig eksisteret. Platon's etiske Idealisme dannede
sig ikke i hans Bevidsthed uden Indflydelse af det græske Aandslivs
Fortid, og hans Etik har, trods al Modsætning til hans Folks Sæd og til
hans Samtids Forestillinger, et afgjort græsk Præg. Kristendommen har i
sin oprindelige Form et afgjort orientalsk Præg, og den peger tilbage
til jødisk Tradition og til persiske Forestillinger. Kant's Etik har en
lang og betydningsfuld Udviklingshistorie. Naar jeg her ogsaa skal tale
om mig selv, har jeg udtrykkelig fremhævet, at de Forudsætninger, paa
hvilke jeg har bygget min Fremstilling af Etiken, skyldes historisk
Udvikling og ikke kan siges at have været tilstede til alle Tider. Naar
\textsc{Westermarck}\footnote{\textit{Origin and Development of Moral
Ideas}. I, p.~18.} siger, at hvis Ordet \guillemotleft
Etik\guillemotright{} skal bruges som Navn for en Videnskab, kan Opgaven
for den kun være at studere den moralske Bevidsthed som et Faktum, føler
jeg mig ikke truffen af denne Bemærkning, der bl.~a. ogsaa er rettet mod
mig; thi jeg har altid antaget en historisk-psykologisk Udvikling af den
moralske Bevidsthed og hævdet, at enhver Etik vil svare til en vis Form
og et vist Trin af denne Udviklingsproces. (Se Etik III, 13; 21). Dette
berøver ikke etiske Ideer deres Betydning eller deres Gyldighed. Man kan
jo dog styre efter Stjernerne, skønt de har en Historie. Og desuden
foregaar den etiske Sejlads paa
% ============================================================
%  BATCH 9 — §4d–§4f
% ============================================================
% --- p. 90 ---
de jordiske Have og betinges derfor ikke blot ved Stjernerne, men ved
Vind og Vove.

Opgaven maa stedse være en Udvikling af det historisk Givne i Retning
af de ideale Normer. Men der maa tillige ske en Orientering med Hensyn
til, om disse Normer nu ogsaa er opstillede med Rette, og om de nu
ogsaa virkelig er Udtryk for Værdierfaringer. Hvis der af en eller
anden Grund rokkes ved selve Grundværdien, bliver det et
psykologisk-historisk Spørgsmaal, om Overgang til en anden Grundværdi
er mulig. Ogsaa her er der Mulighed for alvorlige Kriser og tragiske
Katastrofer.

\subsection{Etiken og de individuelle Relationer.}\label{sec:4d}
% p. 90: printed head reads "d. Etiken og de individuelle Relationer."
% --- letterspaced (spærret) and RUN IN, like §4a, §4b and §4c, and unlike
% §3's lettered heads, which stand on their own line: the head sits at the
% ordinary paragraph indent and the text follows on the same line after an
% em dash, "d. Etiken og de individuelle Relationer. --- Foruden de
% psykologiske Relationer...". The run-in dash is dropped and the
% following paragraph set \noindent.
% VARIANT READING: the Indhold (p. 109) gives "d. Etik og de individuelle
% Relationer" --- without the -en. Both were checked on the images at
% 400 dpi and the difference is real, not an OCR artefact: the body head
% plainly reads "Etiken", the Indhold plainly "Etik". The printed body
% head is kept here; both readings are recorded.
\noindent
Foruden de psykologiske Relationer, der angaar Anerkendelse af en
Grundværdi og af de af den følgende Normer, og de historiske
Relationer, der angaar Muligheden af Gennemførelsen af disse Normer
under givne Betingelser, kommer der i en dybere gaaende etisk
Diskussion endnu en tredie Art af Relationer frem, som angaar det
enkelte bestemte Individs Mulighed for at fyldestgøre Normerne. Naar
etiske Krav ikke skal svæve i Luften, maa de stilles saaledes, at det
enkelte Individs Evne og Trang tages med i Betragtning. Kun da vil der
kunne udløses en Stræben i Retning af, hvad der kræves. Der kan derfor
ikke mekanisk kræves det Samme af Enhver. Af den, der er i Stand til
stor eller opofrende Virken, kræves der andet og mere end af den, der
under indre Splid eller formedelst Afmagt i Sindet maaske lige magter
at holde sig oven Vande. Kun da kræves der i Virkeligheden det samme
etiske Arbejde af alle; thi etisk Arbejde er Udtryk for den Energi,
med hvilken indre og ydre Modstand overvindes, og er denne Modstand
forskellig for forskellige Individers Vedkommende, kan det
% --- p. 91 ---
virkelig udførte Arbejde kun bestemmes ved Forholdet til denne
forskellige Modstand. Hvad der kræves af den Enkelte, kan da snart
ligge over, snart under de gennemsnitlige Krav, der sættes af en almen
Lov. Enhver bliver, som Schiller's Wilhelm Tell siger, beskattet efter
Evne.

I en Afhandling om Forholdslovene i Etiken (Etiske Undersøgelser,
1891) har jeg udførlig begrundet den Individualisering af etiske
Normer, der er en ligesaa nødvendig som vanskelig Konsekvens af
Relationsbegrebet. Den hviler paa det Princip, at Loven er til for
Menneskets Skyld, Mennesket ikke for Lovens Skyld. Det Arbejde, den
Enkelte maa gøre for at fyldestgøre etiske Normer, de Dyder og de
Pligter, der kræves af ham, maa selv tjene til Midler for hans
personlige Udvikling, og dette forudsætter, at hans Evne og Trang
bliver medbestemmende ved de Krav, der ifølge Normen bliver at stille
netop til ham.

De fleste etiske Teorier har miskendt denne individuelle Relations
Betydning. Alligevel har denne atter og atter gjort sig gældende. I
Teologien ligger den til Grund for Begrebet Naade, i Retslæren for
Begrebet Benaadning. Disse Begreber er nemlig uetiske, naar Hensynet
til de individuelle Betingelser ikke ligger til Grund for dem. I
Filosofien har man opstillet det analoge Begreb Billighed, der staar i
Modsætning til abstrakt Retfærdighed, men i Virkeligheden, hvis det
overhovedet er et etisk Begreb, maa siges at udtrykke den højeste
Retfærdighed, da det tager de individuelle Momenter med i Betragtning.
I den nyeste Tids Individualisering af Straffuldbyrdelsen kommer denne
Betragtning ogsaa frem; hvad en Lovovertræder dømmes til, er jo netop
den bestemte Straffuldbyrdelse, der faktisk anvendes paa ham.

Af de individuelle Relationers Betydning følger en vig%
% --- p. 92 ---
tig Konsekvens for social Etik. Det maa være en nødvendig Opgave at
stræbe efter at skaffe alle Medlemmer af Samfundet saadanne Muligheder
for Udvikling, at Valget af Arbejde kan blive frit. Saalænge dette
ikke er naaet, er den etiske Fordring til Samfundet ikke sket Fyldest.
Paa dette Punkt vil mere og mere den sociale Kamp komme til at staa.

\subsection{Kant's Etik og Relationsbegrebet.}\label{sec:4e}
% p. 92: printed head reads "e. Kant's Etik og Relationsbegrebet." ---
% letterspaced (spærret) and RUN IN, like §4a--§4d: "e. Kant's Etik og
% Relationsbegrebet. --- Kant, der selv...". The run-in dash is dropped
% and the following paragraph set \noindent. Matches the Indhold (p. 109)
% word for word. The apostrophe in "Kant's" is printed as a right single
% quote, as everywhere else in this book.
\noindent
Kant, der selv, som vi ovenfor har set, i Relationsbegrebet saa en
Grundlov for al Tænken, vilde dog ikke anerkende det i Etiken. Den
etiske Lov skal ifølge ham ikke i nogensomhelst Henseende eller Grad
være fremgaaet af Erfaring; psykologisk og historisk Relation skal
altsaa ikke gælde for den. Og den skal kræve det Samme af alle, hvilke
individuelle Forudsætninger de saa har. Men derfor staar den etiske
Lov ifølge Kant ogsaa som et uforklarligt Vidunder overfor den
arbejdende, tænkende og stræbende Menneskehed, der stadig er stillet
overfor bestemte Betingelser. Kant udtaler aldeles aabent dette.
\guillemotleft Hvorledes ren Fornuft kan være praktisk, er menneskelig
Fornuft aldeles ude af Stand til at forklare, og al Møje, alt Arbejde
for at søge en Forklaring deraf er spildt\guillemotright.
(Grundlegung\textsuperscript{3}, p.~125).
% p. 92: the raised 3 on "Grundlegung" is an EDITION mark, not a note
% reference: p. 92 carries no note. Set with \textsuperscript, as with
% Wissenschaftslehre² on p. 11.
\guillemotleft Hvorledes en Lov i og for sig og umiddelbart kan være
Bestemmelsesgrund for Villien, det er et for den menneskelige Fornuft
uløseligt Problem\guillemotright. (Kritik der praktischen Vernunft,
p.~128).
% p. 92, printer's defect: a small stray mark --- a broken or turned piece,
% standing at mid x-height and clear of both words --- sits between "det"
% and "er" in "det er et for den menneskelige Fornuft". The text layer
% reads it as a full stop ("det .er"). It is not a punctuation mark;
% verified at 900 dpi. Transcribed as no character, and logged here.
Vel indrømmer Kant, at Loven kun kan virke i os ved at fremkalde en
Følelse af selve Lovens Ophøjethed; men han erklærer saa denne Følelse
for psykologisk uforklarlig.

Men af dette i psykologisk, historisk og individuel Henseende
Relationsløse følger igen, at vi i Grunden ingen
% --- p. 93 ---
\setcounter{footnote}{0}
egentlige etiske Domme kan fælde. Naar vi ingen Adgang har til at
forstaa Forholdet mellem Norm, Motiv og Handling, kan vi aldrig have
Ret til at kalde nogen foreliggende menneskelig Handling god eller
ond. Denne Konsekvens drog Kant allerede i første Udgave af
\guillemotleft Kritik der reinen Vernunft\guillemotright{} (p.~551),
dog kun i en Note under Teksten: \guillemotleft Die eigentliche
Moralität der Handlungen bleibt uns daher, selbst die unseres eigenen
Verhaltens, gänzlich verborgen\guillemotright.\footnote{Maaske skyldes
den i samme Retning gaaende Ytring af \textsc{Droysen} (se ovenfor
p.~81) for en Del Indflydelse af Kant's Tankegang.} Dette er en rigtig
Konsekvens, naar Relationsbegrebet ikke skal kunne anvendes i Etiken.
Anderledes stiller Sagen sig, naar etiske Normer opfattes som Udtryk
for Bestræbelser, der netop udformes til fuldkomnere Former gennem de
nødvendige Relationer, som gør sig gældende ved al Tænken og Handlen.
Netop Relationsbegrebet aabner os her store Muligheder. Det kræver, at
vi sætter det ene Passerben paa et bestemt Sted i den virkelige Verden
(i et psykologisk Faktum, i et historisk Givet, i en individuel
Mulighed) og saa prøver, hvor langt vi kan føre det andet Passerben
ud. Kan vi ikke lige straks sætte det andet Passerben, hvor vi vil,
saa kan vi maaske efterhaanden flytte det længere ud. Maaske kan endog
det første Passerben flyttes. Det gælder i Etiken som i Fysiken om
bestemte points de repère.
% p. 93: "points de repère" is set in ORDINARY ROMAN here, exactly as on
% p. 73 --- no italic in the text layer's font tags and none on the image.
% So no \textit, in keeping with the earlier decision at p. 73.

Ligesom Rummet og Tiden er \guillemotleft Ideer\guillemotright{} (i
Kant's Betydning af Ordet), og ligesom Erfaring (Virkelighed) er en
\guillemotleft Ide\guillemotright, saaledes er ogsaa det etiske Gode en
\guillemotleft Ide\guillemotright, som faar real Betydning, ikke ved at
udelukke alle Relationer, men ved at benytte dem til at stille de
specielle Opgaver og til at bedømme Forsøgene paa deres Løsning.

Kant og Droysen har dog ikke ganske Uret i deres Er%
% --- p. 94 ---
klæringer om, at Handlingers egentlige Moralitet ikke kan opdages
eller begrundes. Thi vi kan aldrig være sikre paa, at vi har taget
alle Relationer med i Betragtning. Det er tilsidst umuligt at
fyldestgøre det Krav, der netop i Kraft af Relationsbegrebet maa
stilles, nemlig at vinde Indsigt i Forholdet mellem Handlingens indre
og ydre Betingelser og i Størrelsen af det etiske Arbejde, det enkelte
Menneske ud fra sine individuelle Betingelser har øvet eller kan øve.
Alle etiske Domme er derfor ufuldkomne og gælder i det Højeste kun
tilnærmelsesvis.

Ogsaa Kant's Antagelse af en for alle Mennesker gældende Norm kan
stadig hævdes. Thi Relationsprincipet bliver i Grunden Et med
Personlighedsprincipet, d.\ v.\ s. det Princip, at ethvert Menneske
stedse skal betragtes og behandles som Formaal, aldrig blot som
Middel. Dette fyldestgøres, naar Kravene til den Enkelte ikke blot
motiveres socialt, men ogsaa individualiseres efter den Enkeltes Evne
og Trang, saa at hans Personlighed kan udvikles gennem det Arbejde,
han gør for at følge Normen.

To af Kant's betydningsfuldeste etiske Tanker kommer saaledes først
gennem Anvendelse af Relationsbegrebet til deres fulde Ret.

\subsection{Religionsfilosofien og Relationsbegrebet.}\label{sec:4f}
% p. 94: printed head reads "f. Religionsfilosofien og
% Relationsbegrebet." --- letterspaced (spærret) and RUN IN, like
% §4a--§4e. The em dash closing the run-in stands at the end of the head's
% own line and the text resumes on the next: "f. Religionsfilosofien og
% Relationsbegrebet. --- / Ogsaa det religiøse Problem...". (The text
% layer drops this dash; it is plain on the image at 400 dpi.) The run-in
% dash is dropped and the following paragraph set \noindent. Matches the
% Indhold (p. 109) word for word.
\noindent
Ogsaa det religiøse Problem, hvorledes man saa stiller sig overfor det,
frembyder flersidig Anvendelse af Relationsbegrebet. Dette Problem
% p. 94, printer's defect: a small raised stray mark stands between
% "Relationsbegrebet." and "Dette Problem" (the text layer reads it as a
% hyphen, "grebet.- Delte"). Verified at 900 dpi: it is a broken or turned
% piece at mid x-height, not a punctuation mark. Transcribed as no
% character, and logged here.
beror i alle sine Former paa en Modsætning mellem det Værdifulde paa
den ene Side, det Værdiløse eller Værdihæmmende paa den anden Side,
--- kortere udtrykt mellem Værdi og Virkelighed, idet Virkeligheden
frembyder baade Værdi og det Modsatte. Den Modsætning, her foreligger,
% p. 94: "Den Modsætning, her foreligger" --- the relative pronoun is
% absent, as printed; transcribed unaltered.
behøver --- efter min Opfat%
% --- p. 95 ---
\setcounter{footnote}{0}%
telse --- ikke at være en Dualisme, som man har
ment\footnote{\textsc{La Harpe}: \textit{La Religion comme conservation
de la valeur, dans ses rapports avec la philosophie générale de Harald
Høffding} (1920), p.~93 ff.}; men den kan ganske vist stige til en
Katastrofe, medføre den Livets Tragik, som ingen Livsanskuelse kan
udelukke Muligheden af, ligesaa lidt som man kan udelukke Muligheden
af Jordskælv. Og selv om man i det Tragiske, i den Lidelse, der fører
til Grænsen af det Menneskelige, vilde finde Tilværelsens sande Natur
udtrykt, saa vilde det netop bekræfte Modsætningers Betydning for
Livsopfattelsen. Det Tragiske forudsætter Relation i skarpeste
Forstand, uden dog at udelukke Fastholden af det Værdifulde. Den
tragiske Helt hævder Menneskehedens Adel i den største Katastrofe. Men
ogsaa bortset fra de store Katastrofer fremtræder Forholdet mellem
Værdi og Virkelighed som det Problem, Religion under alle Former
hviler paa. I alle Former for Religion gælder det om, hvad der er det
Sejrende i Verden, det Gode eller det Onde. Indenfor sin store Fylde
frembyder Virkeligheden Mulighed for Udvikling i forskellige
Retninger, og Brydningen mellem Kræfter, der virker i Retning af, hvad
Mennesker tillægger højest Værdi, og andre Kræfter, der virker i helt
anden, maaske modsat Retning, stiller netop de Spørgsmaal, Religion i
en eller anden Form søger at besvare. Det er denne Brydning, der fører
den religiøse Bevidsthed til at danne to Afslutningsbegreber, Gud og
Verden, og al religiøs Tænken er en Grublen over Forholdet mellem
disse to Begreber. I Askese og i Mystik søger man, snart mere i
Praksis, snart mere i Spekulation at naa ud over alle Relationer,
eller (som vi har set ovenfor p.~47 Note) dog at faa det ene
Forholdsled til at forsvinde; men atter og atter trænger Relation sig
frem og stiller Problemet paany. Livets Trang forstyrrer Asketen,
% --- p. 96 ---
og Virkelighedens Verden trænger sig ind i Mystikerens Celle med sin
Mangfoldighed og sin Uro. Man opdager tilsidst det ene Passerbens
Stilling, selv om det var nok saa godt skjult.

Spørgsmaalet er, om det Værdifulde i Verden (den værdifulde Del af
Virkeligheden) er en fremmed Gæst, som ikke har sin Hjemstavn der,
eller om det ikke netop skulde være en Afsløring af Tilværelsens
inderste Natur. Det eneste rationelle Svar, som her kan gives (og som
jeg har forsøgt at give i \guillemotleft Totalitet som
Kategori\guillemotright{} Kap.~5), maa gaa ud fra den nøje Forbindelse
mellem Begreberne Værdi og Helhed. Alt, hvad der tillægges Værdi, er
Udtryk eller Betingelse for en vis Helheds Bestaaen og Udvikling, og
det Værdifulde opstaar ifølge de samme Love, som gør Helhedsdannelser
(f.\ Eks. Organismer, Personligheder, Samfund) mulige.
(Smlgn.\ ovenfor p.~83 f.). Naar nu Helhed staar mod Helhed, eller naar
kaotiske Elementer gør Helhedsdannelse umulig, saa foreligger netop de
Erfaringer, der fører til at skelne mellem Værdi og Virkelighed.

Naar de monoteistiske Religioner betragter Gud som Ophav til al
Virkelighed, optræder det Ondes Problem med særlig Skarphed. Udenfor
Ophavet til al Virkelighed kan det Onde ikke have sin Oprindelse.
Derfor har dristig og dybsindig Spekulation, som Jacob Böhme's, draget
den Konsekvens, at der maa være en dunkel Grund i Gud selv, hvoraf det
Værdistridige i Verden er opstaaet. Der sættes da indre Relationer i
Gud i Stedet for de ydre Relationer, Erfaringen standser ved. Mere
indviklet bliver Forholdet endnu, naar det, som allerede Böhme
erkendte, kan være en Betingelse for det Godes Udvikling, at det Onde
øver sin Modstand og derved virker til Udløsning af Kræfter, der
ellers vilde ligge brak. En Betragtning af denne Art
% --- p. 97 ---
\setcounter{footnote}{0}
førte Goethe (Aus meinem Leben.\ Fjerde Del) til at fremhæve, hvad han
kaldte det Dæmoniske i Verden, og til at bruge det Motto: nemo contra
deum nisi deus ipse.\footnote{Se nærmere herom
\textit{Religionsfilosofi}, §~88---89.}

Religiøs Tro er et Haab eller en Forvisning om, at der stedse i
Tilværelsen (i den, som Erfaringen frembyder, eller i en anden) maa
være Mulighed for værdifulde Helhedsdannelser. Det er dette, der
udtrykkes ved, at Religion er Tro paa Værdiens Bestaaen. Men ligesom
Helheder opstaar og forgaar, saaledes ogsaa Værdier. Den religiøse
Bevidsthed kommer derfor i Løbet af sin Udvikling til at indse, at det
ikke kan være det sanselig eller empirisk Værdifulde, der bestaar, ja,
at det tilsidst er umuligt at give et Begreb om eller et Billede af,
hvad det er, der bestaar. Denne Erkendelse arbejder sig ikke frem uden
stærk Modstand og kræver en dyb Resignation. Den tilsløres let gennem
det Billedsprog, Religionen anvender.

Selve Begrebet Gud indeholder en Vurdering; Menneskers Gud er det, de
sætter højest. I Udtrykket \guillemotleft
guddommelig\guillemotright{} ses dette tydelig; det er en
Værdibestemmelse. Allerede Platon taler om \guillemotleft det, der gør
Gud guddommelig\guillemotright, hvilket er, at Gud stadig lever i de
store Ideer, hvad Mennesker kun af og til kan (Fajdros p.~249 C).
Plotinos gaar et Skridt videre, idet han siger, at det, der gør
Guderne guddommelige, dog selv maa være den højeste Guddom (Ennead.\ V,
1, 2). Fortsættelsen af denne Tankegang vilde føre al Teologi tilbage
til en Værdilære.\footnote{Smlgn.\ hermed den Maade, paa hvilken
Udtrykket \guillemotleft guddommelig\guillemotright{} er brugt af en
moderne Mystiker (\textit{Oplevelse og Tydning}, p.~89).} ---

Udtrykket \guillemotleft Værdiens Bestaaen\guillemotright{} siger intet
om, hvor megen Værdi der er i Verden, eller af hvilken Art den er. Paa
disse to Punkter ligger Forskellighederne mellem de
% p. 97: the page foot carries the printer's signature line
% "Vidensk. Selsk. Filosof. Medd. I, 3." with the signature numeral 7 in
% the outer corner; omitted as printer's apparatus, like the earlier ones
% at pp. 33, 49, 65 and 81.
% --- p. 98 ---
forskellige Religioner eller religiøse Standpunkter, som det her ikke
er Stedet at gaa ind paa; det maa overlades til Religionshistorie og
Religionspsykologi. For Relationsteorien var det her kun Opgaven at
vise, at Relationsbegrebet uvilkaarlig gør sig gældende i al religiøs
Tænken, den mest elementære saavel som den mest spekulative. Det er,
som vi saa, dels Forholdet mellem Værdi og Virkelighed, dels Forholdet
mellem evige og forgængelige Værdier, der stiller Spørgsmaalene for
religiøs Tænken. Nye Relationer opstaar under Religionernes historiske
Udvikling og trænger ofte de til Grund liggende Relationer tilbage.
Saaledes Forholdet mellem en Religions egentlige Kærne og
Udformningen i Kultus og Dogme. I sin Iver for at faa Garantier for
Værdiens Bestaaen har den religiøse Bevidsthed Tendens til at
mangfoldiggøre Troens Genstand, saa at det ene Dogme bruges til at
støtte det andet. --- En særlig vigtig Relation er den mellem Religion
og Moral. Bevidstheden kan være saa optagen af Værdiens Bestaaen og af
de Dogmer, der menes at garantere den, at der drages Energi bort fra
Arbejdet paa at opdage og frembringe Værdier eller paa at hævde de
allerede fundne Værdier. Man kan dvæle saa længe paa Villiens Grænse,
at man ikke kommer til at virke paa de Omraader, hvor Menneskene ikke
behøver at vente paa Guderne. Og tilsidst trænger fra et etisk
Synspunkt den Betragtning sig frem, at Værdi dog ikke nødvendigvis er
afhængig af Bestaaen. Der kan leves i og for Værdier, selv om ingen
Vished kan vindes om deres Bestaaen. Det Guddommelige er ikke
afhængigt af Tiden. Hvad der vurderes som skønt og godt, kan beholde
sin Værdi, hvilken dets Skæbne end maatte blive.


% ============================================================
%  BATCH 10 — §5 Relationsbegrebet i Kosmologien (head at p. 99); Indhold
% ============================================================
% --- p. 99 ---

\section{Relationsbegrebet i Kosmologien.}\label{sec:5}
% p. 99: printed head reads "5. Relationsbegrebet i Kosmologien." --- bold,
% centred, on its own line, the numeral supplied by the counter. Matches the
% Indhold (p. 109) word for word. p. 99 therefore OPENS a new paragraph.
Tilsteder Relationsbegrebets Gennemførelse en afrundet, om end ikke
afsluttet Opfattelse af Tilværelsen? Vil enhver Helhedsopfattelse her ikke
vise sig umulig overfor et Kaos af Synspunkter, der hvert for sig gør Krav
paa Berettigelse? --- Overfor saadanne Spørgsmaal er der forskellige
Betragtninger at fremføre.

a.~Ingen Orientering er mulig uden ud fra et bestemt Stade. Men ethvert
% pp. 99, 103, 105, 107: §5 is divided by BARE LATIN LETTERS a. b. c. d.,
% set in ordinary roman at the ordinary paragraph indent with the text
% running straight on --- no title, no run-in dash, no letterspacing, no
% line of their own. They are lettered paragraphs, not heads, and the
% Indhold (p. 109) lists no sub-entries whatever under §5. So they are kept
% inline and given no sectioning command and no \addcontentsline; there is
% no title anywhere in the book to give them. Compare §3's bare Greek marks.
Stade kan igen drages ind i Drøftelsen, fordi det staar i bestemte
Relationer og hviler paa bestemte Forudsætninger. Der er en stadig Tendens
til at overse dette. Baade den sunde Menneskeforstand, de religiøse
Verdensanskuelser og altfor ofte ogsaa naturvidenskabelige og filosofiske
Systemer er gaaet ud fra, at der gives Standpunkter, hvis Begrundelse ikke
igen kan fordres, og hvis Berettigelse ikke kan drøftes; hvad der er
Grundlag for Verdensanskuelse, skal selv ikke være betinget. Hver
Verdensanskuelse har sit Urfænomen, for at bruge Goethes Udtryk, og Goethe
paastod udtrykkelig, at Urfænomenet ikke kan forklares, men maa betragtes
som et absolut Sidste. Men hvor man saa flygter hen, vil Relationsbegrebet
kunne gøre sig gældende, og der kan opstaa den Uro, som er uadskillelig fra
al arbejdende Tanke. Tænkningens Idyller er altid udsatte for at opløses og
kan kun være foreløbige Hvilesteder. I nyeste Tid begynder man at blive
klar over, at en rationel Forklaring af alle Tilværelsens Dele eller Sider
ud fra et eneste Urfænomen ikke er mulig. Men ofte hævder man da, at der
her tilsidst maa træffes et Valg mellem to Muligheder: mellem et
matematisk-mekanisk og et psykologisk-kvalitativt Urfænomen.
Naturvidenskaben nærmer sig mere og mere til det matematiske Ideal af en
Identiteternes eller Substitutionernes Verden, hvor Tids- og
% p. 99 foot: the signature mark "7*" stands below the last line; not
% transcribed. Readings verified at 300 dpi against the ABBYY layer's
% t/l confusions: "tilsidst" (layer "tilsidsl"), "psykologisk-kvalitativt"
% (layer "psvkologisk-kvalilativt"), "det matematiske" (layer "del").
% --- p. 100 ---
\setcounter{footnote}{0}
Kvalitetsforskelligheder ikke spiller nogen Rolle, medens Aandsvidenskaben
(eller, som nogle hellere vil kalde den, Kulturvidenskaben) hævder
Realiteten af Kvalitetsforskelle og Udviklingsstadier, der udelukker
Reduktion til Identitet. Mellem de herved bestemte to Typer skulde Valget
staa. I denne Retning har Wundt, James Ward og Henri Bergson udtalt sig.
Men at et saadant Valg skulde være nødvendigt, vilde forudsætte, at der her
foreligger et Enten---Eller, at der altsaa ikke eksisterer andre
Tilværelsestyper end de to antydede. Den Omstændighed, at Identitet,
Kvalitet og Temporalitet ikke kan reduceres til Et, berettiger ikke til et
Enten---Eller. Et er Irreduktibilitet, et Andet er Fraværelse af alle andre
Muligheder. Overfor en Paastand om et Enten---Eller godtgør sig her
% PRINTER'S DEFECT: in this third occurrence a small raised tick, sloping
% like an acute accent, stands immediately after "Eller" at cap height
% (verified at 600 dpi; the ABBYY layer reads it as "Ellei7"). It is a
% stray sort or a broken piece of type, not punctuation --- the sentence
% runs straight on into "godtgør" --- and it cannot be set, so it is logged
% here rather than transcribed.
Betydningen af Spinoza's mærkelige Lære om de uendelig mange Attributer,
under hvilke Tilværelsen fremtræder, men af hvilke vor Erfaring kun viser
os to, --- netop hine to Typer, mellem hvilke Valget skulde staa. I langt
højere Grad, end han er sig det bevidst, peger Spinoza her hen til vor
Erkendelses Begrænsning.\footnote{\textit{Spinoza's \guillemotleft
Ethica\guillemotright}, p.~22---24.} Naar Spinoza, trods den Mangfoldighed
af Synspunkter, hvis Mulighed og Berettigelse han anerkendte, alligevel
hævdede en Enhedslære, var det, fordi ogsaa han havde sit Urfænomen, som
han var overbevist om ikke kunde fornægtes fra noget af de mange
Synspunkter. Han holdt sig til den Kendsgerning, at vi har en Videnskab,
--- at rationel Forstaaelse af Verden er mulig i den Udstrækning, den er.
Det var denne Kendsgerning, han projicerede ud som en absolut Realitet og
kaldte Substansen i Alt. Og den er jo i ethvert Tilfælde et Tilhold, der
gør det muligt at drøfte de forskellige Synspunkter, der kan gøre sig
gældende, navnlig ogsaa med Hensyn til, om deres indbyrdes
% p. 100 note 1: the whole of "Spinoza's »Ethica«" is italic (300 dpi);
% "Spinoza's" is NOT small caps here --- the ABBYY box heights show a
% single face, unlike HOEFER on p. 103. The note is \textit{} (a title),
% not \emph{}.
% --- p. 101 ---
\setcounter{footnote}{0}
Forhold er kontradiktorisk eller blot kontrært. Selve de Grænser,
menneskelig Erkendelse altid er underlagt, kan kun bestemmes ved
videnskabelig Tænken. Og her viser det sig, som flere Gange bemærket i det
foregaaende, at Nødvendigheden af Relationer baade kan være en Skranke og
en Spore for Erkendelsen. Hvor vi ingen bestemte Relationer kan finde, har
vi heller ingen Erkendelse. Vor Søgen hører derfor ikke op; vi ved, hvad vi
mangler: nye Stillinger for vore Passerben. Derfor ender vi med en Spørgen.
% p. 101: spacing.py's RUN flag on "Relationer baade" is justification, not
% spærret --- the line is a loose one in an otherwise ordinary paragraph
% and italics.py sees nothing there. No emphasis set.

Man har nu endog ment, at Relationsbegrebets Gennemførelse vilde gøre al
Sandhed umulig. Bergson skelner mellem relativ og begrænset Erkendelse; den
første forvansker sin Genstands Natur, den anden lader den uberørt, idet
den kun griber en Del af den.\footnote{\textit{Bulletin de la socièté
française de Philosophie.} VIII, p.~341.}
% PRINTER'S DEFECT: the note prints "socièté" --- GRAVE accent on the first
% e, acute on the second (checked at 600 dpi against the acute of
% "Poincaré" three lines above, which is unmistakably the other way).
% The journal is of course the Bulletin de la société française de
% Philosophie; the compositor's reading is kept as printed.
Men med Relationerne falder ogsaa Problemerne bort; der bliver hverken
bestemte Spørgsmaal at stille eller bestemte Svar at give. Konsekvent
henviser Bergson derfor til en umiddelbar Skuen, en Intuition som eneste
Vej, ad hvilken Sandheden kan naas og Tilværelsens inderste Væsen være
given. I den nærmere Gennemførelse heraf benytter han sig dog af stadige
Analogislutninger --- og forsaavidt ogsaa af Relationer, da Analogi jo
betyder Forholdslighed.\footnote{\textit{Henri Bergson's Filosofi},
p.~54.} Et Eksempel paa, hvorledes Relation trænger sig ind, netop som man
har ment at kunne udelukke den! --- Den vigtigste Indvending mod Bergson's
Filosofi er fra et filosofisk Synspunkt ikke den, Poincaré etsteds har
udtrykt i Sætningen: Le monde Bergsonien n'a pas des lois! --- men den, at
denne Filosofi ingen Problemer har. Den er egentlig en Lovprisning af den
problemløse Tilstand.
% p. 101: "Le monde Bergsonien n'a pas des lois!" is set ROMAN, not italic
% (italics.py sees nothing; confirmed at 300 dpi), and is transcribed as
% printed --- "n'a pas des lois" for "n'a pas de lois" is Høffding's or the
% compositor's French, not an error to repair.

% --- p. 102 ---
Den menneskelige Tankes Liv ytrer sig i de Relationer, den formaar at
fremdrage og at bestemme. Enhver saadan Bestemmelse betyder en Sejr over et
Kaos. I en kaotisk Forskelsrække er ingensomhelst bestemte Relationer
paaviselige; absolut Ombyttelighed raader. At Fremdragen og Bestemmen af
Relationer er Tankens væsentlige Egenskab, følger, som flere Gange omtalt i
det foregaaende, deraf, at Relation er Korrelat til Syntese. Tankens,
Bevidsthedens Enhed og Energi ytrer sig i de Relationer, der lader de
enkelte Emner fremtræde i deres Ejendommelighed. En Sammenfatten, ved
hvilken ingen Relationer sættes, er en Modsigelse. Selv i Drømme raader
Forholdsbestemmelsens Lov, saa meget de end kan nærme sig til Kaos.

Det er heller ikke i og for sig en Ufuldkommenhed ved vor Erkendelse, at
den til forskellige Tider og paa forskellige Steder, alt efter de
forskellige Udgangspunkter, kommer til forskellige Resultater. Derved
stilles der blot et nyt Problem, en ny Opgave for Erkendelsen, den nemlig
at udfinde, hvorledes de forskellige Udgangspunkter og dermed de
forskellige Resultater blev mulige. I det foregaaende har vi set, hvorledes
Fysikeren og Historikeren indser Muligheden af forskellige Stader og deraf
følgende forskellige Opfattelser. Det er en Opgave, som naturligvis ikke
altid kan løses, men som nødvendigvis frembyder sig, da Sandheden baade maa
kunne forklare sig selv og Fejltagelserne; det hører med til fuldstændig
Sandhed. Platon har allerede set dette, naar han (Staten X, p.~602
C---603 A) viser, at naar forskellige Iagttagere opfatter samme Ting paa
forskellig Maade, hvad Størrelse og Figur angaar, saa kan man ved Tælling
og Maaling komme til en Opfattelse, ud fra hvilken de forskellige
individuelle Opfattelser kan forklares. Galilei har været inspireret af
denne Tanke ved
% p. 102: the Stephanus reference reads "602 C---603 A" (verified at
% 300 dpi); the ABBYY layer's "602 G" is an OCR misreading.
% --- p. 103 ---
\setcounter{footnote}{0}
Grundlæggelsen af sin Bevægelseslære. Men han kom i Forlegenhed ved, hvad
hans ufuldkomne Teleskop viste ham om det Fænomen, man nu kalder Saturns
Ring. Fænomenet viste sig snart som Kanter eller Hjørner paa Planeten,
snart som to forskellige Biplaneter. Tilsidst opgav Galilei al Forsken
herom. Senere saa Huyghens med sine forbedrede Teleskoper tydelig, at
Saturn var omgiven af en Ring, og viste, hvorledes tidligere Opfattelser
kunde forklares ud fra denne.\footnote{\textsc{Hoefer}: \textit{Histoire de
l'Astronomie}, p.~445---447.} Det er Tankens Triumf paa alle Omraader, naar
der kan vindes en \guillemotleft fuldkommen\guillemotright{} Erkendelse, ud
fra hvilken de \guillemotleft ufuldkomne\guillemotright{} Opfattelser kan
forklares. Naar Spinoza lærer (Eth. II, 36. Dem.), at vore
\guillemotleft ufuldkomne\guillemotright{} Forestillinger opstaar med samme
Nødvendighed som de \guillemotleft fuldkomne\guillemotright, saa er det
netop ved den \guillemotleft fuldkomne\guillemotright{} (ɔ: fuldkomnere)
Forestilling, at denne Nødvendighed indses. Det er Meningen med Spinoza's
berømte Sætning: veritas est norma sui et falsi (Eth. II, 43. Schol.).
% p. 103: the printing uses the Danish abbreviation mark "ɔ:" (U+0254,
% LATIN SMALL LETTER OPEN O, followed by a colon) for "det er"; verified at
% 300 dpi, where the ABBYY layer reads "ø;". Also "Det er Tankens Triumf"
% (layer "l)et") and "at vore" (layer "al vore").
% p. 103: the Latin "veritas est norma sui et falsi" is set ROMAN, not
% italic and not spærret (italics.py sees only the footnote; confirmed at
% 300 dpi). No emphasis set.

Tanken stiller altsaa sine Grænseproblemer i Kraft af sin egen Lov, og det
er i Kraft af denne samme Lov, at den forklarer sine Vildfarelser.

b.~Men selve den Tællen og Maalen, som Platon henviste til, og som Galilei
gjorde principiel Brug af, skyldes jo dog menneskelig Aandsvirksomhed, og
hvilken Ret har vi til at tillægge dens Resultater Gyldighed? Vi vil jo dog
forstaa Verden, som den er, og ikke som den fremstiller sig for menneskelig
Beregning, selv hvor denne forsøger at anlægge nok saa væsentlige
Synspunkter.

Allerede ovenfor (p.~58 ff.) har jeg overfor en saadan Betragtning
indskærpet, at Relationsbegrebet ikke blot maa anvendes mellem Emner
indbyrdes, men ogsaa mellem Emnerne og det erkendende Subjekt, eller
rettere udtrykt,
% p. 103 note 1: HOEFER is SMALL CAPS (the ABBYY box heights split it as
% "H" + "oefer", the standard small-caps signature, and the image confirms
% it at 300 dpi); the title is italic. spacing.py's flags on this page
% ("Ring. Fæno-", "ogsaa") are justification, not spærret.
% --- p. 104 ---
at det erkendende Subjekt selv er et Emne. Ogsaa \guillemotleft the
percipient event\guillemotright, for at bruge Whitehead's Udtryk, hører med
til Begivenhederne i Verden. Vi har ikke Ret til paa Forhaand at hævde en
absolut Modsætning mellem Objekt og Subjekt. Selv om man søger at forvandle
al Videnskab til formal Videnskab og henregner alt, hvad der ikke kan
bringes i denne Form, til \guillemotleft
Antropomorfismer\guillemotright, saa er disse Antropomorfismer nu en Gang
et Faktum, som ogsaa kræver sin videnskabelige Undersøgelse.
% p. 104: »the percipient event« is set ROMAN inside the guillemets --- the
% English phrase is not italicised here (italics.py sees nothing on this
% page; confirmed at 300 dpi). "Whitehead's Udtryk" (layer "Udtrvk") and
% "henregner alt" (layer "all") are OCR misreadings, corrected.

Til Grund for hine Spørgsmaal ligger egentlig den populære og dogmatiske
Definition af Sandhed som vor Erkendelses Overensstemmelse med
Virkeligheden. Dette statiske Sandhedsbegreb er, som ofte paavist,
selvmodsigende. Virkeligheden staar ikke fuldt færdig, saa vi blot behøver
at gaa hen og tage den i Øjesyn. Den bliver først til for os gennem strengt
aandeligt Arbejde i Iagttagelse og Tænkning. Dette Arbejdes Historie er
Sandhedens Historie. Det tjener til selve Tilværelsens Karakteristik, at
den for et vist Trin af menneskelig Erkendelses Udvikling maatte tage sig
ud paa en vis bestemt Maade. Vi behøver blot f.~Eks. at sætte os ind i alle
Forudsætninger for Middelalderens Forkastelse af Antipoder for at se,
hvilken Sandhed der paa dette Trin maatte staa som nødvendig. Hvad der er
Virkelighed for Mennesker paa et givet Trin, beror paa, hvor langt det
aandelige Arbejde paa dette Trin er skredet frem. Senere kan saa denne
Virkelighed staa som et ufuldkomment Tilløb. Sandhedens Evighed beror paa
Kontinuiteten i det aandelige Arbejde.

Det er ikke en aandsaristokratisk Tankegang, vi herved føres til. Netop
ifølge Relationsbegrebet er der Sandhed i enhver ærlig Stræben efter
Klarhed. Vi skaber ikke selv de Forudsætninger, ud fra hvilke vi tænker, og
vi
% --- p. 105 ---
bliver os dem end ikke fuldt bevidste. Der er her en Analogi mellem det
intellektuelle og det etiske Liv (se ovenfor p.~81 ff.). I den
Nødvendighed, med hvilken de \guillemotleft
ufuldkomne\guillemotright{} Forestillinger optræder, er den evige Sandhed
selv tilstede. Og ofte har det endog været en Betingelse for Fremskridt i
Sandhedserkendelse, at der blev gjort fuldt Alvor af de (som det senere
viste sig) falske Forudsætninger. Som William Hamilton har sagt: en levende
Usandhed er bedre end en død Sandhed.
% p. 105: WILLIAM HAMILTON stands in the body in ordinary roman, NOT small
% caps (checked at 300 dpi) --- the small-caps convention of this edition
% is confined to the footnotes.

Ud over den stadige Stræben, som det dynamiske Sandhedsbegreb henpeger paa,
kommer vi ikke. Vor Tanke kan ikke springe over sin egen Skygge. Vort
sidste Tilhold er selve det arbejdende Tankeliv i dets stadige Opdagen og
Bestemmen af Relationer. En Afslutning er her ikke mulig.

Bacon og Comte hævdede, at Mennesket skal forklares ved Verden, Verden ikke
ved Mennesket. Men her er intet Dilemma. Vi har, som vi har set (ovenfor
p.~60 ff.), aldrig et absolut Subjekt overfor et absolut Objekt, men nok et
objektivt bestemt Subjekt (S\textsubscript{o}) overfor et subjektivt
bestemt Objekt (O\textsubscript{s}), og de kræver begge en indgaaende
Karakteristik. Og naar vi saa sammenligner dem, maa denne Sammenligning
foretages af et Subjekt, saa at Problemet stadig kommer igen. Den, der vil
forklare Mennesket ud fra Verden, maa ikke glemme, at Verden kun kendes
gennem menneskelig Forsken, og at al Undersøgelse af Forholdet mellem
Verden og Mennesket anstilles af en menneskelig Tanke.

c.~Naar vi tror, at Tilværelsen er mere omfattende, end hvad vor Tanke til
nogensomhelst Tid eller i nogensomhelst Form kan omfatte, begrundes det
ved, at vi atter og atter erfarer, hvorledes nye Relationer og dermed nye
Spørgsmaal dukker op og kræver Bestemmelser. Vi har
% p. 105: the subscripted (S_o) and (O_s) are set with \textsubscript, the
% convention already used for the same pair at p. 63. "Vort sidste Tilhold"
% (layer "Vori") and "atter og atter" (layer "alter") verified at 300 dpi.
% --- p. 106 ---
derfor, som ovenfor vist, ingen Ret til at tro, at Tilværelsen er udtømt
ved de to Urfænomener eller Attributer, der i Korthed kan betegnes som Aand
og Materie. Det er det drillende ved Relationsbegrebet, at det stadig viser
os et Hul i vore Systemer. Selv om vi har formet de skønneste Totaliteter i
vor Tanke, dukker det Spørgsmaal op, hvad denne Totalitet staar i Forhold
til.

Undersøgelse af Totalitetsbegrebet i dets Betydning for Verdensanskuelsen
ender, som jeg andetsteds (Totalitet som Kategori, Kap. 6) har søgt at
vise, i to Problemer: om Muligheden af at udsige noget om Tilværelsen som
en Helhed, og om Muligheden af at bestemme det Værdifuldes Stilling
indenfor en saadan Helhed, --- altsaa i et kosmologisk og i et religiøst
Problem. Naar Fordringen til Begrundelse stilles skarpt og bestemt, kan
disse Problemer kun løses ad Analogiens Vej. Herom vidner ogsaa en kritisk
Undersøgelse af kosmologiske og religiøse Forestillinger. (Smign. den
erkendelsesteoretiske Del af min Religionsfilosofi). For Filosofen vil
saadanne Analogier ligge nærmest, som er hentede fra det arbejdende Tanke-
og Villiesliv, hvis Historie for os er Sandhedens og Godhedens Historie.
Det skyldes en mystisk Trang hos Spinoza, naar han projicerede den
Tankesammenhæng, der for ham var Urfænomenet, ud som en hvilende Substans.
Vore sidste Analogier maa støtte sig til vore betydningsfuldeste
Erfaringer, og de har Arbejde, Stræben og Kamp til Indhold, svarende til
den stadige Brydning mellem vore to mest fundamentale Kategorier, Syntese
og Relation. Uden Fremtræden af Modsætninger og Stræben efter at overvinde
dem kan Tilværelsen, filosofisk set, ikke tænkes. Blandt disse Modsætninger
gør Modsætningen mellem det Værdifulde og det Værdihæmmende sig, praktisk
set, stærkest gældende.
% p. 106: Høffding's own titles "Totalitet som Kategori" and
% "Religionsfilosofi" are set ROMAN in the body, not italic (italics.py
% sees nothing on this page; confirmed at 300 dpi), unlike the titles in
% the footnotes. Transcribed as printed, without \textit.
% --- p. 107 ---
Kan der ikke vindes nogen begrundet Forvisning om, at alle Modsætninger og
Hæmninger kan harmoniseres, saa kan dog Analogien med det arbejdende Tanke-
og Villiesliv afgive et poetisk og symbolsk Udtryk for den Forvisning, at
det Arbejde, vi er stedt i, ikke er et isoleret Træk i Virkeligheden;
dertil hænger Relation og Kamp med Relationer altfor nøje sammen med den
Erkendelse af Tilværelsen, vi overhovedet er i Stand til at vinde.

d.~Kun i Billedets hvilende Form kan vi, for en Stund, naa ud over den
stadige Uro, som formedelst Relationsbegrebet er karakteristisk for enhver
vaagen Tanke. Det er derfor ogsaa til Billedernes Verden, at den, der søger
Hvile og Fred, tager sin Tilflugt. En berømt dansk Teolog har udtalt, at
uden Brug af Billeder vilde vi komme ud i det Tomme, idet vi maatte holde
os til nogle fattige Begreber, der intet Lys kunde give. Ja, Filosofen har
kun sine fattige Begreber at holde sig til. Men hvis disse Begreber under
stadig fornyet Eftertanke viser sig at være grundede i Tankelivets og
Erfaringens Natur, saa kan han have Ret til at tro, at Tilværelsens Puls
banker i de bedste og klareste Tanker, Mennesker kan udforme, og maaske
ikke banker svagest der, hvor nye Spørgsmaal stilles i Kraft af
Relationsbegrebet. Og han forbeholder sig sin Dom over de Analogier, der
ligger til Grund for Billeddannelsen i det enkelte Tilfælde. Ligesaa lidt
som nogen Tanke vil noget Billede fuldstændig kunne udtrykke den
uafsluttelige Tilværelse, hvoraf vi er Led. I den Kendsgerning, at enhver
Lignelse halter, melder Relationsbegrebet sig atter og atter.

I Obstfelder's \guillemotleft En Præsts Dagbog\guillemotright{} er paa
dybsindig og stemningsfuld Maade Kampen mellem Tanke og Billede skildret,
saaledes som den ofte føres i et higende Sind. Digteren er klar over ---
som Filosofen ogsaa kan være
% p. 107: »En Præsts Dagbog« is set ROMAN inside the guillemets; the title
% is not italicised (confirmed at 300 dpi). "idet vi maatte holde os"
% (layer "inaatte") corrected. spacing.py's single flag on "komme" is
% justification.
% --- p. 108 ---
det --- at det Højeste bor \guillemotleft længere inde i os, end vi selv
kan trænge\guillemotright. Men han er ikke tilfredsstillet ved denne
Overbevisning. Han vil have det frem for sig, \guillemotleft komme udenfor
det, der er inde i os\guillemotright, kunne sige du til det, have et
Billede af det. Hvis det var lykkedes ham at naa til dette Maal, vilde han
have gjort den Erfaring, at intet Billede slaar til, men at der gør sig en
stadig Trang til nye Billeder gældende, ligesom Forskeren bestandig
opdager nye Relationer, hvorved hans hidtil udformede Begreber sprænges.

Poesiens Verden staar aaben for alle, ogsaa for Filosofen. Men Filosofen
kan ikke anerkende noget en Gang for alle dannet Billede. Og han vil være
klar over, at naar han er gaaet over i Poesiens Verden, tier Tankerne,
ligesom for den anakreontiske Digter \guillemotleft Bekymringerne
tier\guillemotright, naar han griber sit Bæger. Der kræves stedse nyt
Arbejde, baade med Billeder og med Tanker, for at Livet kan komme til sin
Ret. Maaske er, hvad en ægyptisk Vismand kaldte \guillemotleft et higende
Sind, hvis Ord alle er skjulte\guillemotright, den inderligste og højeste
Form for Aandsliv.
% p. 108: the body of the treatise ENDS here, with "og højeste Form for
% Aandsliv." about two-thirds down the page (verified on the image). There
% is no colophon or end-rule; printed p. 109 carries the Indhold, and PDF
% p. 112 is a blank verso.

% --- p. 109 ---
% p. 109: the printed INDHOLD, transcribed as a page like any other rather
% than generated by \tableofcontents, so that the edition preserves what
% the book itself prints. Layout of the printed page: the heading INDHOLD
% in letterspaced capitals, a short centred rule, the column head "Side" at
% the right, twenty-four entries with dot leaders and page numbers, then a
% second short centred rule and the two dating lines.
%
% LETTERSPACING: the five NUMBERED entries (1.--5.) and the word INDHOLD
% are set spærret; the lettered (a.--f.) and Greek (α--δ) sub-entries are
% NOT --- they are ordinary roman. (The header note at the top of this file
% says the Indhold is letterspaced "throughout"; that is too strong, and
% this is the correcting observation.) Per the edition's convention the
% five spærret entries are set \emph{} here; the letterspacing itself is
% not imported into the text. The whole list is set \small so that the
% longer entries hold one line, as they do in the printing.
%
% COLLATION against the heads the printed pages actually carry
% (every divergence recorded, neither reading normalised):
%   * §3b's α (p. 33), β (p. 40), γ (p. 44) and §3c's α (p. 45), β (p. 55)
%     carry NO title on their own pages --- the printed pages give the bare
%     Greek mark. The titles below exist only here in the Indhold. Already
%     logged at each site.
%   * §3d δ "Den nye Relativitetsteori", p. 72: the Indhold prints a GREEK
%     DELTA, confirming the printed head at p. 72 and settling the doubt
%     recorded in this file's header. The ABBYY layer's "b." for this entry
%     is an OCR misreading; at 300 dpi the mark is unambiguously δ.
%   * §4d, p. 90: the Indhold reads "Etik og de individuelle Relationer"
%     --- "Etik", not "Etiken" --- where §4a--§4c all have "Etiken". The
%     printed head at p. 90 must be checked against this; both readings to
%     be kept if they differ.
%   * §5's four lettered paragraphs a. (p. 99), b. (p. 103), c. (p. 105),
%     d. (p. 107) are NOT listed in the Indhold at all. The Indhold gives
%     §5 a single line.
%   * The Indhold's entries carry no terminal full stop; the printed heads
%     all do. Every other entry --- wording, lettering and page number ---
%     agrees with the printed head.
\clearpage
{\centering\large INDHOLD\par}
% "INDHOLD" is printed in letterspaced capitals; the letterspacing is not
% reproduced.
\vspace{0.6ex}
{\centering\rule{2.4cm}{0.4pt}\par}
\vspace{1.2ex}
\begingroup
\parindent=0pt
\setlength{\parskip}{0pt}
\small
\mbox{}\hfill Side\par
\hangindent=1.6em \emph{1.~Relationsbegrebets Historie}\dotfill 3\par
\hangindent=1.6em \emph{2.~Relationsbegrebet i Psykologien}\dotfill 16\par
\hangindent=1.6em \emph{3.~Relationsbegrebet i Erkendelsesteorien}\dotfill 26\par
\hangindent=3.2em \hspace*{1.6em}a.~Relationsbegrebet og de andre fundamentale Kategorier\dotfill 26\par
\hangindent=3.2em \hspace*{1.6em}b.~Relationsbegrebet og de formale Kategorier\dotfill 33\par
\hangindent=4.8em \hspace*{3.2em}α.~Rene Kvalitetsrækker\dotfill 33\par
\hangindent=4.8em \hspace*{3.2em}β.~Identitet og Relation\dotfill 40\par
\hangindent=4.8em \hspace*{3.2em}γ.~Negationens og Rationalitetens Relationer\dotfill 44\par
\hangindent=3.2em \hspace*{1.6em}c.~Relationsbegrebet og de reale Kategorier\dotfill 45\par
\hangindent=4.8em \hspace*{3.2em}α.~Kausalitet, Realitet og Relation\dotfill 45\par
\hangindent=4.8em \hspace*{3.2em}β.~Totalitetens og Udviklingens Relationer\dotfill 55\par
\hangindent=3.2em \hspace*{1.6em}d.~Relationen Subjekt---Objekt\dotfill 58\par
\hangindent=4.8em \hspace*{3.2em}α.~Indledning\dotfill 58\par
\hangindent=4.8em \hspace*{3.2em}β.~Omkring Kopernikanismen\dotfill 62\par
\hangindent=4.8em \hspace*{3.2em}γ.~Filosofiske Synspunkter\dotfill 67\par
\hangindent=4.8em \hspace*{3.2em}δ.~Den nye Relativitetsteori\dotfill 72\par
\hangindent=1.6em \emph{4.~Relationsbegrebet i Etiken}\dotfill 77\par
\hangindent=3.2em \hspace*{1.6em}a.~Historieforskningen og Relationsbegrebet\dotfill 77\par
\hangindent=3.2em \hspace*{1.6em}b.~Etiken og de psykologiske Relationer\dotfill 81\par
\hangindent=3.2em \hspace*{1.6em}c.~Etiken og de historiske Relationer\dotfill 87\par
\hangindent=3.2em \hspace*{1.6em}d.~Etik og de individuelle Relationer\dotfill 90\par
\hangindent=3.2em \hspace*{1.6em}e.~Kant's Etik og Relationsbegrebet\dotfill 92\par
\hangindent=3.2em \hspace*{1.6em}f.~Religionsfilosofien og Relationsbegrebet\dotfill 94\par
\hangindent=1.6em \emph{5.~Relationsbegrebet i Kosmologien}\dotfill 99\par
\endgroup
\vspace{5ex}
{\centering\rule{2.4cm}{0.4pt}\par}
\vspace{2ex}
{\centering\small Forelagt paa Mødet den 21. Oktober 1921.\par
Færdig fra Trykkeriet den 14. Novbr. 1921.\par}
% p. 109: "Relationen Subjekt---Objekt" is set here with an UNSPACED em
% dash (the ABBYY boxes leave 9.6 pt between "Subjekt" and "Objekt", just
% the width of the dash), which is how the printed head at p. 58 sets it
% too.

\end{document}
